%%%%%%%%%%%%%%%%%%%%%%%%%%%%%%%%%%%%%%%%%%%%%%%%%%%%%%%%%%%%%%%%%%%%%%%%%%%%%%%%
%                     Bachelor thesis template 
%						 Science Faculty 
%								at 
%			National Autonomous University of Mexico (UNAM)
%%%%%%%%%%%%%%%%%%%%%%%%%%%%%%%%%%%%%%%%%%%%%%%%%%%%%%%%%%%%%%%%%%%%%%%%%%%%%%%%
% based on Harish Bhanderi's PhD/MPhil template, then Uni Cambridge
% http://www-h.eng.cam.ac.uk/help/tpl/textprocessing/ThesisStyle/
% corrected and extended in 2007 by Jakob Suckale, then MPI-iCBG PhD programme
% and made available through OpenWetWare.org - the free biology wiki
%
%                     Under GNU License v3
%
% Adapted for the Engineering School at UNAM by Jesús Velázquez y Marco Ruiz
% Then, adapter fot he Science Faculty at UNAM by Jonathan Urrutia
%
%
%
% All used packages are found in  
%
%					./Latex/Classes/PhDthesisPSnPDF.cls
%
% within this las documments there are some lines to be UnCommented for a printed or a digital version of the output file:
%
%			line 184
%			lines 231-261
%
%	Since this is a template for a thesis made at UNAM (Mexico) the titles may be in spaninish but within PhDthesisPSnPDF.cls it can be change
%

\documentclass[11pt]{Latex/Classes/PhDthesisPSnPDF}

\usepackage{blindtext}         %  For dummy text.
% The \blindtext or \Blindtext commands throughout this template generate dummy text to fill the template out.

\input{Latex/Comands}           % Special commands written by the author
\usepackage{csquotes}
\usepackage{setspace}
\usepackage{ulem}

\usepackage{graphicx} 
\graphicspath{{../Figuras/},{./}} % Directory in which figures are stoorange

%%-------------------------------------------------------------------------------
%                                  Information of the Student                                
%-------------------------------------------------------------------------------
\title{Cálculo numérico de propiedades ópticas de eritrocitos sanos
	y enfermos}
\author{Dana Larissa Luna González}       
\degree{FÍSICA BIOMÉDICA}% Physics
\director{M. en C. Jonathan A. Urrutia Anguiano}% Thesis director
\degreedate{2026}% Year of the dissertation
\lugar{Cd. Universitaria, Cd. de México}% Place of the dissertation
\portadatrue

% ----------------------------- Datos del jurado
\student{Luna\\ González\\ Dana Larissa\\ 776 101 4262\\ Universidad Nacional Autónoma de México\\ Facultad de Ciencias\\ Física Biomédica\\ 421122680}
\secretario{M. en C. \\ Jonathan A. \\ Urrutia \\ Anguiano}
\presidente{Dr \\ President \\ Last name \\ Last name}
\vocal{Dr \\ Vocal \\ Last name \\ Last name}
\supuno{Dr \\ substitute 1 \\ Last name \\ Last name}
\supdos{Dr \\ Substitute 2 \\ Last name \\ Last name}
\pags{pages}

\keywords{tesis,autor,tutor,etc}            % For metadata 
\subject{tema_1,tema_2}                     % Sebjucts for metadata


%-------------------------------------------------------------------------------
%                                   COVER                                   
%-------------------------------------------------------------------------------
\begin{document}
%
\maketitle

%-------------------------------------------------------------------------------
%                                   FRONT MATTER                                
%-------------------------------------------------------------------------------
\frontmatter

% !TeX root = ../tesis.tex

\begin{acknowledgements}
\addcontentsline{toc}{chapter}{\protect\numberline{}Agradecimientos}

Quiero expresar mi más sincero agradecimiento a quienes me brindaron su apoyo y me acompañaron durante mis estudios de licenciatura en la Universidad Nacional Autónoma de México (UNAM). Esta tesis representa tanto el resultado de mi esfuerzo como de las enseñanzas, los consejos y el tiempo que distintas personas compartieron conmigo a lo largo de mi formación académica.

En primer lugar, agradezco a mi asesor, el M. en C. Jonathan A. Urrutia Angiano, por su paciencia, su disposición para orientarme y sus consejos durante todo este proceso. Gracias por dedicar tiempo a resolver mis dudas, por acompañarme en las distintas etapas de este trabajo y por ayudarme a enfrentar los retos que surgieron durante su desarrollo. Su apoyo fue fundamental para avanzar en esta investigación y fortalecer mi formación.

Asimismo, agradezco al Dr. Alejandro Reyes Coronado por su guía y orientación. Sus observaciones y consejos fueron valiosos para el desarrollo de esta tesis y contribuyeron a enriquecer mi perspectiva sobre el trabajo de investigación. Agradezco el tiempo y la atención que me brindó, así como la oportunidad de ser parte de su grupo de investigación.

Agradezco al proyecto DGAPA-UNAM PAPIIT IN101825 por el apoyo económico que me otorgó durante un año para desarrollar mi tesis de licenciatura. Este respaldo me permitió dedicar tiempo y esfuerzo a la investigación y fue una ayuda importante para llevar a cabo este trabajo.

De igual manera, agradezco a quienes integran el grupo de Nanoplasmónica por su apoyo, sus consejos y su disposición para escucharme una y otra vez. Gracias por compartir sus conocimientos, por abrir un espacio para conversar sobre mis dudas e ideas y por acompañarme en los momentos de dificultad. El intercambio de ideas y el aprendizaje compartido hicieron de este proceso una experiencia más enriquecedora, tanto en lo académico como en lo personal.

Finalmente, agradezco a la Universidad Nacional Autónoma de México por la formación recibida y por las oportunidades de aprendizaje que me ofreció durante la licenciatura. Me llevo de esta etapa conocimientos, experiencias y enseñanzas que seguirán acompañándome en mi desarrollo académico y profesional.


\end{acknowledgements}
       %
% --- JAUA ---
% Inventate unos agradecimientos chafones y genéricos al grupo y pregúntale a taata  cómo procedemos con PAPIIT porque técniamente tuviste agradecimentos aunque no se presentre el trabajo coo suyo.
% Taambién quitamos la declaración que no usamos en ciencias y el quote luego te sacas uno chido :)
%   
% \include{0-2-Declaration/1-Declaration}
% \include{0-3-Quote/1-Quote}                  
% !TeX root = ../tesis.tex

% Thesis Abstract -----------------------------------------------------

%\begin{abstractslong}    %uncommenting this line, gives a different abstract heading
\begin{abstracts}        %this creates the heading for the abstract page
\addcontentsline{toc}{chapter}{\protect\numberline{}Resumen}

Los eritrocitos presentan propiedades ópticas que dependen de su geometría, tamaño, composición y orientación, por lo que alteraciones morfológicas y variaciones en la concentración de hemoglobina pueden modificar la manera en que absorben y esparcen la luz. Esta sensibilidad ha motivado el uso de técnicas ópticas, como la citometría de flujo, para la caracterización y diferenciación celular. Sin embargo, gran parte de los estudios numéricos del esparcimiento de luz por eritrocitos se han concentrado en longitudes de onda cercanas a 632.8 nm, fuera de las principales bandas de absorción de la hemoglobina. En este trabajo de tesis de licenciatura se estudian numéricamente las propiedades ópticas de eritrocitos sanos y de dos alteraciones morfológicas representativas, el esferocito y el microcito, en el intervalo espectral de 400 a 600 nm. Las distintas morfologías se modelaron mediante óvalos de Cassini, mientras que sus propiedades ópticas se incorporaron a través de funciones dieléctricas dependientes de la longitud de onda y de la concentración de hemoglobina, obtenidas a partir de datos experimentales mediante las relaciones de Kramers-Kronig sustractivas. Las simulaciones se realizaron mediante el método \textit{Finite Integration Technique} en CST Studio Suite. Los parámetros óptimos para realizar las simulaciones se establecieron mediante un análisis paramétrico y comparación con la solución analítica de Mie para una esfera, obteniéndose errores relativos inferiores al 5\%, y posteriormente se contrastó el patrón angular de un eritrocito con resultados numéricos reportados en la literatura. Los resultados muestran que la respuesta óptica depende del efecto combinado de la geometría y de las propiedades dieléctricas. La banda de Soret, alrededor de 431 nm, presenta las mayores diferencias en la eficiencia de extinción entre las morfologías consideradas. Aunque el esparcimiento se concentra predominantemente en la dirección frontal, las mayores diferencias relativas aparecen en regiones laterales y posteriores. Finalmente, el análisis mediante ventanas angulares permitió identificar una región alrededor de $\theta\approx130^\circ$ en la que el contraste entre eritrocitos sanos y patológicos aumenta de manera consistente para las tres longitudes de onda analizadas. Estos resultados sugieren condiciones espectrales y angulares potencialmente útiles para futuros esquemas de diferenciación celular mediante citometría de flujo.



\end{abstracts}
%\end{abstractlongs}


% ----------------------------------------------------------------------       

%-------------------------------------------------------------------------------
%                                INDICES                                    |
%-------------------------------------------------------------------------------
%
\setcounter{secnumdepth}{3} % organisational level that receives a numbers
\setcounter{tocdepth}{3}    % print table of contents for level 3

\tableofcontents            % Print main index

%: ----------------------- list of figures/tables ------------------------
%\listoffigures              % Genera el ínidce de figuras, comentar línea si no se usa
%\listoftables               % Genera índice de tablas, comentar línea si no se usa


%-------------------------------------------------------------------------------
%                                MAIN MATTER                                   %-------------------------------------------------------------------------------
% the main text starts here with the introduction, 1st chapter,...
\mainmatter

\def\baselinestretch{1}                   % Line spacing

% !TeX root = ../tesis.tex

\chapter*{Introducción}
\addcontentsline{toc}{chapter}{\protect\numberline{}Introducción}	  		% Comment if you don't want the introduction to appear on the table of content. It will not have a number
\label{chapter:intro}

% this file is called up by thesis.tex
% content in this file will be fed into the main document

%
% --- JAUA ---
% Luz-materia puede confundir con propiedaes electrónicas de materiales, o de excitacines específicas (como plasmones o multipolos). Mejor habla del pronblema de espacimento por partículas.
% También omitamo biofotónica porque nunca usaste ese térmio en la tsis. Enuncia mejor diversas aplicaciones como imagenologhía, snesado, etc.
%

La interacción de la luz con la materia depende de propiedades como el tamaño, la geometría, la composición, la estructura interna y la orientación del sistema considerado \cite{bectondickinsonIntroductionFlowCytometry2000}. En este contexto, la biofotónica estudia el uso de la luz para analizar sistemas biológicos y desarrollar métodos de diagnóstico, monitoreo y tratamiento \cite{baldiniShiningLightFuture2025}. En el área de diagnóstico, la biofotónica aprovecha el hecho de que las alteraciones fisiológicas o patológicas en sistemas de células o en células individuales pueden modificar la manera en que interactúan con la luz, proporcionando información potencialmente útil para su caracterización \cite{baldiniShiningLightFuture2025}. %
% --- JAUA ---
% Corta acá el párrafo y comienza con el siguiente. Rescribe acorde
%

Un ejemplo de ello es el estudio del esparcimiento de luz en la sangre, cuya respuesta óptica puede mostrar cambios en las propiedades morfológicas y composicionales de sus componentes celulares asociados a distintas enfermedades \cite{ergulComputationalStudyScattering2010}.

Entre los principales componentes celulares de la sangre se encuentran los eritrocitos, cuya función fisiológica fundamental es el transporte de oxígeno \cite{palTextbookMedicalPhysiology2021}. En el intervalo %
% --- JAUA ---
% intervalo suena muy IA pocha. En el rango/espectro, o alfo .
% visibke, los eritrocitos, constituidos principalmente por hemoglobina, dominan la.....
%
 visible, estas células dominan la respuesta electromagnética de la sangre y están constituidas principalmente por hemoglobina \cite{bosschaartLiteratureReviewNovel2014}. En condiciones fisiológicas %
 % --- JAUA ---
 % Considerando el caso de una persona sana, los eritrocios...
 %
  presentan una geometría discoidal bicóncava; sin embargo,%
  % --- JAUA ---
  % innecesario el punto y coma, no??
  % cambia tamño por volumen
  distintas alteraciones hematológicas pueden modificar tanto su morfología como su composición \cite{palTextbookMedicalPhysiology2021}. Estas modificaciones pueden manifestarse, entre otros parámetros, como variaciones en la forma, el tamaño y la concentración de hemoglobina corpuscular media (CHCM) \cite{palTextbookMedicalPhysiology2021}.

%
% --- JAUA ---
% Mehir di que estas alteracion corresponden a respuestas espectrales y patrones de interferencia distintas, que podrían emplearse para caracterisar las enfermendas En particular, en el espectro visible, la hemoglobina cuenta con una estructa de bandas a longitudes de ondas caracterísicas que son relevantes para su fincón disiológica....
%
  El interés en estudiar estas modificaciones radica en que la respuesta electromagnética de los eritrocitos depende del efecto conjunto de su geometría y de su composición. 

En el intervalo visible, la hemoglobina presenta una estructura espectral característica, dominada por la banda de Soret y las bandas $\alpha$ y $\beta$, asociadas a transiciones electrónicas de la molécula \cite{dayerBandAssignmentHemoglobin2010}. Estas bandas producen variaciones en las partes real e imaginaria del índice de refracción y, por tanto, en los procesos de absorción y esparcimiento. Por ello, cambios en la morfología y en la CHCM pueden modificar tanto la respuesta espectral como la distribución angular de la luz esparcida por los eritrocitos \cite{bectondickinsonIntroductionFlowCytometry2000}. La sensibilidad de la respuesta óptica a estas propiedades ha motivado el desarrollo de técnicas ópticas para la identificación y diferenciación celular. Entre ellas, la citometría de flujo, que permite analizar grandes poblaciones de células a partir de la luz esparcida \cite{bectondickinsonIntroductionFlowCytometry2000}. Sin embargo, las señales registradas por un detector de citometría de flujo corresponden a la medición del campo esparcido dentro de una región angular determinada por la geometría del sistema de medición \cite{mckinnonFlowCytometryOverview2018}. En consecuencia, el conocimiento de la distribución angular completa del esparcimiento puede aportar información adicional para identificar regiones de observación particularmente sensibles a cambios en la morfología y en la composición de los eritrocitos.

La descripción teórica de esta respuesta constituye un problema electromagnético complejo %
% --- JAUA ---
% Es una ecuación diferencial que no oscila tonto en sus condicones de feontera, por qué sería complejo??  Hay que tener cuodado con los teŕminos que usamos,. 
%
 debido tanto a la geometría bicóncava de los eritrocitos como a que sus dimensiones son mayores que la longitud de onda de la radiación incidente. Por esta razón, el estudio del esparcimiento de luz por eritrocitos se ha abordado mediante distintos métodos numéricos. Lu et al. \cite{luSimulationsLightScattering2005}, por ejemplo, emplearon el método de \textit{Finite Difference Time Domain} (FDTD); Karlsson et al. \cite{karlssonNumericalSimulationsLight2005} emplearon el método de Discrete Dipole
Approximation (DDA) y Tsinopoulos et al. \cite{tsinopoulosLightScatteringAggregated2002} desarrollaron un enfoque basado en el Boundary Element Method (BEM). No obstante, los trabajos mencionados anteriormente se concentran principalmente en longitudes de onda cercanas a 632.8 nm, fuera de las principales bandas de absorción de la hemoglobina. Esto limita el estudio de la dependencia espectral de la respuesta óptica en regiones donde las propiedades ópticas de los eritrocitos presentan variaciones más pronunciadas. Por lo tanto, resulta de interés extender el análisis a un intervalo espectral que incluya la banda de Soret y las bandas $\beta$ y $\alpha$, considerando además las modificaciones morfológicas y composicionales asociadas a eritrocitos patológicos.

%
% --- JAUA ---
% Estos  cuatro parrafos estan bien escritos pero muy repetitvos.  Recuerda que acá no hay un resumen estricto de tu trabajjo entonees:
% Unidica los siguientes tres párrafos en uno solo: reduce el segundo párrafo porque no necesitamos información tan detallada. Y el pultimo párrafo reescielo omitiendo Capítulo 1, Capítulo 2.... Mejor pon sus títulos e hipervínculos a cada uno de ellos.

%
% --- JAUA ---
% De nuevo: todo esto lo tenemos que trabajar a profundidad de nuevo: te lo pongo muchas veces para que no creas que hemos acabado :)
%

A partir de esto, el objetivo de esta tesis es estudiar  cómo las modificaciones en la morfología y en la concentración de hemoglobina alteran las propiedades ópticas de los eritrocitos y %
% --- JAUA ---
% no pongas esto porque en principio y a lo hiciste -> si es posible 
identificar condiciones espectrales y angulares que aumenten el contraste entre células sanas y patológicas. Para ello, se estudian tres morfologías representativas: un eritrocito sano, un esferocito y un microcito. El eritrocito sano se modela mediante óvalos de Cassini, que permiten reproducir su forma bicóncava con un número reducido de parámetros geométricos; el esferocito se representa mediante una geometría esférica, mientras que el microcito conserva la forma bicóncava pero presenta dimensiones reducidas. A cada morfología se le asigna además una CHCM representativa, con el fin de incorporar las variaciones en su composición asociadas a estas alteraciones.

Para describir las propiedades ópticas de los eritrocitos se emplean datos experimentales del índice de refracción de la hemoglobina reportados en la literatura para distintas concentraciones. Debido a que estos datos se encuentran restringidos a ventanas espectrales finitas, se expanden las funciones dieléctricas al rango de 400 a 600 nm mediante un ajuste de osciladores de Lorentz para la parte imaginaria y la reconstrucción de la parte real mediante las relaciones de Kramers-Kronig sustractivas. Posteriormente, se estudia la convergencia numérica del método empleando como referencia la solución analítica de Mie para una esfera, para luego contrastar el modelo de eritrocito con resultados numéricos reportados en la literatura. Una vez establecidas las condiciones de simulación, se calculan las eficiencias y secciones transversales de extinción, absorción y esparcimiento, así como la sección transversal diferencial de esparcimiento y el factor de anisotropía en el intervalo espectral comprendido entre 400 y 600 nm. Finalmente, se analizan las bandas de Soret, $\beta$ y $\alpha$ y se estudia la distribución angular de la luz esparcida con el propósito de localizar intervalos de detección en los que aumente el contraste entre las distintas morfologías.

De esta manera, el objetivo general de esta tesis es caracterizar numéricamente las propiedades ópticas de eritrocitos sanos y con alteraciones morfológicas y determinar regiones espectrales y angulares potencialmente útiles para su discriminación. Como objetivos específicos se plantea construir modelos geométricos y ópticos representativos de las distintas células, verificar la convergencia y consistencia del método numérico, estudiar la dependencia espectral de las propiedades ópticas, comparar sus patrones angulares de esparcimiento y localizar ventanas de detección en las que las diferencias entre las morfologías sean más pronunciadas. 

Esta tesis está dividida en cuatro capítulos:  en el Capítulo 1 se presenta el formalismo electromagnético empleado para describir el esparcimiento de luz por partículas, incluyendo las ecuaciones de Maxwell, las secciones transversales y las relaciones de Kramers-Kronig y Kramers-Kronig sustractivas. En el Capítulo 2 se estudian la morfología, composición y principales alteraciones de los eritrocitos, así como la citometría de flujo y la construcción de las funciones dieléctricas empleadas en las simulaciones. En el Capítulo 3 se introduce el método de integración finita y se establece la configuración numérica mediante un análisis de convergencia. En el Capítulo 4 se presentan y discuten los resultados correspondientes al eritrocito sano, el esferocito y el microcito, incluyendo su respuesta espectral, sus patrones angulares de esparcimiento y la identificación de regiones de mayor contraste. Finalmente, se presentan las conclusiones generales y las principales líneas de trabajo a futuro.            
% !TeX root = ../tesis.tex

\chapter{Esparcimiento de luz por partículas}
\label{chapter:theory}

\vspace*{7em}
%
% --- JAUA ---
% Ya que tenemos una idea del texto final, vsmos a cambiar iun poco las estrucutras del primer párrafo de cada sección. La idea es irlos ligando entre sí. La propuesta es; Primera oración enunciando el tema general; ssegunda: para qué es util con nuestro propópsito de la tesis. A partir de la tercera oración ya comienzas con los contenidos de este capítlo. Hay que revisar esto en todo inicio de sección/subsección, etc.
%


En este capítulo se estudia la interacción electromagnética entre la luz y la materia mediante una formulación basada en la función dieléctrica dependiente de la frecuencia de la luz considerada. En la primera parte se presentan las ecuaciones de Maxwell y las relaciones constitutivas para sistemas lineales, homogéneos e isótropos y se analiza la solución de ondas planas, la cual se emplea para describir el flujo de energía transportado por una onda electromagnética mediante el vector de Poynting y su promedio temporal. En la segunda parte se derivan las secciones transversales de absorción, esparcimiento y extinción, que permiten cuantificar la interacción entre una onda incidente y una partícula \cite{bohrenAbsorptionScatteringLight2008}. Adicionalmente, en la Sección \ref{section:rKK} se estudian las relaciones de Kramers–Kronig aplicadas a la función dieléctrica. Estas relaciones permiten vincular las partes real e imaginaria de cualquier función óptica lineal compleja, condicionada a pruebas de causalidad y analiticidad en el contexto de funciones complejas \cite{lucariniKramersKronigRelationsOptical2005}. Finalmente, se discuten las limitaciones de aplicar las relaciones de Kramers-Kronig a datos experimentales con un rango finito de frecuencias y se presenta la derivación de las relaciones de Kramers–Kronig sustractivas. 


\section{Formalismo}
\label{section:basics}
Las ecuaciones de Maxwell, junto con la fuerza de Lorentz\footnote{La fuerza de Lorentz es la fuerza que experimenta una partícula de carga \( q \) que se desplaza con velocidad \( \vb{v} \) en presencia de un campo eléctrico \( \vb{E} \) y un campo magnético \( \vb{B} \). En unidades del SI está dada por $\vb{F}=q(\vb{E}+\vb{v}\times\vb{B})$ \cite{zangwillModernElectrodynamics2013}. }, describen a la electrodinámica clásica, que se centra en el origen y comportamiento de los campos electromagnéticos \cite{zangwillModernElectrodynamics2013}. En unidades del Sistema Internacional (SI), las ecuaciones de Maxwell se expresan en forma diferencial como
\cite{griffithsIntroductionElectrodynamics2023}
%
	\begin{subequations} \label{eqs:Maxwell}
	\begin{tcolorbox}[
	ams align, breakable]
	\nabla \cdot\vb{E} &= \frac{\rho_{\text{tot}}}{\varepsilon_0}, &\mbox{(Ley de Gauss eléctrica)}  
	\label{seq:GE} \\
	\nabla \cdot\vb{B} &= 0,						&\mbox{(Ley de Gauss magnética)}   
	\label{seq:GM} \\
	\nabla \times\vb{E} &= -\pdv{\vb{B}}{t}, 	&\mbox{(Ley de Faraday-Lenz)}		
	\label{seq:FL}\\
	\nabla \times\vb{B} &= \mu_0 \vb{J}_{\text{tot}} +\varepsilon_0\mu_0 \pdv{\vb{E}}{t}, &
	\mbox{(Ley de Ampère-Maxwell)} \label{seq:AM}
	\end{tcolorbox}\end{subequations}\noindent
%
donde se obvian las dependencias espaciales y temporales y $\vb{E}(\vb{r},t)$ es el campo eléctrico,  $\vb{B}(\vb{r},t)$ es el campo magnético; $\rho_{\text{tot}}$  es la densidad volumétrica de carga total, $\vb{J}_{\text{tot}}$ la densidad volumétrica de corriente  total; $\epsilon_0$ es la permitividad eléctrica en el vacío y $\mu_0$ la permeabilidad magnética en el vacío. 

\noindent Las ecuaciones de Maxwell se pueden reescribir en términos del vector de desplazamiento eléctrico $\vb{D}$ y el campo $\vb{H}$, dados por \cite{griffithsIntroductionElectrodynamics2023}
%
\begin{equation}
	\vb{D}(\vb{r},t)=\epsilon_0\vb{E}(\vb{r},t)+\vb{P}(\vb{r},t) \hspace{1cm}\text{y}\hspace{1cm}
	\vb{H}(\vb{r},t)=\frac{\vb{B}(\vb{r},t)}{\mu_0}-\vb{M}(\vb{r},t), \label{eq:D_H}
\end{equation}
%
donde $\vb{P}$ es la polarización y $\vb{M}$ es la magnetización de un material determinado, que corresponden al momento dipolar eléctrico y magnético por unidad de volumen, respectivamente, inducidos en el material. La relación lineal más general entre $\vb{E}$ y $\vb{D}$ está dada por \cite{jacksonClassicalElectrodynamics2021}
%
\begin{tcolorbox}[ams align]
	\vb{D}(\vb{r},t)=\int_{\infty}^{\infty}\int_{V'}\stackrel{\leftrightarrow}{\varepsilon}(\vb{r}-\vb{r}',t-t')\vdot\vb{E}(\vb{r}',t')\dd{V}'\dd{t'}\
	\label{eq:d} 
\end{tcolorbox}
%	
\noindent donde $\stackrel{\leftrightarrow}{\varepsilon}$ es un tensor de rango dos y representa una función de respuesta\footnote{Las funciones de respuesta proporcionan información sobre el efecto de las interacciones electromagnéticas con el medio y sobre la estructura y propiedades del propio medio \cite{kirzhnitsGeneralPropertiesElectromagnetic1987}.} del material en el espacio y el tiempo \cite{jacksonClassicalElectrodynamics2021} y $V'$ es el volumen donde se encuentra el material. La Ec. \eqref{eq:d} muestra que el campo de desplazamiento eléctrico al tiempo $t$ depende del campo eléctrico en todos los demás tiempos $t'$. Adicionalmente, el campo de desplazamiento eléctrico en un punto $\vb{r}$ depende de los valores del campo eléctrico en puntos vecinos $\vb{r}'$ \cite{jacksonClassicalElectrodynamics2021}. A estas características, se les conoce como la no localidad temporal y espacial, respectivamente, entre $\vb{D}$ y $\vb{E}$ \cite{jacksonClassicalElectrodynamics2021}.

En un medio isótropo y homogéneo, se cumple que $\stackrel{\leftrightarrow}{\varepsilon} = \varepsilon\; \text{diag}(1,1,1)$ donde $\varepsilon$ es una función escalar~\cite{jacksonClassicalElectrodynamics2021} y $\text{diag}(1,1,1)$ es la matriz identidad; en este caso, tanto $\varepsilon$ como $\vb{E}$ son independientes de $\vb{r}'$ y dependen únicamente de la diferencia temporal, es decir \cite{jacksonClassicalElectrodynamics2021}
%
\begin{tcolorbox}[ams align]
	\vb{D}(\vb{r},t)=\int_{-\infty}^{\infty}\varepsilon(t-t')\vb{E}(\vb{r},t')\dd{t'},
	\label{eq:d1} 
\end{tcolorbox}
%	
%
% --- JAUA ---
% Tu pie de página queda cortado. Habría que definir primero la tranformada de fuourier como simnolo, y luego el teorema de la convolución. Te recomiendo poner un $\omega$ luego de escribir frecuencias en el texrto principal.
%
\noindent de modo que la Ec. \eqref{eq:d1} es una convolución\footnote{La convolución de dos funciones $f(t), g(t)$ está dada por $f\ast g=\int_{-\infty}^{\infty}f(t-t')g(t')\dd{t'}$ \cite{arfkenMathematicalMethodsPhysicists2011}.} en el tiempo. Al aplicar la transformada de Fourier y el teorema de la convolución\footnote{El teorema de la convolución es $\mathscr{F}[f\ast g]=\mathscr{F}[f]\mathscr{F}[g]$, donde $\mathscr{F}$ es la transformada de Fourier de una función. En este trabajo se considera la convención de la transformada de Fourier como \cite{arfkenMathematicalMethodsPhysicists2011}
	\begin{equation*}
		 f(\omega)=\mathscr{F}[f(t)]=\int_{-\infty}^{\infty}f(t)e^{i\omega t}\dd{t},\;\;\; \text{y}\;\;\; f(t)=\mathscr{F}^{-1}[f(\omega)]=\frac{1}{2\pi}\int_{-\infty}^{\infty}f(\omega)e^{i\omega t}\dd{\omega},
	\end{equation*} 
donde la dependencia indica si se trata de la transformada de Fourier directa (dependencia en la frecuencia) o inversa (dependencia en el tiempo).}, la Ec. \eqref{eq:d1} se reescribe como un producto en el espacio de frecuencias. Realizando el mismo procedimiento con las relaciones constitutivas que involucran a los campos magnéticos $\mathbf{B}$ y $\mathbf{H}$, se llega a un resultado análogo. Esto es \cite{bohrenAbsorptionScatteringLight2008}
%
\begin{tcolorbox}[ams align]
		\vb{D}(\omega)=\varepsilon(\omega)\vb{E}(\omega)\hspace{2cm}\text{y}\hspace{2cm}\vb{B}(\omega)=\mu(\omega)\vb{H}(\omega).
	\label{eq:d3_b3} 
\end{tcolorbox}
%
\noindent Así, en el espacio de frecuencias, la respuesta del medio queda completamente caracterizada por una función dieléctrica y una permeabilidad magnética dependiente de la frecuencia. Adicionalmente, la relación lineal entre $\vb{D}$ y $\vb{E}$ se debe a que tanto $\vb{P}$  como  $\vb{M}$ tienen a su vez esta relación con los campos $\vb{E}$ y $\vb{H}$, respectivamente. Entonces, en un medio lineal, homogéneo e isótropo, $\vb{P}$ es paralelo a 
 $\vb{E}$ y $\vb{M}$ es paralelo a $\vb{H}$, con coeficientes de proporcionalidad $\varepsilon_0 \chi(\omega)$ y $\chi_m(\omega)$, respectivamente, es decir \cite{griffithsIntroductionElectrodynamics2023}
\begin{align}
	\vb{P}=\varepsilon_0\chi(\omega)\vb{E}\hspace{2cm}\text{y}\hspace{2cm}
	\vb{M}=\chi_m(\omega)\vb{H},
	\label{eq:pol_mag}
\end{align}
%
donde $\chi$ y $\chi_m$ son cantidades adimensionales que corresponden a la susceptibilidad eléctrica y magnética, respectivamente. Al sustituir las Ecs. \eqref{eq:pol_mag} en las Ecs. \eqref{eq:D_H} y al comparar con las Ecs. \eqref{eq:d3_b3} se obtiene que
%
\begin{equation}
	\varepsilon(\omega)=\varepsilon_0[1+\chi(\omega)]\hspace{1cm}\text{y}\hspace{1cm}	\mu(\omega)=\mu_0[1+\chi_m(\omega)].
\end{equation}
%.
%
% --- JAUA ---
% en asencia de fuentes externas e inducidas, no? de lo contrario no son cero las totales. no??
%
En ausencia de fuentes externas ($ \rho_{\text{tot}} = 0, \vb{J}_{\text{tot}} = \vb{0}$) y en un medio lineal, homogéneo e isótropo, las Ecs. \eqref{eqs:Maxwell} pueden desacoplarse y, al aplicar la transformada de Fourier temporal, se obtienen las ecuaciones de Helmholtz para $\vb{E}$ y para $\vb{B}$ \cite{jacksonClassicalElectrodynamics2021}
%
\begin{equation}
	\laplacian{\vb{E}}+k^2\vb{E}=0\hspace{2cm}\text{y}\hspace{2cm} \laplacian{\vb{B}}+k^2\vb{B}=0,
	\label{eqs:helmholtz}
\end{equation}
%
con $k$ una constante denominada número de onda. Una solución a la ecuación de Helmholtz son las ondas planas, en las que los campos están descritos por superficies de fase constante \cite{griffithsIntroductionElectrodynamics2023}
%
\begin{equation}
	\vb{E}(\vb{r},t) =\vb{E}_0 e^{i(\vb{k}\cdot\vb{r} -\omega t)}\hspace{2cm}\text{y}\hspace{2cm} \vb{B}(\vb{r}, t) =\vb{B}_0 e^{i(\vb{k}\cdot\vb{r} -\omega t)},	
	\label{eqs:ondas_planas}
\end{equation}
%
\noindent donde $\vb{E}_0$ y $\vb{B}_0$ corresponden a las amplitudes de los campos, $\omega$ es la frecuencia angular de la onda y $\vb{k}$ el vector de onda. Para que las Ecs. \eqref{eqs:ondas_planas} satisfagan las Ecs. \eqref{eqs:helmholtz}, se impone la relación de dispersión de una onda plana dada por $k=\sqrt{\mu\varepsilon}\;\omega$, donde $\varepsilon$ y $\mu$ corresponden a la permitividad eléctrica y la permeabilidad magnética del medio \cite{jacksonClassicalElectrodynamics2021}. La relación de dispersión se puede reescribir en términos del índice de refracción del material dado por \cite{jacksonClassicalElectrodynamics2021}
\vspace{-0.1cm}
%
\begin{tcolorbox}[ams align]
	n(\omega) = \sqrt{\dfrac{\varepsilon(\omega)\mu(\omega)}{\varepsilon_0\mu_0 }},
	\label{eq:indice} 
\end{tcolorbox}
%	
\noindent con lo que se obtiene
%
\begin{equation}
	k(\omega) =\dfrac{\omega n(\omega)}{c},
	\label{eq:numero_onda} 
\end{equation}

\noindent donde $c=1/\sqrt{\varepsilon_0\mu_0}$ es la velocidad de la luz en el vacío.

El flujo de energía que transporta una onda electromagnética en la dirección de propagación se describe mediante el vector de Poynting  $\vb{S}$ \cite{griffithsIntroductionElectrodynamics2023}
\begin{tcolorbox}[ams align]
	\vb{S}=\frac{1}{\mu_0}(\vb{E}\times\vb{B}).
	\label{eq:vect_Poynting} 
\end{tcolorbox}
\noindent El vector de Poynting tiene unidades de energía por unidad de área y por unidad de tiempo. Para una onda armónica\footnote{Es decir, ondas electromagnéticas $u(\vb{r},t)$ cuya dependencia temporal se describe como $u(\vb{r},t)=~\hat{u}(\vb{r})\,\text{exp}(-i\omega t)$~\cite{zangwillModernElectrodynamics2013}.}, $\vb{E}$ y $\vb{B}$ oscilan en el tiempo, de modo que $\vb{S}$ también varía en el tiempo, por lo que se emplea el vector de Poynting promedio \cite{bohrenAbsorptionScatteringLight2008}
\begin{tcolorbox}[ams align]
	\langle\vb{S}\rangle_t = (1/2) \text{Re}\{\vb{E} \times (\vb{B}^*/\mu_0)\},
	\label{eq:vect_Poynting_prom_arm} 
\end{tcolorbox}
\noindent donde $	\langle\vb{S}\rangle_t$ denota el promedio temporal, $\Re{\cdot}$ la parte real de un número complejo y $*$ la operación complejo conjugado.

%
% --- JAUA ---
% Por qué noescribimos esto en términos de H?? Me parece más directo con la seción anterior y como son medios bien comportados da iugal un poco. 
%













	% !TeX root = ../tesis.tex


\section{Secciones transversales}
\label{section:SC}

%
% --- JAUA ---
% Hayq ue pensar mejor como ligar estas dos secciones. Le falta una oracion que las una ya sea al ginal del anterior o al incio de esta.



%
El problema de interacción de la luz con la materia se puede estudiar al considerar una onda plana [Ec. \eqref{eqs:ondas_planas}] interactuando con una partícula de un material descrito por su índice refracción. La respuesta electromagnética de la partícula puede describirse clásicamente mediante dos procesos fundamentales: la absorción y el esparcimiento, cuyo efecto conjunto se denomina extinción \cite{bohrenAbsorptionScatteringLight2008}. Experimentalmente, estos fenómenos se evidencian al comparar la potencia medida por un detector con y sin la presencia de una partícula iluminada por un campo electromagnético incidente  ($\vb{E}_{\text{inc}}, \vb{H}_{\text{inc}}$) la cual genera un campo electromagnético esparcido ($\vb{E}_{\text{sca}}, \vb{H}_{\text{sca}}$) como se ejemplifica en la Fig. \ref{scattering}, donde se muestra una partícula de geometría arbitraria iluminada por un campo $(\mathbf{E}_{\text{inc}},  \mathbf{H}_{\text{inc}})$ representado por vectores en rojo paralelos. La interacción del campo con la partícula genera un campo electromagnético esparcido $(\mathbf{E}_{\text{sca}}, \mathbf{H}_{\text{sca}})$ representado por los vectores en rojo que surgen de la partícula y se propagan en distintas direcciones. Este campo puede ser medido con un detector~\cite{bohrenAbsorptionScatteringLight2008}.
%
\begin{figure}[t]
	\centering
	\includegraphics[width=8.3cm]{1-Esparcimiento/scattering.pdf}
	\caption{Diagrama del problema de esparcimiento. El campo electromagnético ($\vb{E}_{\text{inc}}, \vb{H}_{\text{inc}}$) incide sobre una partícula arbitraria, produciendo el campo electromagnético esparcido ($\vb{E}_{\text{sca}}, \vb{H}_{\text{sca}}$).}
	\label{scattering}
\end{figure}
%

Para cuantificar la extinción, absorción y el esparcimiento de manera macroscópica, en esta sección se introducen las secciones transversales de extinción, absorción y esparcimiento, al considerar a la partícula embebida en un medio no absorbente e iluminada por una onda plana~\cite{bohrenAbsorptionScatteringLight2008}. Bajo estas condiciones, la energía transportada por los campos electromagnéticos que es absorbida por la partícula $W_{\text{abs}}$, se calcula al integrar el promedio temporal del vector de Poynting  $\langle\vb{S}\rangle_t$ [Ec. \eqref{eq:vect_Poynting_prom_arm}] sobre una superficie cerrada $A$. Para simplificar el análisis, dicha superficie se toma como una esfera de radio $r$ que encierra completamente a la partícula, como se muestra en la Fig. \ref{WA}, es decir, 
\begin{tcolorbox}[ams align]
	W_{\text{abs}}=-\int_A \langle\vb{S}\rangle_t\cdot\vu{e}_r \text{ dA},
	\label{flujopoynting}
\end{tcolorbox}
\noindent donde $\langle\vb{S}\rangle_t$ y $\vu{e}_r$ son paralelos y el signo negativo garantiza que $W_{\text{abs}}$ sea una cantidad positiva definida, asumiendo que únicamente hay procesos de absorción y no de emisión \cite{bohrenAbsorptionScatteringLight2008}. Es posible descomponer el vector de Poynting total en tres contribuciones \cite{bohrenAbsorptionScatteringLight2008}
%		
\begin{subequations}%
	\begin{tcolorbox}[
		ams align, breakable]
\langle\vb{S}_{\text{inc}}\rangle_t&=  \frac{1}{2}\:\mbox{Re}\{\vb{E}_{\text{inc}}\times\vb{H}_\text{inc}^*\},\\%
	\langle\vb{S}_{\text{sca}}\rangle_t&=  \frac{1}{2}\:\mbox{Re}\{\vb{E}_{\text{sca}}\times\vb{H}_{\text{sca}}^*\},\\%
	\langle\vb{S}_{\text{ext}}\rangle_t&= \frac{1}{2}\:\mbox{Re}\{\vb{E}_{\text{inc}}\times\vb{H}_{\text{sca}}^*+\vb{E}_{\text{sca}}\times\vb{H}_\text{inc}^*\},\label{eq:S}%
\end{tcolorbox}
\end{subequations}
%
\noindent donde los subíndices \enquote*{inc}, \enquote*{sca} y \enquote*{ext} se refieren al campo incidente, al campo esparcido y a la interacción entre los dos anteriores, respectivamente. Al integrar estas contribuciones sobre $A$ se obtiene que $W_{\text{abs}}=~W_{\text{inc}}-~W_{\text{sca}}+~W_{\text{ext}}$, donde \cite{bohrenAbsorptionScatteringLight2008}
%
\begin{subequations}\label{eqs:W}
	\begin{align} 
		W_{\text{inc}}  &= -\int_{A}\langle\vb{S}_{\text{inc}}\rangle_t\cdot\vu{e}_r\text{ dA}, \label{seq:Wi}\\
		W_{\text{sca}} &= \int_{A}\langle\vb{S}_{\text{sca}}\rangle_t\cdot\vu{e}_r\text{ dA}, \label{seq:Ws} \\ 
		W_{\text{ext}} &= -\int_{A}\langle\vb{S}_{\text{ext}}\rangle_t\cdot\vu{e}_r\text{ dA}. \label{seq:Wext}
	\end{align}
\end{subequations}
%
%
\begin{figure}[tb]
	\centering
	\includegraphics[width=6.6cm]{1-Esparcimiento/WA.pdf}
	\caption{Esquema de la superficie de integración $A$ como una esfera de radio $r$ centrada en el origen y que encierra a una partícula.}
	\label{WA}
\end{figure}
%
Los signos de las Ecs. \eqref{eqs:W} se eligen de forma que sean cantidades positivas considerando las direcciones de los vectores de Poynting y de $\vu{e}_r$. Como la energía del campo incidente que cruza la esfera es la misma que la que entra, $W_{\text{inc}}$ se anula, entonces %
% --- JAUA ---
% No vi que mencionaramos que el medio donde stá la partícula es no-absorbente
%
% Está al inicio del párrafo
\begin{equation}
	W_{\text{ext}}=W_{\text{abs}}+W_{\text{sca}}.
	\label{Wext_tot}
\end{equation}
%%

Al asumir que la contribución dominante del campo esparcido es la dipolar, una partícula iluminada por una onda plana $\vb{E}_{\text{inc}} = E_0\,\vu{e}_x$, con $\vb{H}_{\text{inc}} = (k/\mu\omega)\,\vu{e}_z\times\vb{E}_{\text{inc}}$, puede describirse mediante un momento dipolar inducido $\vb{p}$ caracterizado por una polarizabilidad $\alpha$. Por lo tanto, el campo esparcido se describe como una onda esférica de frecuencia angular $\omega$ dada por \cite{bohrenAbsorptionScatteringLight2008}
%
\begin{align}
	\vb{E}_{\text{sca}}&=\frac{e^{ik(r-z)}}{-ikr}\vb{X}\:E_0 \qquad\text{y}\qquad
	\vb{H}_{\text{sca}}=\frac{k}{\omega\mu}\vu{e}_r\times\vb{E}_{\text{sca}},
	\label{eq:EH_s}
\end{align}
%
donde $\vb{X}$ es el vector de amplitud de esparcimiento
%
\begin{equation}
	\vb{X}=\frac{ik^3}{4\pi}\alpha \left[ \vu{e}_r\times(\vu{e}_r\times \vu{e}_x)\right].
	\label{eq:Xvec}
\end{equation}
%
A partir de las Ecs.~\eqref{eq:EH_s} y de las ecuaciones del campo electromagnético incidente ($\vb{E}_{\text{inc}} , \vb{H}_{\text{inc}} $), es posible determinar $W_{\text{ext}}$, que se obtiene a partir de $\langle\vb{S}_{\text{ext}}\rangle_t$, definido en la Ec.~\eqref{eq:S}. Para ello se considera\footnote{En donde de empleó la identidad del triple producto vectorial dada por $\vb{A}\times(\vb{B}\times\vb{C})=\vb{B}(\vb{A}\cdot\vb{C})-\vb{C}(\vb{A}\cdot\vb{B})$~\cite{griffithsIntroductionElectrodynamics2023}.}
%
\begin{align*}
	\vb{E}_{\text{inc}}\times\vb{H}_{\text{sca}}^*&=E_0\vu{e}_x\times\left(\frac{k}{\omega\mu}\vu{e}_r\times \vb{E}_{\text{sca}}\right)^*\\
	&=\frac{k|E_0|^2}{\omega\mu}\frac{e^{-ik(r-z)}}{ikr}[\vu{e}_r(\vu{e}_x\cdot\vb{X}^*)-\vb{X}^*(\vu{e}_x\cdot\vu{e}_r)],
\end{align*}
%
y análogamente,
%
\begin{equation*}
	\vb{E}_{\text{sca}}\times\vb{H}_{\text{inc}}^*=\frac{k|E_0|^2}{\omega\mu}\frac{1}{-ikr}[\vu{e}_z(\vb{X}\cdot\vu{e}_x)-\vu{e}_x(\vb{X}\cdot\vu{e}_z)].
\end{equation*}
%
De esta forma,
%
\begin{align}
	\langle\vb{S}_{\text{ext}}\rangle_t=\frac{1}{2}\:\mbox{Re}&\biggl\{\frac{k|E_0|^2}{\omega\mu}\frac{e^{-ik(r-z)}}{ikr}[\vu{e}_r(\vu{e}_x\cdot\vb{X}^*)-\vb{X}^*(\vu{e}_x\cdot\vu{e}_r)]+ \nonumber\\ 
	&+\frac{k|E_0|^2}{\omega\mu}\frac{1}{-ikr}[\vu{e}_z(\vb{X}\cdot\vu{e}_x)-\vu{e}_x(\vb{X}\cdot\vu{e}_z)]\biggl\}.\label{eq:S_1}
\end{align}
%
Sustituyendo la Ec. \eqref{eq:S_1} en la Ec. \eqref{seq:Wext} y empleando las relaciones $\vb{X}\cdot \vu{e}_r=~0, \vu{e}_r\cdot~\vu{e}_z=~\cos\theta$ y $\vu{e}_r\cdot~\vu{e}_x=\cos\varphi\sin\theta$, se obtiene
%
\begin{align}
		W_{\text{ext}} = -\int_{A}\langle\vb{S}_{\text{ext}}\rangle_t\cdot\vu{e}_r\text{ dA}&=\frac{|E_0|^2k}{\omega\mu}\text{Re}\biggl\{\frac{e^{-ikr}}{ikr}\int_A e^{ikz}\left(\vu{e}_x\cdot\vb{X}^*\right)\dd{a}+ \nonumber\\
		& +\frac{e^{-ikr}}{ikr}\int_A e^{-ikz}\left(\vb{X}\cdot\vu{e}_x\right)\cos\theta\dd{a}+\nonumber\\
		&-\frac{e^{-ikr}}{ikr}\int_Ae^{-ikz}\left(\vb{X}\cdot\vu{e}_z\right)\sin\theta\cos\varphi \dd{a}\biggl\}, \label{eq:Wext_0}
\end{align}
%
que al integrar por partes\footnote{La Ec. \eqref{eq:Wext_0} presenta integrales del tipo
$$\int_{-1}^{1}e^{ikr\mu}f(\mu)\dd{\mu},$$
con $\mu=\cos\theta$, que al integrar por partes dos veces resulta en \cite{bohrenAbsorptionScatteringLight2008}
$$\int_{-1}^{1}e^{ikr\mu}f(\mu)\dd{\mu}=f(\mu)\frac{e^{ikr\mu}}{ikr}\Big|_{-1}^{1}+\frac{1}{(kr)^2}\left[\dv{f}{\mu}e^{ikr\mu}\Big|_{-1}^{1}-\int_{-1}^1\dv[2]{f}{x}\right],$$
y como $kr\ll1$, entonces \cite{bohrenAbsorptionScatteringLight2008}
$$\int_{-1}^{1}e^{ikr\mu}f(\mu)\dd{\mu}\approx \frac{1}{ikr}\left[f(1)e^{ikr}-f(-1)e^{-ikr}\right].$$
}, simplificando se tiene que \cite{bohrenAbsorptionScatteringLight2008}
%
\begin{equation*}
	W_{\text{ext}}=I_{\text{inc}}\frac{4\pi}{k^2}\:\text{Re}\{(\vb{X}\cdot\vu{e}_x)_{\theta=0}\},
\end{equation*}
%
con $I_{\text{inc}}$ la irradiancia de la onda incidente, que corresponde a $|\langle\vb{S}_{\text{inc}}\rangle_t|=k|E_0|^2/(\omega\mu)$, y por consiguiente\footnote{Debido al teorema óptico, la extinción sólo depende de la amplitud de esparcimiento en la dirección de propagación y es el efecto combinado de la absorción en la partícula y el esparcimiento por la partícula en todas las direcciones \cite{bohrenAbsorptionScatteringLight2008}.},
%
\begin{equation}
	C_{\text{ext}}=\frac{W_{\text{ext}}}{I_{\text{inc}}}=\frac{4\pi}{k^2}\:\text{Re}\{(\vb{X}\cdot\vu{e}_x)_{\theta=0}\}, \label{C_ext}
\end{equation}
%
que es la sección transversal de extinción y que posee dimensiones de área.
 Al sustituir las Ecs.~(\ref{eq:EH_s}) en la Ec. (\ref{seq:Ws}) se obtiene
%
\begin{equation}
	C_{\text{sca}}=\int_0^{2\pi}\int_0^{\pi}\frac{|\vb{X}|^2}{k^2}\:\sin\theta\: \text{d}\theta\:\text{d}\phi,
	\label{C_sca_general}
\end{equation}
%%
donde el integrando corresponde a la sección transversal diferencial de esparcimiento, una magnitud que puede determinarse experimentalmente \cite{kinnunenEffectSizeShape2011} y que se denota por $\dd{C_{\text{sca}}}/\dd{\Omega}$ \cite{bohrenAbsorptionScatteringLight2008}. La variable $\Omega$ es el ángulo sólido subtendido por superficie de la partícula, que representa el área de la superficie $\Omega$ de una esfera unitaria cubierta por la proyección de la superficie sobre la esfera~\cite{weissteinAnguloSolido}. %

%
Finalmente, la Ec. (\ref{Wext_tot}) puede ser reescrita como \cite{bohrenAbsorptionScatteringLight2008}
%
\begin{tcolorbox}[ams align]
	C_{\text{ext}}=C_{\text{abs}}+C_{\text{sca}},
	\label{C} 
\end{tcolorbox}
%
\noindent donde $C_{\text{abs}}=W_{\text{abs}}/I_{\text{inc}}$ y $C_{\text{sca}}=W_{\text{sca}}/I_{\text{inc}}$ corresponden a las secciones transversales de absorción y esparcimiento, respectivamente.
Las Ecs. (\ref{C_ext})–(\ref{C}) son propiedades macroscópicas y medibles que proveen información sobre la energía absorbida y esparcida por una partícula \cite{bohrenAbsorptionScatteringLight2008}. Existen diversos métodos experimentales para determinar la sección transversal de extinción, entre ellos la espectroscopía basada en la ley de Beer–Lambert \cite{bohrenAbsorptionScatteringLight2008} y la microscopía holográfica digital \cite{subediMeasurementExtinctionCross2014}. Ambos se fundamentan en medir la superposición entre la onda incidente y la esparcida por la partícula en la dirección de propagación.


	% !TeX root = ../tesis.tex


\section{Relaciones de Kramers-Kronig}
\label{section:rKK}

El estudio del esparcimiento y la absorción de partículas mediante secciones transversales proporciona una descripción macroscópica de su respuesta óptica. Este enfoque se complementa con una caracterización microscópica del material, representada por su función dieléctrica~$\varepsilon(\omega)$~\cite{lucariniKramersKronigRelationsOptical2005}. En el espacio de frecuencias, la función dieléctrica de un material es, en general, una función compleja: su parte real $\text{Re}\{\varepsilon(\omega)\}$ describe cómo la polarización inducida en el material modifica la propagación de la onda electromagnética en el medio, mientras que su parte imaginaria $\text{Im}\{\varepsilon(\omega)\}$ describe la disipación o pérdida de energía debida a procesos de absorción \cite{maierPlasmonicsFundamentalsApplications2007}. Desde un punto de vista físico, $\text{Im}\{\varepsilon(\omega)\}$ se origina en el desfase entre el campo eléctrico incidente y la respuesta inducida de los dipolos del material~\cite{royHilbertTransformBrief2024}. Experimentalmente, es usual realizar mediciones de únicamente una de las dos contribuciones de $\varepsilon(\omega)$ \cite{terdaleIntensitymodulatedOpticalFiber2022,brindzaThinSampleRefractive2014,liuRealPartRefractive2015}. Por ejemplo, técnicas como la espectroscopía de absorción o de transmisión permiten obtener información disipativa, es decir, de $\text{Im}\{\varepsilon(\omega)\}$~\cite{lucariniKramersKronigRelationsOptical2005}. Este tipo de técnicas motiva el uso de las relaciones de Kramers-Kronig (KK), las cuales establecen una conexión  entre las partes real e imaginaria de cualquier función óptica lineal compleja \cite{lucariniKramersKronigRelationsOptical2005}. La base de las relaciones KK reside en el principio de causalidad \cite{lucariniKramersKronigRelationsOptical2005}: la respuesta del sistema no puede preceder temporalmente a la excitación que la genera. En el contexto del esparcimiento de ondas electromagnéticas, esto implica que el campo esparcido no puede existir antes de que la onda incidente alcance el centro de esparcimiento \cite{lucariniKramersKronigRelationsOptical2005}. Este principio, junto con la condición de analiticidad\footnote{Una función compleja $f(x+iy)=u(x,y)+i\,v(x,y)$ es analítica si $u$ y $v$ tienen primeras derivadas continuas y satisfacen las ecuaciones de Cauchy-Riemann \cite{arfkenMathematicalMethodsPhysicists2011}.} de la función óptica en el semiplano superior del plano complejo de frecuencias, garantiza que las partes real e imaginaria de funciones como la susceptibilidad eléctrica, la función dieléctrica, el índice de refracción o la reflectividad no sean independientes; por el contrario, que estén acopladas a través de una transformada de Hilbert \cite{lucariniKramersKronigRelationsOptical2005}. Esto implica que el conocer la respuesta disipativa [$\text{Im}\{\varepsilon(\omega)\}$] en todo el espectro de frecuencias es, en principio, suficiente para determinar de manera unívoca la respuesta óptica asociada al esparcimiento [$\text{Re}\{\varepsilon(\omega)\}$], y viceversa.

En esta sección se considera un material homogéneo e isótropo, caracterizado por una función óptica lineal compleja, en particular la función dieléctrica, definida en un espacio de frecuencias complejo. El objetivo es derivar las relaciones de KK asociadas a dicha función. Para ello, se emplean las transformadas de Fourier directa e inversa de los campos $\vb{E}$ y $\vb{D}$, a partir de las cuales se obtiene una expresión para la permitividad eléctrica relativa. Posteriormente, se imponen condiciones de causalidad sobre esta expresión que permiten  garantizar su analiticidad en el semiplano superior del plano complejo de frecuencias y con ello deducir las relaciones de~KK.

Las transformadas de Fourier temporales directa e inversa de $\vb{D}$ se definen como~\cite{jacksonClassicalElectrodynamics2021}

%
\begin{equation}
	\vb{D}(\vb{r},\omega)=\int_{-\infty}^{\infty}\vb{D}(\vb{r},t')e^{i\omega t'}\text{d}t',\;\;\text{y}\;\;\vb{D}(\vb{r},t)=\frac{1}{2\pi}\int_{-\infty}^{\infty}\vb{D}(\vb{r},\omega)e^{-i\omega t}\text{d}\omega,\label{eq:TFD}
\end{equation}
%
respectivamente. Al sustituir la Ec. para $\vb{D}$ en  \eqref{eq:d3_b3} en la transformada de Fourier inversa de las Ecs. \eqref{eq:TFD}, se obtiene
\vspace{-0.5cm}

%
\begin{equation}
	\vb{D}(\vb{r},t)=\frac{1}{2\pi}\int_{-\infty}^{\infty}\varepsilon(\omega)\vb{E}(\vb{r},\omega)e^{-i\omega t}\text{d}\omega,\label{eq:TF_D_intermedia_1}
\end{equation}
%
y al expresar $\vb{E}(\vb{r},\omega)$ mediante su transformada de Fourier, la Ec. \eqref{eq:TF_D_intermedia_1} se reescribe como
%
\begin{align}
	\vb{D}(\vb{r},t)&=\frac{1}{2\pi}\int_{-\infty}^{\infty}\varepsilon(\omega)e^{-i\omega t}\text{d}\omega\int_{-\infty}^{\infty}e^{i\omega t'}\vb{E}(\vb{r},t')\text{d}t',\\
	&=\frac{1}{2\pi}\int_{-\infty}^{\infty}\int_{-\infty}^{\infty}\varepsilon(\omega)\vb{E}(\vb{r},t') e^{i\omega (t-t')} \text{d}\omega \text{d}t',
	\label{eq:TF_D_intermedia_2}
\end{align}
%
donde además se consideró que los órdenes de integración se pueden intercambiar dado que la integración se realiza en el mismo intervalo ($-\infty$, $\infty$). Para imponer las propiedades de causalidad en una única función respuesta $G(\tau)$, con $\tau = t - t'$, donde $t$ es el tiempo de observación y $t'$ todos los demás, se escribe al campo eléctrico como\footnote{Donde se realiza el cambio de variable $\tau = t - t'$ y se emplea la función constante \cite{arfkenMathematicalMethodsPhysicists2011}
	%
\begin{equation}
	\mathscr{F}[1]=\int_{-\infty}^{\infty}e^{-i\omega\tau}=\delta(\tau),
	\end{equation}
%
con $\delta(\tau)$ la función delta de Dirac, que cumple con $\int_{-\infty}^{\infty}f(\tau)\delta(\tau)\text{d}\tau=f(0)$ \cite{arfkenMathematicalMethodsPhysicists2011}.
}
\begin{equation}
	\vb{E}(\vb{r},t)=\int_{-\infty}^{\infty}\vb{E}(\vb{r},t-\tau)\delta(\tau)=\frac{1}{2\pi}\int_{-\infty}^{\infty}\int_{-\infty}^{\infty}\vb{E}(\vb{r},t-\tau)e^{-i\omega\tau}\text{d}\omega\text{d}\tau,
\end{equation}
%
con $\delta(\tau)$ la delta de Dirac. Al sumar $\vb{E}(\vb{r},t')-\vb{E}(\vb{r},t')$ y multiplicar por $\varepsilon_0/\varepsilon_0$ a la Ec. \eqref{eq:TF_D_intermedia_2}, simplificando se obtiene
%
\begin{equation}
	\vb{D}(\vb{r},t)=\varepsilon_0\left[\vb{E}(\vb{r},t)+\int_{-\infty}^{\infty}G(\tau)\vb{E}(\vb{r},t-\tau)\text{d}\tau\right],\label{eq:TF_D_final} 
\end{equation}
%
\noindent donde $G(\tau)$ es \cite{jacksonClassicalElectrodynamics2021}
%
\begin{equation}
	G(\tau)=\frac{1}{2\pi}\int_{-\infty}^{\infty}\left[\frac{\varepsilon(\omega)}{\varepsilon_0}-1\right]e^{-i\omega\tau}\text{d}\omega,
	\label{eq:G_tdependent} 
\end{equation}
%
\noindent con $\varepsilon(\omega)/\varepsilon_0$ la permitividad eléctrica relativa. Al aplicar la transformada de Fourier a la Ec.~\eqref{eq:G_tdependent}, se obtiene una expresión para la permitividad eléctrica relativa
%
\begin{equation}
	\frac{\varepsilon(\omega)}{\varepsilon_0}=1+\int_{-\infty}^{\infty}G(\tau) e^{-i\omega\tau}\text{d}\tau.
	\label{eq:epsrelativa} 
\end{equation}
%
\noindent Las Ecs. \eqref{eq:TF_D_final} y \eqref{eq:G_tdependent} muestran que el campo de desplazamiento eléctrico al tiempo $t$ depende del campo eléctrico en todos los demás tiempos $t'$. A esta característica, se le conoce como la no localidad temporal entre $\vb{D}$ y $\vb{E}$ \cite{jacksonClassicalElectrodynamics2021}. 
\vspace{0.1cm}

De modo que las Ecs. \eqref{eq:TF_D_final} y \eqref{eq:epsrelativa} sean causales, es necesario imponer condiciones que garanticen que la Ec. \eqref{eq:G_tdependent} se anule para $\tau<0$, lo que se traduce en que al tiempo $t$, únicamente valores del campo eléctrico previos a ese tiempo determinan el vector de desplazamiento eléctrico~\cite{jacksonClassicalElectrodynamics2021}. De esta forma, las Ecs. \eqref{eq:TF_D_final} y \eqref{eq:epsrelativa} se reescriben como
%
\begin{subequations}%
	\begin{tcolorbox}[
		ams align, breakable]
		\vb{D}(\vb{r},t)&=\varepsilon_0\left[\vb{E}(\vb{r},t)+\int_{0}^{\infty}G(\tau)\vb{E}(\vb{r},t-\tau)\text{d}\tau\right], \label{seq:D_final} \\ 
 		\frac{\varepsilon(\omega)}{\varepsilon_0}-1&=\int_{0}^{\infty}G(\tau) e^{-i \omega\tau}\text{d}\tau,\label{seq:G}	\end{tcolorbox}
\end{subequations}
%
\noindent donde $\varepsilon(\omega)$ es una función compleja, cuyas partes real e imaginaria se relacionan como $\varepsilon(-\omega)/\varepsilon_0=\varepsilon^*(\omega^*)/\varepsilon_0$.

Dado que $\vb{E}(\vb{r},t)$ y $\vb{D}(\vb{r},t)$ son funciones reales, $G(\tau)$ también lo es. Si se considera a la Ec. \eqref{seq:G} como una representación de $	\varepsilon(\omega')/\varepsilon_0$ en el plano complejo, entonces la Ec. \eqref{seq:G} es analítica en el semiplano superior complejo siempre que $G(\tau)$ sea finita para toda $\tau$ \cite{jacksonClassicalElectrodynamics2021}. Para que dicha analiticidad se extienda hasta el eje real, es necesario imponer la condición $G(\tau)\rightarrow 0$ cuando $\tau\rightarrow \infty$ \cite{jacksonClassicalElectrodynamics2021} \footnote{Esto se cumple para dieléctricos; en conductores, en cambio, se tiene que $G(\tau)\rightarrow \sigma/\epsilon_0$ cuando $\tau\rightarrow \infty$~\cite{jacksonClassicalElectrodynamics2021}. }. Como $\varepsilon(\omega')/\varepsilon_0-1$ es una función analítica en el semiplano superior complejo incluyendo el eje real, la función $	[\varepsilon(\omega')/\varepsilon_0-1]/(\omega'-\omega)$ también lo es, salvo por la singularidad en $\omega'=\omega$. De acuerdo con el teorema integral de Cauchy\footnote{El teorema integral de Cauchy establece que si una función $f(z)$ es analítica dentro y sobre un contorno cerrado $C$, entonces $\oint_C f(z)/(z-z_0)\dd{z}= 0$ \cite{arfkenMathematicalMethodsPhysicists2011}.}, para un contorno cerrado $C$ que no encierre $\omega_0$ se cumple \cite{arfkenMathematicalMethodsPhysicists2011}
%
\begin{equation}
	\oint_C \frac{[\varepsilon(\omega')/\varepsilon_0-1]}{\omega'-\omega}\dd{\omega'}= 0.
	\label{eq:Cauchy}
\end{equation}
%
Para determinar el contorno apropiado de modo que se pueda emplear la Ec. \eqref{eq:Cauchy}, se considera a $\omega$ un punto sobre el eje real y a $C$ la unión de cuatro curvas con representaciones paramétricas como se muestra en la Fig. \ref{fig:integration_contour0} y que están dadas por \cite{bohrenAbsorptionScatteringLight2008}
%
\begin{align*}
	C_1&: \; \omega'=\Omega, & -A\leq \Omega \leq \omega-a,\\
	C_2&: \; \omega'=\omega-ae^{-i\Omega}, & 0\leq\Omega\leq\pi,\\
	C_3&: \;\omega'=\Omega, & \omega+a\leq\Omega\leq A,\\
	C_4&: \;\omega'=Ae^{i\Omega}, & 0\leq\Omega\leq \pi.
\end{align*}
%
\begin{figure}[h!]
	\centering
	\includegraphics[width=8cm]{1-Esparcimiento/integration.pdf}
	\caption{Contorno de integración $C$ conformado por cuatro curvas.}
	\label{fig:integration_contour0}
\end{figure}
%
Así, la Ec. \eqref{eq:Cauchy} se reescribe como
\begin{align*}
	\int_A^{\omega-a}\frac{[\varepsilon(\Omega)/\varepsilon_0-1]}{\Omega-\omega}\dd{\Omega}&+	\int_{\omega+a}^{A}\frac{[\varepsilon(\Omega)/\varepsilon_0-1])}{\Omega-\omega}\dd{\Omega}+\\
	&+	\int_0^{\pi}\frac{iA\;e^{i\Omega}\; [\varepsilon(A\;e^{i\Omega})/\varepsilon_0-1]}{A\;e^{i\Omega}-\omega}\dd{\Omega}-\int_0^{\pi}i\;[\varepsilon(\omega-ae^{-i\Omega})/\varepsilon_0-1]\dd{\Omega}=0.
\end{align*}
Integrando por partes la Ec. \eqref{seq:G}, se obtiene  
%
\begin{equation*}
	\frac{\varepsilon(\omega)}{\varepsilon_0}-1=\frac{i G^{(0)}(0)}{\omega}-\frac{ G^{(1)}(0)}{\omega^2}+\frac{i G^{(2)}(0)}{\omega^3}+...,
\end{equation*}
%
donde $G^{(i)}(\tau)$ indica la $i$-ésima derivada y donde se empleó que $G(\tau)\rightarrow 0$ cuando $\tau\rightarrow \infty$. Además, dado que $G(\tau)$ es causal, $G(0)=0$; por lo que el primer término se anula al igual que los términos que decaen más rápido que $1/\omega^2$ a frecuencias altas. Con ello, por el lema de Jordan\footnote{El lema de Jordan establece que si una función es continua en un contorno semicircular de radio $R$ centrado en el origen y acotada cuando $R\rightarrow \infty$, entonces la integral sobre el contorno semicircular es 0 \cite{arfkenMathematicalMethodsPhysicists2011}. } la integral sobre el segmento $C_4$ es \cite{bohrenAbsorptionScatteringLight2008}
%
\begin{equation}
	\int_0^{\pi}i\;[\varepsilon(\omega-ae^{-i\Omega})/\varepsilon_0-1]\dd{\Omega}=0.
\end{equation}
%
 Por otro lado, al considerar el límite de $a\to 0$ y al emplear nuevamente el teorema de Cauchy, la integral sobre el segmento $C_2$ es \cite{jacksonClassicalElectrodynamics2021} 
 %
 \begin{equation*}
	\int_0^{\pi}i\;[\varepsilon(\omega-ae^{-i\Omega})/\varepsilon_0-1]\dd{\Omega}=i\pi[\varepsilon(\Omega)/\varepsilon_0-1].
 \end{equation*}
 %
 Finalmente, las integrales sobre los segmentos $C_1$ y $C_3$ equivalen a 
%
\begin{equation}
		\int_A^{\omega-a}\frac{[\varepsilon(\Omega)/\varepsilon_0-1]}{\Omega-\omega}\dd{\Omega}+	\int_{\omega+a}^{A}\frac{[\varepsilon(\Omega)/\varepsilon_0-1])}{\Omega-\omega}\dd{\Omega}=\text{P.V.}\int_{-\infty}^{\infty}\frac{[\varepsilon(\Omega)/\varepsilon_0-1]}{\Omega-\omega} \dd{\Omega},
\end{equation}
%
donde P.V. denota el valor principal de la integral \cite{arfkenMathematicalMethodsPhysicists2011}. Por lo tanto, se obtiene %
%
\begin{equation}
	\PV{\int_{-\infty}^{\infty}\frac{[\varepsilon(\omega')/\varepsilon_0-1]}{\omega'-\omega} \dd{\omega'}}=i\pi[\varepsilon(\omega')/\varepsilon_0-1].
	\label{eq:KK_compleja}
\end{equation}
%
 Al separar la Ec. \eqref{eq:KK_compleja} en parte real e imaginaria, se obtienen las relaciones de KK
%
\begin{align}
	\frac{\text{Re}\{\varepsilon(\omega)\}}{\varepsilon_0}&=1+\frac{1}{\pi}\,\PV{\int_{-\infty}^{\infty}\frac{\text{Im}\{\varepsilon(\omega')\}/\varepsilon_0}{\omega'-\omega}\dd{\omega'}}, \label{seq:ReKK_inf}\\
	\frac{\text{Im}\{\varepsilon(\omega)\}}{\varepsilon_0}&=-\frac{1}{\pi}\,\PV{\int_{-\infty}^{\infty}\frac{[\text{Re}\{\varepsilon(\omega')\}/\varepsilon_0-1]}{\omega'-\omega}\dd{\omega'}}.\label{seq:ImKK_inf}
\end{align}
%


\noindent Matemáticamente, la relación entre las Ecs.~\eqref{seq:ReKK_inf} y~\eqref{seq:ImKK_inf}
se expresa mediante una transformada
de Hilbert\footnote{La transformada de Hilbert $\hat{g}(\omega)$ se define como la convolución de una función $g(\omega)$ con $1/(\pi\omega)$, es decir \cite{hallAdvancedSignalIntegrity2008}
$$\hat{g}(\omega)=g(\omega)\ast\frac{1}{\pi\omega}=\frac{1}{\pi}\int_{-\infty}^{\infty}\frac{g(\omega')}{\omega-\omega'}\dd{\omega'}.$$}. En el dominio armónico, esta relación
puede interpretarse como un desfase\footnote{La transformada inversa de Fourier de $1/(\pi\omega)$ es $i\,\text{sgn}(t)$, donde sgn($\cdot$) denota la función signo. En el plano complejo, multiplicar cualquier número por $i$ equivale a una rotación de $\pi/2$ \cite{hallAdvancedSignalIntegrity2008}.}~de~$\pi/2$~entre los términos asociados al
almacenamiento y a la disipación de energía~\cite{zangwillModernElectrodynamics2013,nussenzveigCausalityDispersionRelations1972}. A partir de que la propiedad de simetría $\varepsilon(-\omega)/\varepsilon_0=\varepsilon^*(\omega^*)/\varepsilon_0$, las Ecs. \eqref{seq:ReKK_inf} y \eqref{seq:ImKK_inf} se reescriben en frecuencias positivas como

%
\begin{subequations}%
	\begin{tcolorbox}[
		ams align, breakable]
		\frac{\text{Re}\{\varepsilon(\omega)\}}{\varepsilon_0}&=1+\frac{2}{\pi}\,\PV{\int_{0}^{\infty}\frac{\omega' [\text{Im}\{\varepsilon(\omega')\}/\varepsilon_0]}{\omega'^2-\omega^2}\dd{\omega'}}, \label{seq:ReKK}\\ 
		\frac{\text{Im}\{\varepsilon(\omega)\}}{\varepsilon_0}&=-\frac{2\omega}{\pi}\,\PV{\int_{0}^{\infty}\frac{[\text{Re}\{\varepsilon(\omega')\}/\varepsilon_0-1]}{\omega'^2-\omega^2}\dd{\omega'}}.
		\label{seq:ImKK}
	\end{tcolorbox}
\end{subequations}
%
\vspace{-0.2cm}









	
	\section{Relaciones de Kramers-Kronig sustractivas}
\label{section:SSKK}
\vspace{-0.4cm}

A pesar de que las relaciones de Kramers–Kronig se han aplicado ampliamente en el análisis de datos experimentales \cite{nakovUnifiedFrameworkNumerical2023}, su implementación práctica enfrenta dos limitaciones fundamentales: los errores experimentales y el intervalo finito de frecuencias en el que suelen obtenerse los datos~\cite{nakovUnifiedFrameworkNumerical2023}. De hecho, cuando las relaciones de KK se evalúan sobre un rango espectral limitado, el error en la reconstrucción de la parte real o imaginaria de la función óptica tiende a incrementarse monótonamente con la frecuencia \cite{miltonFiniteFrequencyRange1997}. Para mitigar este problema, Bachrach y Brown \cite{bachrachExcitonOpticalPropertiesTlBr1970} propusieron las relaciones de Kramers–Kronig sustractivas (SSKK), %
% --- JAUA ---
% nunca explicas por qué SSKK; te lo digo poque en textos más reientes explicas que los acrónimos vienen por sus siglas en inglés.
%
 las cuales incorporan un punto de anclaje experimental es decir, el valor conocido de la función óptica en una frecuencia particular, con el fin de mejorar significativamente la precisión del análisis vía KK \cite{lucariniKramersKronigRelationsOptical2005}. En esta sección se desarrollan las relaciones SSKK para la función dieléctrica considerando un punto de referencia en una frecuencia $\omega_1$, cuyo valor experimental se asume conocido y se puede escribir como %
 % --- JAUA ---
 % Recomendación, decir expxlixitamente que evalues en \omega_q la Ec. XX y YY, dando com oresultado
 %
%
\begin{align}
		\frac{\text{Re}\{\varepsilon(\omega_1)\}}{\varepsilon_0}&=1+\frac{2}{\pi}\,\PV{\int_{0}^{\infty}\frac{\text{Im}\{\varepsilon(\omega')\}/\varepsilon_0}{\omega'-\omega_1}\dd{\omega'}}, \label{seq:omega1_re}\\
		\frac{\text{Im}\{\varepsilon(\omega_1)\}}{\varepsilon_0}&=-\frac{2\omega_1}{\pi}\,\PV{\int_{0}^{\infty}\frac{[\text{Re}\{\varepsilon(\omega')\}/\varepsilon_0-1]}{\omega'-\omega_1}\dd{\omega'}}.\label{seq:omega1_im}	
\end{align}
%
Al restar la Ec. \eqref{seq:omega1_re} de la relación  de KK correspondiente para una $\omega$ abitraria se tiene %
% --- JAUA ---
%
%y acaá ya puedes decir que restas omega1_re de la Ec (KKRe) directamente
%
\begin{align}
	\frac{\text{Re}\{\varepsilon(\omega)\}}{\varepsilon_0}-\frac{\text{Re}\{\varepsilon(\omega_1)\}}{\varepsilon_0}&=\frac{2}{\pi}\PV{\int_{0}^{\infty} \frac{\omega'(\omega'^2-\omega_1^2)(\text{Im}\{\varepsilon(\omega')\}/\varepsilon_0)}{(\omega'^2-\omega^2)(\omega'^2-\omega_1^2)}\dd{\omega'}}+\nonumber\\
	&-\frac{2}{\pi}\PV{\int_{0}^{\infty} \frac{\omega'(\omega'^2-\omega^2)(\text{Im}\{\varepsilon(\omega')\}/\varepsilon_0)}{(\omega'^2-\omega^2)(\omega'^2-\omega_1^2)}\dd{\omega'}},\nonumber\\
	&=\frac{2(\omega^2-\omega_1^2)}{\pi}\,\PV{\int_{0}^{\infty}\frac{\omega'\text{Im}\{\varepsilon(\omega')\}/\varepsilon_0}{(\omega'^2-\omega^2)(\omega'^2-\omega_1^2)}\dd{\omega'}},\label{seq:SSKKRe0}
\end{align}
%

\noindent y de manera análoga para la parte imaginaria, al restar la Ec. \eqref{seq:omega1_im} de la relación  de KK correspondiente para una $\omega$ arbitraria se obtiene
%
\begin{align}
		\frac{\text{Im}\{\varepsilon(\omega)\}}{\varepsilon_0}-\frac{\text{Im}\{\varepsilon(\omega_1)\}}{\varepsilon_0}&=\frac{2}{\pi}\PV{\int_0^{\infty}\frac{\omega(\omega'^2-\omega_1^2)(\text{Re}\{\varepsilon(\omega')\}/\varepsilon_0-1)}{(\omega'^2-\omega^2)(\omega'^2-\omega_1^2)}\dd{\omega'}}+\nonumber\\
		&-\frac{2}{\pi}\PV{\int_0^{\infty}\frac{\omega_1(\omega'^2-\omega^2)(\text{Re}\{\varepsilon(\omega')\}/\varepsilon_0-1)}{(\omega'^2-\omega^2)(\omega'^2-\omega_1^2)}\dd{\omega'}},\nonumber\\
		&=-\frac{2}{\pi}\,\PV{\int_{0}^{\infty}\frac{(\omega-\omega_1)(\omega'^2+\omega\,\omega_1)(\text{Re}\{\varepsilon(\omega')\}-1)/\varepsilon_0}{(\omega'^2-\omega^2)(\omega'^2-\omega_1^2)}\dd{\omega'}}.\label{seq:SSKKIm0}
\end{align}
%
Así, al simplificar las Ecs. \eqref{seq:SSKKRe0} y \eqref{seq:SSKKIm0}, se obtienen las relaciones SSKK \cite{lucariniKramersKronigRelationsOptical2005}
%
\begin{subequations}%
	\begin{tcolorbox}[
		ams align, breakable]
		\frac{\text{Re}\{\varepsilon(\omega)\}}{\varepsilon_0}&=\frac{\text{Re}\{\varepsilon(\omega_1)\}}{\varepsilon_0}+\frac{2(\omega^2-\omega_1^2)}{\pi}\,\PV{\int_{0}^{\infty}\frac{\omega'\text{Im}\{\varepsilon(\omega')\}/\varepsilon_0}{(\omega'^2-\omega^2)(\omega'^2-\omega_1^2)}\dd{\omega'}},\label{seq:SSKK_Re}\\ 
		\frac{\text{Im}\{\varepsilon(\omega)\}}{\varepsilon_0}&=\frac{\text{Im}\{\varepsilon(\omega_1)\}}{\varepsilon_0}+ \nonumber\\
		&\hspace{1cm}-\frac{2(\omega-\omega_1)}{\pi}\,\PV{\int_{0}^{\infty}\frac{[\omega'^2+\omega\,\omega_1][\text{Re}\{\varepsilon(\omega')\}-1]/\varepsilon_0}{(\omega'^2-\omega^2)(\omega'^2-\omega_1^2)}\dd{\omega'}}.\label{seq:SSKK_Im}
	\end{tcolorbox}
\end{subequations}\vspace*{1em}

Las relaciones SSKK permiten mitigar los errores asociados al truncamiento espectral al incorporar, mediante puntos de anclaje, contribuciones provenientes de frecuencias fuera de la ventana de medición \cite{lucariniKramersKronigRelationsOptical2005}. Debido a la propiedad de causalidad de las relaciones KK, las relaciones SSKK preservan las propiedades de analiticidad en el semiplano superior complejo y constituyen un resultado exacto \cite{nussenzveigCausalityDispersionRelations1972}.
	
% !TeX root = ../tesis.tex

\chapter{Eritrocitos y sus patologías}

%
% --- JAUA ---
%  Hay que cambiar un poco la intro acá. Empezar cretomando que tu sistema de interés es un eritrosito y sus propiueddades opticas. Por tanto, en est trabajo presentas su morfologio y cimposiciones.... También debes inlcuir que presentas el esquema de mediciones de las propuedaes opticas que se emplea actualmente en algunas áreas de la medicina. 
%

En este capítulo se estudia la morfología, la composición y las principales alteraciones patológicas de los eritrocitos, y se derivan las funciones dieléctricas correspondientes a distintas concentraciones de hemoglobina a partir de datos reportados en la literatura. En la primera parte se describe la composición y las características generales de la sangre, para posteriormente centrarse en los eritrocitos, analizando su forma, su variabilidad morfológica y los cambios en la concentración de hemoglobina corpuscular media (CHCM) asociados a diversas enfermedades. En la segunda parte se introduce un modelo geométrico del eritrocito basado en óvalos de Cassini, el cual permite representar de manera realista su forma bicóncava \cite{ereminaDifferentShapeModels2006}. A continuación, se presentan y analizan las propiedades ópticas de los eritrocitos y del plasma en el rango ultravioleta–visible (UV-Vis), con énfasis en la dependencia de la función dieléctrica con la CHCM. Se discuten los datos experimentales disponibles para las partes real e imaginaria del índice de refracción de la hemoglobina a distintas concentraciones y su formulación en términos del incremento refractivo específico \cite{friebelModelFunctionCalculate2006}. Con el fin de obtener representaciones analíticas de los datos experimentales, se ajusta un modelo de osciladores de Lorentz a la parte imaginaria de la función dieléctrica  y la parte real se reconstruye empleando las relaciones SSKK. Finalmente, se interpretan las principales características espectrales de la hemoglobina, en particular la banda de Soret en el UV-Vis y las bandas Q en el visible, que dominan la respuesta óptica de los eritrocitos \cite{dayerBandAssignmentHemoglobin2010}.

\section{Morfología y alteraciones}
\label{section:results}


La sangre humana está compuesta principalmente por eritrocitos —también conocidos como glóbulos rojos (RBCs, por sus siglas en inglés)—, leucocitos, plaquetas y plasma sanguíneo~\cite{mescherJunqueirasBasicHistology2024}. El plasma constituye el medio continuo en el que se encuentran suspendidos estos elementos y está formado por agua, electrolitos, proteínas plasmáticas, carbohidratos, lípidos y diversas vesículas extracelulares \cite{mescherJunqueirasBasicHistology2024}. Si bien el plasma presenta absorción y esparcimiento óptico, originados principalmente por proteínas como la albúmina y el fibrinógeno, así como por las plaquetas, su contribución a la respuesta óptica global de la sangre en el rango espectral de 250 nm a 1100 nm es despreciable \cite{meinkeOpticalPropertiesPlatelets2007}. No obstante, bajo ciertas condiciones patológicas, la  absorción del plasma puede incrementarse \cite{bosschaartLiteratureReviewNovel2014}. Dentro de ese mismo rango espectral, la respuesta electromagnética de la sangre está dominada por los eritrocitos, cuya contribución puede superar a la de los demás componentes hasta por tres órdenes de magnitud, debido al contraste entre su índice de refracción y el del plasma \cite{bosschaartLiteratureReviewNovel2014}. Por ello, para describir la interacción luz–sangre, en este trabajo se consideran únicamente a los eritrocitos y al plasma: el primero actúa como la partícula esparcidora  y el segundo como el medio circundante.

 
 Los eritrocitos son células terminalmente diferenciadas\footnote{Una célula terminalmente diferenciada se define como aquella que, al adquirir funciones especializadas, pierde su capacidad de proliferar de manera irreversible. Por ejemplo, los eritrocitos y queratinocitos de los mamíferos, sufren una detención irreversible de su crecimiento al perder su núcleo durante la diferenciación terminal \cite{crescenziReactivationCellCycle2003}.} que carecen de núcleo y están constituidas casi por completo de hemoglobina, la proteína responsable de sus características de absorción óptica \cite{bosschaartLiteratureReviewNovel2014}. A su vez, la hemoglobina presenta firmas espectrales definidas en el rango entre 250–1100 nm, las cuales dependen de su estado químico, principalmente oxihemoglobina ($\text{HbO}_2$) y desoxihemoglobina (Hb) \cite{bosschaartLiteratureReviewNovel2014}. A través de la hemoglobina, los eritrocitos desempeñan su función fisiológica fundamental de transporte de oxígeno \cite{palTextbookMedicalPhysiology2021}. La eficiencia de este proceso se encuentra relacionada con la geometría de los eritrocitos. Estas células poseen una membrana sostenida por un citoesqueleto altamente flexible que les confiere su morfología característica \cite{palTextbookMedicalPhysiology2021}, que se muestra en la micrografía electrónica de barrido (SEM por sus siglas en inglés) mostrada en la Fig.~\ref{subfig:SEM}, tomada de la Ref. \cite{alummoottilScanningElectronAtomic2023}. Como se observa en el diagrama de un corte axial de un eritrocito en la Fig.~\ref{subfig:ery}, los eritrocitos sanos presentan un contorno ovalado y un diámetro promedio de aproximadamente 7.5~$\mu$m. Su forma incluye una depresión central que reduce el espesor en la región media, dando lugar a un disco bicóncavo con un grosor aproximado de 2.6~$\mu$m en los bordes y 0.75~$\mu$m en el centro \cite{palTextbookMedicalPhysiology2021}. Esta geometría no solo es fundamental para su función fisiológica, al proporcionar una alta relación superficie-volumen y facilitar el intercambio de gases \cite{palTextbookMedicalPhysiology2021}, sino que también influye directamente en su comportamiento óptico y en la manera en que esparcen la luz \cite{bosschaartLiteratureReviewNovel2014}.

\begin{figure}[b]
	\centering %
	% --- JAUA ---
	% También has el diagrama de la figura b) en Blender para que no quede desentonada.
	%
	
	\sidesubfloat[First image]{\hspace{-0.2cm}{
	\includegraphics[width=0.35\textwidth]{2-Eritrocitos/SEM_erythrocyte.pdf}\label{subfig:SEM}}}\hspace{0.5cm}
	\sidesubfloat[Second image]{\hspace{-0.5cm}{	\includegraphics[width=0.37\textwidth]{/2-Eritrocitos/erythrocyte.pdf}\label{subfig:ery}}}
	\caption{Estructura de un eritrocito sano. \textbf{a)} Micrografía SEM de un eritrocito de un pacientes sanos. La micrografía muestra un eritrocito con una membrana ligeramente granular; imagen extraída y adaptada de la Ref. \cite{alummoottilScanningElectronAtomic2023}. \textbf{b)} Diagrama de un corte de un eritrocito que muestra las dimensiones de la célula. }
	\label{fig:erythrocytes}
\end{figure}


\begin{figure}[t]
	\centering
	\sidesubfloat[First image]{\hspace{-0.1cm}{
		\includegraphics[width=0.33\textwidth]{2-Eritrocitos/erythrocyte_health1.pdf}\label{subfig:ery_health}}}
	\sidesubfloat[Second image]{\hspace{-0.3cm}{	\includegraphics[width=0.33\textwidth]{2-Eritrocitos/spherocytosis1.pdf}\label{subfig:spherocyte}}}
	
	\vspace{0.3cm}
	
	\sidesubfloat[Third image]{\hspace{-0.1cm}{	\includegraphics[width=0.33\textwidth]{2-Eritrocitos/microcyte1.pdf}\label{subfig:microcyte}}}	\sidesubfloat[Fourth image]{\hspace{-0.3cm}{	\includegraphics[width=0.27\textwidth]{2-Eritrocitos/macrocyte1.pdf}\label{subfig:macrocyte}}}
	
	
	\caption{Morfologías de eritrocitos sanos y patológicos. \textbf{a)} Eritrocito sano con forma bicóncava y radio menor a 3.75~$\mu$m. \textbf{b)} Esferocito: eritrocito con el mismo volumen que un eritrocito sano pero con radio menor. \textbf{c)} Microcito: eritrocito de menor tamaño que uno sano. Conserva la forma bicóncava. \textbf{d)} Macrocito: eritrocito de mayor tamaño que uno sano. Conserva la forma bicóncava.  }
	\label{fig:erythrocytes_diseases}
\end{figure}

La geometría, el tamaño y el índice de refracción de los eritrocitos son susceptibles a alteraciones con relevancia clínica, influyendo tanto en su funcionalidad fisiológica como en su comportamiento óptico \cite{ergulComputationalStudyScattering2010}. Bajo diversas patologías y condiciones hematológicas, los eritrocitos pueden experimentar modificaciones significativas en su morfología y composición. Una de las condiciones más frecuentes es la anemia, definida como una disminución en la capacidad de transporte de oxígeno de la sangre  \cite{keohaneRodaksHematologyClinical2015}. Esta no se considera una enfermedad en sí misma, sino una manifestación clínica de un trastorno o deficiencia subyacente, y puede reflejarse en alteraciones en el número de eritrocitos, su tamaño o su contenido de hemoglobina \cite{keohaneRodaksHematologyClinical2015}. Su clasificación suele basarse en parámetros hematológicos, como la CHCM, que cuantifica la concentración promedio de hemoglobina en el eritrocito \cite{bosschaartLiteratureReviewNovel2014}, y en su morfología, incluyendo el radio y el volumen celular. La Fig.~\ref{subfig:ery_health} presenta un eritrocito sano, con las dimensiones indicadas previamente en la Fig.~\ref{subfig:ery}. Un eritrocito sano presenta valores de CHCM entre 30 y 32 g/dL \cite{saxenaClinicalEvaluationDifferent2018}. En contraste, la Fig.~\ref{subfig:spherocyte} muestra un eritrocito afectado por esferocitosis heriditaria, caracterizado por tener un volumen igual o similar al de un eritrocito sano, pero con una superficie celular reducida debido a defectos en las proteínas responsables de mantener la unión adecuada entre el citoesqueleto y la bicapa lipídica \cite{donatoEsferocitosisHereditariaRevision2015}. Como consecuencia, la célula adquiere una forma esférica con una menor relación superficie–volumen y posee valores de CHCM mayores a 32 g/dL \cite{saxenaClinicalEvaluationDifferent2018}. Esta alteración incrementa la tensión de membrana y la vuelve más frágil e incapaz de adaptarse a variaciones fisiológicas \cite{donatoEsferocitosisHereditariaRevision2015}. La Fig.~\ref{subfig:microcyte} muestra un eritrocito con microcitosis, definido como un eritrocito de tamaño reducido (diámetro menor que $ 6~\mu$m y, a menudo, CHCM menor que $ 30~\text{g/dL}$) \cite{saxenaClinicalEvaluationDifferent2018}. Suelen asociarse con hipocromía y con trastornos que afectan el metabolismo del hierro o la síntesis de hemoglobina \cite{ergulComputationalStudyScattering2010}. Por otro lado, en la Fig.~\ref{subfig:macrocyte} se observa un eritrocito con macrocitosis, identificado por un diámetro mayor a $ 8~\mu\text{m}$ y CHCM de aproximadamente $34~\text{g/dL}$ \cite{saxenaClinicalEvaluationDifferent2018}. Aunque su forma puede seguir siendo bicóncava, su tamaño aumentado suele relacionarse con alcoholismo crónico, enfermedades hepáticas, deficiencias vitamínicas, mieloma o leucemia \cite{ergulComputationalStudyScattering2010}. 
  
Las modificaciones en la geometría (radio celular y volumen) y la composición interna (CHCM) que caracterizan a estas alteraciones, modifican el índice de refracción y las propiedades de esparcimiento óptico de los eritrocitos \cite{bosschaartLiteratureReviewNovel2014}. Por lo tanto, en el estudio de las propiedades ópticas de la sangre en estas condiciones se requiere la implementación de modelos que contemplen con precisión la morfología y la composición interna de los glóbulos rojos patológicos, estableciendo un contraste con las propiedades de los eritrocitos sanos.

\section{Citometría de flujo}
\label{section:flow_cytometry}

La citometría de flujo es una técnica que permite analizar de manera rápida las propiedades físicas de partículas individuales ---generalmente células--- a medida que fluyen frente a uno o varios láseres, mientras se encuentran en suspensión en una solución buffer. Durante el paso de cada célula por el láser, se registra la luz esparcida. Debido a su alta velocidad de adquisición y capacidad para analizar miles de células por segundo, la citometría de flujo constituye una herramienta ampliamente utilizada en hematología, inmunología, biología molecular y diagnóstico clínico.

Cuando una célula atraviesa el haz láser, parte de la radiación incidente es esparcida debido al contraste entre el índice de refracción de la célula y el medio circundante. La distribución angular de la luz dispersada depende de diversos factores, entre ellos el tamaño, la forma, la orientación y la distribución espacial del índice de refracción de la célula. En la mayoría de los citómetros convencionales se registran dos señales de esparcimiento. El esparcimiento frontal (Forward Scatter, FSC), medida en un pequeño intervalo angular alrededor de la dirección de propagación del haz incidente. Tradicionalmente, esta señal se relaciona con el tamaño efectivo de la célula, aunque también depende de su forma y de su índice de refracción y el esparcimiento lateral (Side Scatter, SSC), medida aproximadamente a $90^{\circ}$ especto al haz incidente. Esta señal se asocia comúnmente con la complejidad interna o granularidad celular, ya que es sensible a las variaciones del índice de refracción presentes dentro de la célula. Es importante señalar que los detectores no registran la intensidad dispersada en un único ángulo, sino la luz colectada dentro de un intervalo angular determinado por la geometría de las lentes y las aperturas del sistema óptico.

Un citómetro de flujo está compuesto por tres sistemas principales: el sistema de fluidos, el sistema óptico y el sistema electrónico.
\begin{itemize}
	\item Sistema fluídico: El sistema de fluidos tiene como objetivo transportar las partículas hasta el punto de intersección con el haz láser. Para ello, la muestra celular se inyecta en una corriente de fluido de envoltura (sheath fluid), la cual la confina en una corriente estrecha mediante el principio de enfoque hidrodinámico. Este proceso mantiene el flujo de la muestra centrado dentro del fluido de envoltura y permite que las partículas atraviesen la región de iluminación una a una, evitando que dos células sean analizadas simultáneamente.
	
	En años recientes se han desarrollado citómetros con enfoque acústico, en los cuales ondas ultrasónicas complementan el enfoque hidrodinámico para mejorar la alineación de las células, incrementar la velocidad de adquisición y reducir el riesgo de obstrucción del sistema.
	
	\item Sistema óptico: El sistema óptico está formado por uno o varios láseres, lentes de colección  y detectores. Los láseres iluminan las partículas que atraviesan la región de análisis, mientras que las lentes colectan la luz esparcida. La señal de esparcimiento frontal suele detectarse mediante un fotodiodo, ya que la intensidad de la luz dispersada en esta dirección es considerablemente mayor.
	
	\item Sistema electrónico: El sistema electrónico convierte las señales ópticas registradas por los detectores en señales digitales que posteriormente son procesadas por el software del instrumento. 
\end{itemize}

Las mediciones conjuntas de FSC y SSC permiten representar cada célula mediante un punto en un diagrama bidimensional, facilitando la identificación y separación de diferentes poblaciones celulares. En el caso de los leucocitos, por ejemplo, estas dos señales son suficientes para distinguir las principales subpoblaciones presentes en sangre periférica. En la instrumentación convencional, las señales de dispersión se obtienen mediante detectores colocados en regiones angulares específicas, principalmente en la dirección frontal y alrededor de $90^{\circ}$. Sin embargo, la distribución angular de la luz dispersada contiene información adicional sobre la geometría y las propiedades ópticas de la célula. En consecuencia, la exploración de otros intervalos angulares de detección podría mejorar la capacidad para diferenciar células con distintas morfologías, motivando el análisis desarrollado en este trabajo.
\section{Modelado y propiedades ópticas}
\label{section:optical_properties}

Para modelar la morfología bicóncava característica de los eritrocitos, se han propuesto diversos enfoques analíticos, entre los que destacan los modelos de Skalak \cite{skalakStrainEnergyFunction1973}, Lu \cite{luSimulationsLightScattering2005}, Fung \cite{fungHighresolutionDataGeometry1981},
 y los óvalos de Cassini \cite{ereminaDifferentShapeModels2006}. En este trabajo se emplea el modelo basado en óvalos de Cassini, debido a su simplicidad analítica y su capacidad para reproducir fielmente la morfología promedio de un eritrocito sano, optimizando el costo computacional en simulaciones de esparcimiento \cite{ereminaDifferentShapeModels2006}. Los óvalos de Cassini son curvas planas definidas como el lugar geométrico de los puntos cuyo producto de distancias a dos focos fijos es constante; esto los distingue de las elipses, donde se conserva la suma de dichas distancias \cite{wriedtLightScatteringSingle2006}. Al rotar un óvalo de Cassini alrededor de su eje de simetría, se genera un sólido de revolución con la geometría de un disco bicóncavo. En coordenadas cartesianas, asumiendo simetría respecto al eje $z$, la superficie se describe mediante la expresión \cite{ereminaDifferentShapeModels2006}
 %
 \begin{tcolorbox}[ams align]
 	z(x) = \pm c \sqrt{-a^2 - x^2 \pm \sqrt{4x^2a^2 + b^4}},
 	\label{eq:cassini}
 \end{tcolorbox}
 %
\noindent donde los signos $\pm$ externos definen la simetría respecto al plano XY, mientras que los internos delimitan las ramas del óvalo. Los parámetros geométricos se interpretan de la siguiente manera: $a$ representa la semidistancia focal, $b$ es la raíz cuadrada del producto constante de las distancias a los focos, y $c$ es un factor de escala que ajusta el espesor máximo sin cambiar su radio \cite{wriedtLightScatteringSingle2006}. La relación entre estos parámetros delimita la forma del eritrocito: un incremento en $a$ acentúa la depresión central (bicóncavidad), mientras que valores menores de $a$ derivan en geometrías convexas o incluso esferoidales. Por su parte, $b$ escala las dimensiones globales de la curva, controlando así el volumen total \cite{ereminaDifferentShapeModels2006}. En la Fig. \ref{fig:Cassini_ovals} se muestra un modelo con parámetros $a=2.6~\mu$m, $b=2.7~\mu$m y $c=0.813~\mu$m, los cuales corresponden a un eritrocito sano con un volumen de $80.62$ $\mu$m$^3$. La Fig. \ref{subfig:bidimesional} muestra la sección transversal en donde se identifican en color rojo los dos focos ($F_1,F_2$) del óvalo separados uno del otro a una distancia focal de $2a$. Además, se muestra el valor de $c$ sobre el eje $z$ y se muestran en color gris las distancias entre los puntos focales y un punto arbitrario $P$, cuyo producto determina el valor de $b$. La Fig. \ref{subfig:front} presenta el corte del sólido de revolución resultante.
%
\begin{figure}[t]
	\centering
	\sidesubfloat[First image]{\hspace{-0.5cm}{
		\includegraphics[width=0.53\textwidth]{2-Eritrocitos/Cassini.pdf}\label{subfig:bidimesional}}}\hspace{0.1cm}
	\sidesubfloat[Second image]{\hspace{-0.2cm}{	\includegraphics[width=0.35\textwidth]{2-Eritrocitos/axial_ery.pdf}\label{subfig:front}}}
	\caption{\textbf{a)} Óvalo de Cassini con $a=2.6\,\mu\text{m}, b=2.7\,\mu\text{m}$ y $c=0.813\,\mu\text{m}$ correspondiente a un corte transversal de \textbf{b)} el sólido de revolución que modela a un eritrocito sano con volumen de $80.62\,\mu\text{m}^3$. Los focos $(F_1,F_2)$ se identifican con color azul y se encuentran separados a una distancia de $2a$. Se muestra en color rojo el valor de $c$ para el caso particular sobre el eje $z$. Además, se observan en color gris las distancias entre los puntos focales y un punto arbitrario $P$. El producto de distancias de los focos $F_1$ y $F_2$ a $P$ determinan el parámetro $b$. }	\label{fig:Cassini_ovals}
\end{figure}
%

Para estudiar los procesos de esparcimiento y absorción de la luz en eritrocitos es necesario considerar tanto la descripción morfológica proporcionada por los óvalos de Cassini, que define el volumen ocupado por la partícula, como su respuesta electromagnética, determinada por el contraste entre las propiedades ópticas del eritrocito y las del plasma \cite{bosschaartLiteratureReviewNovel2014}. En el caso de los eritrocitos, su función dieléctrica está modulada la hemoglobina, ya su vez por variaciones de la CHCM, lo que modifica la respuesta óptica de la célula \cite{bosschaartLiteratureReviewNovel2014}. Experimentalmente, los valores del índice de refracción de la hemoglobina se han obtenido a partir de mediciones de transmitancia y reflectancia mediante el uso de esferas integradoras y técnicas espectroscópicas \cite{friebelDeterminationComplexRefractive2005, meinkeOpticalPropertiesPlatelets2007}. No obstante, la mayoría de los datos disponibles en la literatura para distintas concentraciones de hemoglobina se encuentran restringidos a ventanas espectrales finitas.

Los datos experimentales de la parte real e imaginaria del índice de refracción de la oxihemoglobina empleados en este trabajo corresponden a tres valores de la CHCM\footnote{La CHCM cuantifica la concentración total de hemoglobina intracelular y no distingue entre oxihemoglobina y desoxihemoglobina. En consecuencia, el uso de propiedades ópticas de la oxihemoglobina es consistente con la definición de CHCM y con su empleo en modelos ópticos de sangre bajo condiciones fisiológicas promedio \cite{bosschaartLiteratureReviewNovel2014}.}: 28.7~g/dL~\cite{friebelModelFunctionCalculate2006}, 30.6~g/dL \cite{friebelDeterminationComplexRefractive2005} y 34.0~g/dL \cite{friebelDeterminationComplexRefractive2005}. En este trabajo se optó por emplear las propiedades ópticas de la oxihemoglobina, dado que las condiciones hematológicas analizadas se caracterizan principalmente por alteraciones en el tamaño, volumen y concentración de hemoglobina de los eritrocitos, sin implicar necesariamente variaciones en el estado de oxigenación de la hemoglobina~\cite{keohaneRodaksHematologyClinical2015}. Los valores de CHCM considerados fueron seleccionados con el propósito de representar los distintos escenarios patológicos descritos en la sección anterior. La parte real del índice de refracción $\eta_{\text{Hb}}$ para todas las concentraciones se reporta en términos de la relación \cite{friebelModelFunctionCalculate2006}
%
\begin{equation}
	\eta_{\text{Hb}}(\lambda,c_{\text{Hb}})=n_{\text{H}_{2}\text{O}}(\lambda)\,[\beta(\lambda)c_{\text{Hb}}+1],
	\label{eq:IR_eri}
\end{equation}
%
 de $n_{\text{Hb}}$ con $c_{\text{Hb}}$. En cuanto a la parte imaginaria, para la concentración de 28.7 g/dL se dispone de datos en el intervalo espectral de 250 nm a 1100 nm. En contraste, para las concentraciones de 30.6 g/dL y 34 g/dL la información experimental de la parte imaginaria sólo está disponible en el rango de 300 a 600 nm. Para el plasma, se emplearon los datos de la parte real reportados en la Ref. \cite{liuMeasurementRefractiveIndex2019} y de la parte imaginaria en la Ref.~\cite{meinkeOpticalPropertiesPlatelets2007}, ambos en el intervalo espectral comprendido entre 400 nm y 750 nm.
donde $n_{\text{H}_{2}\text{O}}(\lambda)$ es el índice de refracción del agua, $\beta(\lambda)$ es el incremento refractivo específico de la hemoglobina y $c_{\text{Hb}}$ es la concentración de hemoglobina. En la Fig.~\ref{subfig:beta} se muestra la dependencia espectral de $\beta(\lambda)$ en función de la longitud de onda $\lambda$ de la luz incidente, donde los puntos indican los datos experimentales y la línea continua una interpolación. En la región ultravioleta se observa dispersión anómala, mientras que en los intervalos 310–355 nm y 500–1100 nm las variaciones de $\beta$ no son significativas: la desviación estándar del promedio en dichos rangos es del orden de $3.6\times10^{-5}$, comparable con el error experimental reportado para mediciones a una longitud de onda dada \cite{friebelModelFunctionCalculate2006}. En la Fig.~\ref{subfig:refrac} se presentan las curvas de la parte real del índice de refracción obtenidas a partir de la Ec.~\eqref{eq:IR_eri} para las distintas concentraciones: la curva verde para 28.7~g/dL, la azul para 30.6 g/dL y la roja para 34.0 g/dL en función de la longitud de onda $\lambda$ de la luz incidente; a largo de este trabajo se mantendrá el código  de color asociado a la concentración de CHCM. La curva gris punteada representa la parte real del índice de refracción del agua, incluida como referencia. La comparación entre las curvas correspondientes a distintas concentraciones de hemoglobina muestran la dependencia lineal
%
\begin{figure}[t]
	\centering
	\sidesubfloat[First image]{\hspace{-0.5cm}{
			\includegraphics[width=7.5cm]{2-Eritrocitos/beta.pdf}\label{subfig:beta}}}\hspace{0.1cm}
	\sidesubfloat[Second image]{\hspace{-0.4cm}{	\includegraphics[width=7.5cm]{2-Eritrocitos/refractiveIndex.pdf}\label{subfig:refrac}}}\vspace{-0.2cm}
	\caption{Datos empleados reportados por Meinke \textit{et al}. \cite{meinkeOpticalPropertiesPlatelets2007} para el índice de refracción de la hemoglobina a distintas concentraciones como función de la longitud de onda $\lambda$. \textbf{a)}~Incremento refractivo específico $\beta(\lambda)$ de la hemoglobina. Los marcadores corresponden a datos experimentales y la línea continua a una interpolación. \textbf{b)} Parte real del índice de refracción de la hemoglobina $\eta_{\text{Hb}}$ para distintas concentraciones: 28.7 g/dL, 30.6 g/dL y 34.0~g/dL indicados mediante las líneas continuas verde, azul y roja, respectivamente. La línea gris punteada corresponde a la parte real del índice de refracción del agua obtenido de la Ref. \cite{haleOpticalConstantsWater1973}. }
	\label{fig:dataEri}
\end{figure}
%

Con el fin de disponer de una representación continua y analítica de la función dieléctrica en todo el rango del visible, la parte imaginaria de la función dieléctrica de los eritrocitos y del plasma se ajustó mediante un modelo de osciladores de Lorentz\footnote{En este modelo, los electrones del material se describen como osciladores armónicos amortiguados con frecuencias de resonancia características, sometidos a la acción de un campo electromagnético externo \cite{bohrenAbsorptionScatteringLight2008}.}. Esta elección, aunque no es única, resulta particularmente conveniente porque proporciona una expresión analítica que, por construcción, satisface el principio de causalidad \cite{maggioreModernIntroductionClassical2023} y, por lo tanto, es compatible con las relaciones de KK. El ajuste se realizó por mínimos cuadrados entre los datos experimentales y el modelo propuesto, empleando un número finito de osciladores (ver el Apéndice \ref{section:apendix1}). Para identificar las frecuencias de resonancia que corresponden a los osciladores ajustados, en las siguientes figuras se incluye el valor de la energía correspondiente a la posición espectral de cada resonancia. Este valor está determinado por la expresión $\hbar\omega=hc/\lambda$, donde $\hbar=h/2\pi$ con $h$ la constante de Planck y $\lambda$ la longitud de onda de la luz incidente. En la Fig.~\ref{subfig:ajuste_ery} se muestran los datos experimentales correspondientes a la parte imaginaria de la función dieléctrica ($\text{Im}\{\varepsilon\}$) de un eritrocito sano con una CHCM de 30.6 g/dL, extraídos de la Ref. \cite{friebelDeterminationComplexRefractive2005} (puntos de color azul oscuro), junto con el ajuste obtenido mediante osciladores de Lorentz (línea azul claro continua) en función de la longitud de onda $\lambda$ (eje inferior) y la energía $\hbar\omega$ (eje superior) de la luz incidente. De manera análoga, en la Fig.~\ref{subfig:ajuste_plasma} se presentan los datos experimentales de $\text{Im}\{\varepsilon\}$ del plasma (puntos de color violeta oscuro) y su correspondiente ajuste mediante osciladores de Lorentz (línea violeta claro continua). En ambos casos, las líneas grises verticales indican las frecuencias de resonancia asociadas a los osciladores empleados.
%
\begin{figure}[h!]
	\centering
	\sidesubfloat[First image]{\hspace{-0.5cm}{
		\includegraphics[width=7.5cm]{2-Eritrocitos/ajusteLorentz30.pdf}\label{subfig:ajuste_ery}}}\hspace{0.1cm}
	\sidesubfloat[Second image]{\hspace{-0.5cm}{	\includegraphics[width=7cm]{2-Eritrocitos/ajustePlasma.pdf}\label{subfig:ajuste_plasma}}}\vspace{-0.1cm}
	\caption{Parte imaginaria $\text{Im}\{\varepsilon\}$ en función de la longitud de onda $\lambda$ (eje inferior) y la energía $\hbar\omega$ (eje superior) de la luz incidente. \textbf{a)} Eritrocitos con una concentración de 30.6 g/dL. Los datos experimentales se obtuvieron de la Ref.~\cite{friebelDeterminationComplexRefractive2005} y están indicados mediante marcadores de color azul oscuro. La línea azul claro continua representa el ajuste con lorentzianas obtenido.  \textbf{b)} Plasma, cuyos datos experimentales representados por marcadores de color violeta oscuro se extrajeron de la Ref.~\cite{meinkeOpticalPropertiesPlatelets2007}. La línea violeta claro continua representa el ajuste con lorentzianas obtenido. En ambas figuras las líneas grises verticales indican las frecuencias de resonancia asociadas a los osciladores de Lorentz empleados en el ajuste.}
	\label{fig:ajuste}
\end{figure}
%

Dado que los datos experimentales disponibles cubren únicamente ventanas espectrales finitas, los ajustes empleados no consideran transiciones electrónicas en otras regiones por lo que cálculos con KK no estimarían correctamente la parte real calculada a partir de los datos experimentales. Para mitigar este problema se emplearon las relaciones SSKK, las cuales requieren conocimiento del valor de $\text{Re}\{\varepsilon\}$ o $\text{Im}\{\varepsilon\}$ a una frecuencia específica $\omega_1$, lo que permite reducir los errores asociados al truncamiento espectral \cite{lucariniKramersKronigRelationsOptical2005}. En este trabajo se determinó la energía $\hbar\omega_1$ mediante un ajuste minimizando el error absoluto entre la parte real experimental de la función dieléctrica y la obtenida a partir de las relaciones SSKK en el rango espectral de 400 a 600 nm (ver Apéndice \ref{section:apendix2}). Se obtuvieron valores de 2.85 eV, 2.85 eV  y 2.91 eV para las concentraciones de hemoglobina de 28.7~g/dL, 30.6 g/dL y 34.0 g/dL, respectivamente, y de 2.88~eV para el plasma. En la Fig.~\ref{subfig:epsSSKK_eri} se muestra la comparación entre la parte real de la función dieléctrica ($\text{Re}\{\varepsilon\}$) reconstruida mediante SSKK, indicada mediante la línea azul claro continua y los datos experimentales para una CHCM de 30.6~g/dL indicados mediante los puntos azul oscuro, observándose una desviación promedio de 0.014. La línea azul oscuro continua que se observa es únicamente una guía al ojo. De manera análoga, la Fig. \ref{subfig:epsSSKK_Plasma} muestra la parte real de la función dieléctrica obtenida mediante el análisis de SSKK como una línea violeta claro continua donde se observa una desviación promedio de 0.011 respecto a los datos experimentales indicados mediante puntos violeta oscuro (gráfica principal). Asimismo, se muestra una ampliación de la parte real de la función dieléctrica obtenida mediante el análisis de SSKK en el rango de 350 a 650 nm (gráfica interna). La diferencia observada entre lo obtenido mediante SSKK y los datos experimentales, particularmente en el ultravioleta, se debe a que se ha reportado que el índice de refracción de soluciones de hemoglobina y del plasma presenta un incremento pronunciado en el ultravioleta profundo, asociado principalmente a la absorción del agua \cite{bosschaartLiteratureReviewNovel2014}.  Este comportamiento no se reproduce completamente en el presente análisis, ya que los espectros de absorción disponibles para la hemoglobina y el plasma ---y, por ende, la parte imaginaria de la función dieléctrica empleada--- sólo se extienden hasta aproximadamente 250~nm y 300~nm, respectivamente.  A pesar de esta limitación, la aplicación de las relaciones SSKK permite obtener una reconstrucción cuya desviación promedio con respecto a los datos experimentales es casi 1.5 veces menor que la reportada en la Ref. \cite{bosschaartLiteratureReviewNovel2014}, donde se emplea el valor de $\hbar\omega_1$ igual a 1.5 eV [ver Apéndice \ref{section:apendix2}]. En principio, la extensión de los datos experimentales a un rango espectral más amplio \cite{lucariniKramersKronigRelationsOptical2005} o el uso de múltiples frecuencias de anclaje mediante relaciones MSKK\footnote{Las relaciones de Kramers–Kronig generalizadas con sustracciones múltiples (MSKK) permiten introducir varios valores de la función dieléctrica conocidos para mejorar la precisión del análisis \cite{palmerMultiplySubtractiveKramers1998}.} permitiría mejorar la reconstrucción de la función dieléctrica en la región UV.
%
\vspace{-0.1cm}
\begin{figure}[h]
	\centering
		\sidesubfloat[First image]{\hspace{-0.5cm}{
			\includegraphics[width=7.5cm]{2-Eritrocitos/sskk30.pdf}\label{subfig:epsSSKK_eri}}}\hspace{-0.1cm}
	\sidesubfloat[Second image]{\hspace{-0.5cm}{	\includegraphics[width=7.5cm]{2-Eritrocitos/sskkPlasma.pdf}\label{subfig:epsSSKK_Plasma}}}\vspace{-0.2cm}
	\caption{Comparación entre la parte real de la función dieléctrica $\text{Re}\{\varepsilon\}$ reconstruida mediante las relaciones SSKK (línea color claro continua) y los datos experimentales (marcadores color oscuro) en función de la longitud de onda $\lambda$ de la luz incidente para: \textbf{a)} eritrocitos con una CHCM de 30.6 g/dL, a partir de los datos reportados en la Ref. \cite{friebelDeterminationComplexRefractive2005}, donde la línea azul oscuro continua que se observa es una guía al ojo y \textbf{b)} plasma, empleando los datos experimentales de la Ref. \cite{meinkeOpticalPropertiesPlatelets2007} (gráfica principal), donde la línea violeta oscuro continua que se observa es una guía al ojo. Se muestra además una ampliación de la parte real de la función dieléctrica obtenida mediante el análisis de SSKK en el rango de 350 a 650 nm (gráfica interna).}
	\label{fig:epsSSKK}
\end{figure}
%

Cabe destacar que, aunque para la concentración de 28.7 g/dL se dispone de datos experimentales tanto de la parte real como de la imaginaria en un intervalo espectral relativamente amplio, en este trabajo se empleó de manera consistente la función dieléctrica reconstruida mediante el mismo procedimiento (ajuste con osciladores de Lorentz seguido de SSKK) para las tres concentraciones consideradas. Esta elección garantiza que todas las funciones dieléctricas estén sujetas a los mismos efectos de truncamiento espectral, de modo que las comparaciones posteriores de las propiedades ópticas calculadas reflejen exclusivamente la variación de la CHCM y no diferencias metodológicas en el tratamiento de los datos.

En la Fig.~\ref{fig:epsEri} se muestran $\text{Re}\{\varepsilon\}$ [Fig. \ref{subfig:epsReEri}] e $\text{Im}\{\varepsilon\}$ [Fig. \ref{subfig:epsImEri}] en función de la longitud de onda $\lambda$ de la luz incidente para las distintas concentraciones de hemoglobina consideradas: verde para 28.7~g/dL, azul para 30.6 g/dL y rojo para 34.0 g/dL. En todas las curvas destaca la banda de Soret, característica de la hemoglobina, con un máximo situado típicamente en el intervalo 410–430~nm~\cite{bosschaartLiteratureReviewNovel2014}, cuya intensidad es considerablemente mayor que la de las transiciones en el visible. En esta última región se observan las denominadas bandas Q, correspondientes a la banda $\alpha$ (a mayor longitud de onda) y la banda $\beta$ (a menor longitud de onda), cuyos máximos se localizan de manera general en los rangos aproximados 600–650 nm y 520–580 nm \cite{dayerBandAssignmentHemoglobin2010}, respectivamente, dependiendo de la concentración y del estado de oxigenación\footnote{Estas bandas se originan en transiciones electrónicas $\pi \rightarrow \pi^*$ del anillo de porfirina. Tanto la banda de Soret (UV–Vis) como las bandas Q (visible) son sensibles al estado de oxidación del hierro y a cambios conformacionales de la hemoglobina, por lo que caracterizan su estructura y función \cite{dayerBandAssignmentHemoglobin2010}.}. Las regiones sombreadas en la Fig. \ref{fig:epsEri} señalan las bandas espectrales mencionadas anteriormente. De manera adicional, tanto en la parte real como en la imaginaria se aprecia un relieve alrededor de 450 nm, que coincide con el flanco de la banda de Soret y marca la transición entre el régimen de dispersión anómala en el UV-Vis y el comportamiento más suave característico del visible.
%
\begin{figure}[H]
	\centering
	\sidesubfloat[First image]{\hspace{-0.8cm}{
			\includegraphics[width=0.47\textwidth]{2-Eritrocitos/parteReEri.pdf}\label{subfig:epsReEri}}}\hspace{0.1cm}
	\sidesubfloat[Second image]{\hspace{-0.4cm}{	\includegraphics[width=0.47\textwidth]{2-Eritrocitos/parteImEri.pdf}\label{subfig:epsImEri}}}\vspace{-0.2cm}
	\caption{Función dieléctrica de eritrocitos en función de la longitud de onda $\lambda$ de la luz incidente para distintas concentraciones de hemoglobina, indicadas mediante líneas continuas verde, azul y roja correspondientes a 28.7 g/dL, 30.6~g/dL y 34 g/dL, respectivamente. Las regiones sombreadas señalan las bandas espectrales que dominan la respuesta óptica de los eritrocitos en el UV-Vis. \textbf{a)} Parte real de la función dieléctrica. \textbf{b)} Parte imaginaria de la función dieléctrica.}
	\label{fig:epsEri}
\end{figure}
%

% !TeX root = ../tesis.tex

\chapter{\textit{Finite Integration Technique}}
\label{chapter:FIT}

\vspace*{7em}

%
% --- JAUA ---
% Igual que en pasado: Primero menciona que debes calcular las porpiedades del eritroxito pweo que no se cuenta con una solución analítica, por lo tanto recurres a métpdps niméricos.
%

Debido a la geometría bicóncava de los eritrocitos, su respuesta electromagnética no puede describirse mediante una solución analítica, por lo que se aborda de manera numérica. En este capítulo se presenta la formulación del método de \textit{Finite Integration Technique} (FIT) y su aplicación al estudio del esparcimiento de luz por una esfera. En primer lugar, se describe el método a partir de las ecuaciones de Maxwell en su forma integral y su discretización sobre mallas duales. Se introducen las ecuaciones de malla de Maxwell, junto con las relaciones constitutivas discretas que fundamentan el esquema numérico, así como su formulación en el dominio del tiempo. Posteriormente, se analiza la convergencia de las simulaciones realizadas en \textit{Computer Simulation Technology} (CST) mediante la comparación con la solución analítica de la teoría de Mie \cite{bohrenAbsorptionScatteringLight2008}. Este análisis se lleva a cabo evaluando la influencia de los principales parámetros numéricos —como el tamaño del dominio computacional, el mallado y las condiciones de frontera— sobre la precisión y el costo computacional. A partir de estos resultados, se establecen los criterios de discretización que se emplean como base en las simulaciones presentadas en los capítulos siguientes.




\section{Ecuaciones de malla de Maxwell}
\label{section:MGE}

El método \textit{Finite Integration Technique} (FIT) se basa en la discretización espacial de las ecuaciones de Maxwell a partir de su forma integral, obteniendo así las Ecuaciones de malla de Maxwell \cite{clemensDiscreteElectromagnetismFinite2001}. En este método las magnitudes electromagnéticas se asocian a los objetos geométricos de la malla sobre los cuales se definen dichas integrales \cite{dassaultsystemesCSTMicrowaveStudio2003}.

El primer paso del proceso de discretización empleado en FIT consiste en definir el mallado de integración. Para ello, se emplea una discretización tipo Yee \cite{yeeNumericalSolutionInitial1966}. En este contexto, el problema electromagnético ---generalmente formulado en una región con frontera abierta---  se restringe a un dominio computacional finito $\Omega\in\mathbb{R}^3$. Posteriormente, dicho dominio se descompone en un número finito de subregiones $V_i$, tales como tetraedros o hexaedros, que encajan exactamente entre sí. Esto implica que la intersección entre dos celdas distintas  es vacía o corresponde a un elemento geométrico compartido de menor dimensión (una cara, una arista o un vértice). Esta descomposición da lugar a un sistema coordenado ortogonal $G$, que constituye la malla computacional del problema~\cite{clemensDiscreteElectromagnetismFinite2001}. Para simplificar el desarrollo del método, se asume que $\Omega$ tiene forma paralelepipédica y que su discretización se realiza mediante una malla cartesiana obtenida como producto tensorial de subdivisiones unidimensionales, de modo que \cite{hadabSolutionMaxwellsEquation2023}
%
\begin{align}
	G=\{(x_i,y_j,z_k); i\in\{1,...,I\};j\in\{1,...,J\};k\in\{1,...,K\}\}.
\end{align}
%
Los nodos de la malla $(x_i,y_j,z_k)$ se indexan mediante los índices $i,j,k$ a lo largo de los ejes $x,y,z$. De esta forma, el mallado contiene un total de $N_p=IJK$ nodos y $(I-1)(J-1)(K-1)$ celdas \cite{hadabSolutionMaxwellsEquation2023}. Para facilitar la implementación numérica, los nodos se enumeran mediante un único índice $n$ definido como \cite{hadabSolutionMaxwellsEquation2023}
%
\begin{equation}
	n=1+(i-1)M_x+(j-1)M_y+(k-1)M_z,
	\label{eq:index_system}
\end{equation}
% 
donde $M_x=1,M_y=I$ y $M_z=IJ$. En la Fig. \ref{fig:G_complex} se muestra el sistema coordenado $G$ para un conjunto de celdas cartesianas. Asimismo, se muestran las componentes de los nodos en los extremos de la malla, así como la identificación del volumen $V_n$ y el área $S_{u,m}$ asociados a los elementos vinculados a los nodos $n$ y $m$, respectivamente.
%
\begin{figure}[h!]
	\centering
	\includegraphics[width=8cm]{3-FIT/G_complex.pdf}
	\caption{Representación de una malla tridimensional de tipo Yee ($G$) en un sistema de coordenadas cartesiano. Se muestra la discretización del volumen de interés donde cada vértice actúa como un nodo $n$, con $N = I \times J \times K$ nodos totales vinculados mediante el sistema de indexación dado por la Ec. \eqref{eq:index_system}. }
	\label{fig:G_complex}
\end{figure}
%


Una vez definido $G$, el método FIT se formula a partir de las ecuaciones de Maxwell en su forma integral que, en unidades del SI, están dadas por \cite{griffithsIntroductionElectrodynamics2023}

%
\begin{subequations} \label{eqs:Maxwell_int}
	\begin{tcolorbox}[
		ams align, breakable]
		\oint_{S}\vb{D} \cdot \dd{\vb{a}}&=Q_{\text{ext}}, &\mbox{(Ley de Gauss eléctrica)}  
		\label{seq:GE_int} \\
		\oint_{S}\vb{B} \cdot \dd{\vb{a}} &= 0,						&\mbox{(Ley de Gauss magnética)}   
		\label{seq:GM_int} \\
		\oint_{\partial S}\vb{E} \cdot d\vb{l} &= -\frac{\text{d}}{\dd{t}}\int_{S}\vb{B}\cdot\dd{\vb{a}}, 	&\mbox{(Ley de Faraday-Lenz)}		
		\label{seq:FL_int}\\
		\oint_{\partial S}\vb{H} \cdot d\vb{l} &= \int_{S}\left(\vb{J}_{\text{ext}}+\frac{\partial \vb{D}}{\partial t}\right)\cdot\dd{\vb{a}}, &
		\mbox{(Ley de Ampère-Maxwell)} \label{seq:AM_int}
\end{tcolorbox}\end{subequations}\noindent
%
donde se obvian las dependencias espaciales y temporales; $\rho_{\text{ext}}$ es la densidad volumétrica de carga externa y  $\vb{J}_{\text{ext}}$ es la densidad volumétrica de corriente  externa. A partir de las ecuaciones de Maxwell en su forma integral, se construye la discretización del método FIT. En particular, al aplicar la ley de Faraday–Lenz [Ec. \ref{seq:FL_int}] sobre una cara $S_{z,n}$ asociada a una celda con volumen $V_n$ se obtiene~\cite{hadabSolutionMaxwellsEquation2023}
%
\begin{equation}
	\oint_{\partial S}\vb{E} \cdot d\vb{l}=(x_{i+1}-x_i)E_{x,n}+(y_{j+1}-y_j)E_{y,n+1}+(x_{i+1}-x_i)E_{u,n+I}+(y_{j+1}-y_j)E_{y,n},
	\label{eq:lineint_Yee}
\end{equation}
%
donde la integral de línea se evalúa sobre los cuatro segmentos que conforman el contorno de la cara $S_{z,n}$. Por consiguiente, la integral de área considera el producto de los segmentos 
%
\begin{equation}
	\int_{S}\vb{B}\cdot\dd{\vb{a}}=B_{z,n}(x_{i+1}-x_i)(y_{j+1}-y_j),
	\label{eq:surfint_Yee}
\end{equation}
%
de modo que la Ec. \eqref{seq:FL_int} se reescribe como %
%
\begin{align}
	(x_{i+1}-x_i)E_{x,n}+(y_{j+1}-y_j)E_{y,n+1}&+(x_{i+1}-x_i)E_{u,n+I}+ \nonumber\\
	&+(y_{j+1}-y_j)E_{y,n}
	=-\frac{\dd}{\dd{t}}B_{z,n}(x_{i+1}-x_i)(y_{j+1}-y_j).
	\label{eq:Faraday_grid}
\end{align}
%
En la Fig.~\ref{subfig:FIT_cell_electric} se muestra la superficie $S_{z,n}$ de una celda en la cual los vectores en rojo y azul oscuro representan, respectivamente, la circulación del campo eléctrico a lo largo de las aristas y el flujo magnético que atraviesa la superficie. 
%
\begin{figure}[tb]
	\centering
	\sidesubfloat[First image]{\hspace{-0.8cm}{
			\includegraphics[width=7.8cm]{3-FIT/FIT_cell.pdf}\label{subfig:FIT_cell_electric}}}\hspace{0.1cm}
	\sidesubfloat[Second image]{\hspace{-0.5cm}{	\includegraphics[width=7.5cm]{3-FIT/FIT_cell_fluxB.pdf}\label{subfig:FIT_cell_magnetic}}}\vspace{-0.1cm}
	\caption{\textbf{a)} Cara $S_{z,n}$ de la celda $V_{n}$ del complejo $G$ en la que las flechas rojas indican la circulación del campo eléctrico a lo largo de las aristas que conforman el contorno de la cara, mientras que la flecha azul representa el flujo del campo magnético a través de dicha superficie, asociado a $B_{z,n}$. \textbf{b)} Celda $V_{n}$ del complejo $G$ que muestra el flujo de campo magnético que atraviesa cada una de sus superficies. Los subíndices asociados a cada flujo corresponden a la superficie de la cara que atraviesan.}
	\label{fig:FIT_cell}
\end{figure}
%

La formulación anterior es válida para cualquier cara de $G$ y puede extenderse a superficies formadas por la unión de varias caras, ya que la suma de las circulaciones sobre los contornos individuales equivale a la circulación sobre el contorno total de la superficie \cite{arfkenMathematicalMethodsPhysicists2011}. Con el fin de extender esta relación a todo el dominio, se introduce una formulación matricial que permite representar de manera compacta tanto las magnitudes electromagnéticas como los operadores integrales discretos.
En este marco, los campos electromagnéticos se organizan como vectores de dimensión $3N$. En particular, los campos $\vb{E}, \vb{D}, \vb{B}$ y $\vb{H}$ se representan mediante los vectores $\vb{e}, \vb{d}, \vb{b}$ y $\vb{h}$, respectivamente, definidos como \cite{hadabSolutionMaxwellsEquation2023}
\begin{equation}
	\vb{e}=\begin{bmatrix}
		E_{x,1}\\
		E_{x,2}\\
		\vdots\\
		E_{z,N} 
	\end{bmatrix},\;\;
	\vb{d}=\begin{bmatrix}
		D_{x,1}\\
		D_{x,2}\\
		\vdots\\
		D_{z,N} 
	\end{bmatrix},\;\;
	\vb{b}=\begin{bmatrix}
		B_{x,1}\\
		B_{x,2}\\
		\vdots\\
		B_{z,N} 
	\end{bmatrix},\;\;
	\vb{h}=\begin{bmatrix}
		H_{x,1}\\
		H_{x,2}\\
		\vdots\\
		H_{z,N} 
	\end{bmatrix},
\end{equation}
%
y de manera análoga se definen los vectores $\vb{j}$ y $\vb{q}$, correspondientes a la discretización de $\vb{J}_{\text{ext}}$ y $Q_{\text{ext}}$, respectivamente. Para construir los operadores integrales discretos, se introducen las matrices $P_u$, con $u=\{x,y,z\}$, que codifican la conectividad de la malla. Estas matrices se definen de forma general como \cite{hadabSolutionMaxwellsEquation2023}
%
\begin{equation}
	(P_u)_{m,p}=
	\begin{cases}
		-1, & \text{si } p=m,\\
		1, & \text{si } p=m+M_u,\\
		0, & \text{en cualquier otro caso},
	\end{cases}
	\qquad (m,p)\in \{1,\dots,N\}^2.
\end{equation}
A partir de estas, se construye el operador rotacional $\mathbb{C}\in\mathbb{R}^{3N\times 3N}$, dado por \cite{hadabSolutionMaxwellsEquation2023}
%
\begin{equation}
	\mathbb{C}=\begin{pmatrix}
		0 & -P_z & P_y\\
		P_z & 0 & -P_x\\
		-P_y & P_x & 0
	\end{pmatrix}.
\end{equation}
Adicionalmente, se introducen matrices diagonales que contienen la información geométrica del mallado. En particular, la matriz $\mathbb{D}_l$ agrupa las longitudes asociadas a las aristas, mientras que la matriz $\mathbb{D}_S$ contiene las áreas de las caras. La primera se define como \cite{hadabSolutionMaxwellsEquation2023}
%
\begin{equation}
	\mathbb{D}_l=\text{diag}(\Delta x_{1},\dotso,\Delta x_{N},\Delta y_{1},\dotso,\Delta y_{N},\Delta z_{1},\dotso,\Delta z_{N}),
\end{equation}
%
donde las longitudes asociadas a cada dirección se definen de manera general como \cite{hadabSolutionMaxwellsEquation2023}
%
\begin{equation}
	\Delta u_n =
	\begin{cases}
		\displaystyle \int_{u_\ell}^{u_{\ell+1}} \left\lVert \frac{\partial \vb{r}}{\partial u} \right\rVert \dd{u}, & \text{si } \ell < N_u,\\
		0, & \text{si } \ell = N_u,
	\end{cases}
\end{equation}
%
con $\ell=\{i,j,k\}$ tal que $u_\ell$ es la coordenada correspondiente ($x_i$, $y_j$ o $z_k$), y $N_u$ el número total de nodos en la dirección $u$ (es decir, $N_x=I$, $N_y=J$, $N_z=K$) \cite{hadabSolutionMaxwellsEquation2023}. De manera análoga, la matriz de áreas está dada por \cite{hadabSolutionMaxwellsEquation2023}
%
\begin{equation}
	\mathbb{D}_S=\text{diag}(S_{x,1},\dotso,S_{x,N}, S_{y,1},\dotso,S_{y,N},S_{z,1},\dotso,S_{z,N}),
\end{equation}
%
con $S_{x,n}=\Delta y_n \Delta z_n, S_{y,n}=\Delta x_n \Delta z_n$ y $S_{z,n}=\Delta x_n \Delta y_n$ \cite{hadabSolutionMaxwellsEquation2023}. Finalmente, la ley de Faraday–Lenz (Ec. \eqref{seq:FL}) empleando la discretización tipo Yee se reescribe en forma matricial como \cite{hadabSolutionMaxwellsEquation2023}
%
\begin{equation}
	\mathbb{C}\mathbb{D}_l\vb{e}=-\mathbb{D}_S\frac{\partial}{\partial t}\vb{b}.
\end{equation}
%

El segundo operador a discretizar es el operador de la divergencia. Para ello, se parte de la Ec. \eqref{seq:GM_int} aplicada a una celda de volumen $V_n$, como se muestra en la Fig.~\ref{fig:FIT_cell}, donde las flechas azules representan los flujos magnéticos asociados a cada una de las caras de la celda. Al aplicar la ley de Gauss magnética siguiendo el esquema de Yee (Ec. \ref{eq:surfint_Yee}), la ecuación integral puede reescribirse como una suma de los flujos magnéticos que atraviesan las seis caras de la celda, dando lugar a~\cite{hadabSolutionMaxwellsEquation2023}
%
\begin{align}
	B_{z,n}(x_{i+1}-x_i)(y_{j+1}-y_j)
	&-B_{z,n+1}(x_{i+1}-x_i)(y_{j+1}-y_j)\nonumber\\
	&+B_{x,n}(y_{j+1}-y_j)(z_{k+1}-z_k)\nonumber\\
	&-B_{x,n+1}(y_{j+1}-y_j)(z_{k+1}-z_k)\nonumber\\
	&+B_{y,n}(x_{i+1}-x_i)(z_{k+1}-z_k)\nonumber\\
	&-B_{y,n+1}(x_{i+1}-x_i)(z_{k+1}-z_k)=0.
	\label{eq:GM-Yee}
\end{align}
%
La Ec. \eqref{eq:GM-Yee} puede escribirse en forma matricial para extenderse en todo el dominio mediante la introducción del operador de divergencia \cite{hadabSolutionMaxwellsEquation2023}
%
\begin{equation}
	\mathbb{S} = \begin{pmatrix} P_x & P_y & P_z \end{pmatrix},
\end{equation}
%
el cual actúa sobre el vector de campo magnético ponderado por las áreas de las caras. De este modo, la ecuación discreta de la ley de Gauss magnética se expresa como \cite{hadabSolutionMaxwellsEquation2023}
%
\begin{equation}
	\mathbb{S}\mathbb{D}_S\vb{b} = 0.
\end{equation}
%%

La discretización de las ecuaciones de Maxwell restantes requiere la introducción de un segundo sistema coordenado $\tilde{G}$, dual al complejo primario $G$ \cite{hadabSolutionMaxwellsEquation2023}
%
\begin{align}
	\tilde{G}=\{(\tilde{x}_i,\tilde{y}_j,\tilde{z}_k); i\in\{1,...,I\};j\in\{1,...,J\};k\in\{1,...,K\}\}.
\end{align}
%
 Este complejo dual se construye tomando los centros de las celdas de $G$ como vértices de las celdas de $\tilde{G}$. Con esta definición se obtiene una correspondencia uno a uno entre las aristas de $G$ y las caras de $\tilde{G}$, y viceversa, con la malla dual desplazada media celda respecto a la malla primaria \cite{clemensDiscreteElectromagnetismFinite2001}. En consecuencia, los nodos de la malla dual ($\tilde{x}_i,\tilde{y}_j,\tilde{z}_k$) se definen como los puntos medios entre nodos consecutivos de la malla primaria \cite{hadabSolutionMaxwellsEquation2023}
 %
 \begin{equation}
 	 \tilde{u}_{\ell} =\dfrac{u_{l+1}+u_l}{2}, \begin{cases}
 		 l \in \{1,...,N_u-1\}\\
 		\tilde{u}_{N_u}=0
 	\end{cases}.
 \end{equation}
 %
 En la Fig. \ref{fig:FIT_cell_dual} se muestra la relación de ortogonalidad entre una celda del complejo primario $G$ (celda amarilla) y su correspondiente celda dual $\tilde{G}$ (celda gris). Además, se muestran las cantidades asociadas a cada complejo ($\vb{E}$ y $\vb{B}$ asociados a $G$, mientras que $\vb{D}$ y $\vb{H}$ asociados a $\tilde{G}$) y su localización espacial en cada celda.
 
 %
 \begin{figure}[h!]
 	\centering
 	\includegraphics[width=8cm]{3-FIT/FIT_dual_grid.pdf}
 	\caption{Celda $V_{n}$ del complejo $G$ y celda dual $\tilde{V}_{n}$ del complejo $\tilde{G}$ desplazada media celda en todas las direcciones y ortogonal a la primera. Se muestran las cantidades asociadas a una cara $S_{z,n}$ y $\tilde{S}_{\tilde{z},n}$ del complejo $G$ y $\tilde{G}$, respectivamente. Se muestran en azul las cantidades relacionadas con los campos magnéticos y en rojo con los campos eléctricos.}
 	\label{fig:FIT_cell_dual}
 \end{figure}
 %
 \noindent De manera análoga al caso del complejo primario, se introducen matrices diagonales que contienen la información geométrica de la malla dual. En particular, la matriz dual de longitudes de integración  $\tilde{\mathbb{D}}_{\ell}$ se define como \cite{hadabSolutionMaxwellsEquation2023}
 %
 \begin{equation}
 	 \tilde{\mathbb{D}}_{\ell}=\text{diag}(\Delta\tilde{x}_{1},\dotso,\Delta\tilde{x}_{N},\Delta\tilde{y}_{1},\dotso,\Delta\tilde{y}_{N},\Delta\tilde{z}_{1},\dotso,\Delta\tilde{z}_{N}),
 \end{equation}
 %
 donde las longitudes asociadas a cada dirección se definen de manera general como \cite{hadabSolutionMaxwellsEquation2023}
 %
 \begin{equation}
 	\Delta \tilde{u}_n =
 	\begin{cases}
 		\dfrac{\Delta u_n+ \Delta u_{n-M_u}}{2}, &  2<\ell < N_u-1,\\
 		\dfrac{\Delta u_n}{2}, & \ell = 1,\\
 		\dfrac{\Delta u_{n-M_u}}{2}, & \ell=M_u.
 	\end{cases}
 \end{equation}
 %
De forma similar, la matriz de áreas de la malla dual $\tilde{\mathbb{D}}_S$ está dada por \cite{hadabSolutionMaxwellsEquation2023}
 %
 \begin{equation}
 	\mathbb{D}_S=\text{diag}(\tilde{S}_{x,1},\dotso,\tilde{S}_{x,N}, \tilde{S}_{y,1},\dotso,\tilde{S}_{y,N},\tilde{S}_{z,1},\dotso,\tilde{S}_{z,N}),
 \end{equation}
 %
 con $\tilde{S}_{x,n}=\Delta \tilde{y}_n \Delta \tilde{z}_n, \tilde{S}_{y,n}=\Delta \tilde{x}_n \Delta \tilde{z}_n$ y $\tilde{S}_{z,n}=\Delta \tilde{x}_n \Delta \tilde{y}_n$ \cite{hadabSolutionMaxwellsEquation2023}. 
 
 Los operadores integrales discretos en la malla dual se relacionan directamente con los del complejo primario. En particular, el operador rotacional dual $\tilde{\mathbb{C}}$ y el operador divergencia dual $\tilde{\mathbb{S}}$ están dados por \cite{hadabSolutionMaxwellsEquation2023}
 %
 \begin{align}
 	\tilde{\mathbb{C}}&=\mathbb{C}^T\\
 	\tilde{\mathbb{S}}&=(P_x^T,P_y^T,P_z^T).
 \end{align}
 %
 De esta forma, para el par de sistemas $\{G,\tilde{G}\}$, el conjunto completo de ecuaciones matriciales discretas —conocidas como \textit{Maxwell Grid Equations} (MGE)— se escribe como \cite{hadabSolutionMaxwellsEquation2023}
%
\begin{subequations} \label{eqs:MGE}
	\begin{tcolorbox}[
		ams align, breakable]
	\tilde{\mathbb{S}}\tilde{\mathbb{D}}_S \vb{d}&= \vb{q}, &\mbox{(Ley de Gauss eléctrica)}  
		\label{seq:GE_MGE} \\
	\mathbb{S}\mathbb{D}_S\vb{b}&=0,						&\mbox{(Ley de Gauss magnética)}   
		\label{seq:GM_MGE} \\
		\mathbb{C}\mathbb{D}_l\vb{e}&=-\mathbb{D}_S\frac{\partial}{\partial t}\vb{b},	&\mbox{(Ley de Faraday-Lenz)}		
		\label{seq:FL_MGE}\\
		\tilde{\mathbb{C}}\tilde{\mathbb{D}}_l\vb{h} &= \tilde{\mathbb{D}}_S \left( \frac{\text{d}}{\dd{t}}\vb{d}+\vb{j}\right). &
		\mbox{(Ley de Ampère-Maxwell)} \label{seq:AM_MGE}
\end{tcolorbox}\end{subequations}\noindent
%

Además de las MGE, el método FIT incorpora las relaciones constitutivas discretas, que describen la respuesta del material. Estas relaciones se escriben en forma matricial como \cite{hadabSolutionMaxwellsEquation2023}
%
\begin{equation}
	\vb{d}=\tilde{\mathbb{D}}_{\varepsilon}\vb{e},\;\;\; \text{y}\;\;\;\vb{b}=\mathbb{D}_{\mu}\vb{h}
\end{equation}
%
donde $\tilde{\mathbb{D}}_{\varepsilon}$  y $\mathbb{D}_{\mu}$ son las matrices discretas de permitividad y la permeabilidad, respectivamente. Para medios isótropos, estas matrices son diagonales \cite{clemensDiscreteElectromagnetismFinite2001}.


\input{3-FIT/2-Leap-frog Scheme}
\section{Convergencia}
\label{section:convergence}

Las simulaciones se realizaron empleando el software \textit{Computer Simulation Technology} (CST, versión 2025) en una computadora de escritorio equipada con un procesador Intel Core i7 de $13^a$ generación (16 núcleos, 24 hilos) y 64 GB de memoria RAM, bajo el sistema operativo Windows 11 Home. El análisis de convergencia se llevó a cabo mediante la comparación entre la solución analítica para una esfera, descrita por la teoría de Mie \cite{bohrenAbsorptionScatteringLight2008}, y los resultados numéricos obtenidos con CST. Para este propósito, se consideró una esfera de radio 2,680 nm, cuya función dieléctrica corresponde a la de un eritrocito con una CHCM de 30.6 g/dL, inmersa en plasma con $\varepsilon_{\text{Plasma}}=~1.84$. La esfera se iluminó con una onda plana propagándose en la dirección $z$ y polarizada en $x$, considerando longitudes de onda en el rango de 400 a 600 nm. En CST se evaluó la influencia de distintos parámetros numéricos sobre la convergencia de la solución: el tamaño del dominio computacional ---definido por la distancia entre la partícula y la caja de simulación (\textit{bounding box})---, el nivel de reflexión estimado (\textit{estimated reflection level}) ---relacionado con la capa absorbente Perfect Matching Layer (PML)---, así como el mallado global y el mallado local (tanto en volumen como en superficie). Con este propósito, se realizó una variación paramétrica para identificar la combinación de parámetros que proporcionara la mejor concordancia con la solución analítica de la teoría de Mie. En la Fig.~\ref{subfig:CST_matriz} se muestra la geometría del problema, donde la esfera se encuentra inmersa en el dominio delimitado por la caja de simulación. Por otro lado, en la Fig.~\ref{subfig:CST_mesh} se muestran algunos de los parámetros de discretización considerados. El mallado global corresponde a la discretización de todo el dominio (matriz y partícula), mientras que, en caso de emplearse un refinamiento local, este se aplica exclusivamente a la partícula. Asimismo, se define la distancia $h$ como la separación mínima entre la superficie de la partícula y la caja de simulación.

\begin{figure}[h]
	\centering
	\sidesubfloat[First image]{{
			\includegraphics[width=5cm]{3-FIT/Matriz.pdf}\label{subfig:CST_matriz}}}\hspace{0.1cm}
	\sidesubfloat[Second image]{{	\includegraphics[width=4.55cm]{3-FIT/Mesh.pdf}\label{subfig:CST_mesh}}}\vspace{-0.1cm}
	\caption{\textbf{a)} Vista tridimensional del dominio computacional en CST.  \textbf{b)} Corte transversal del dominio computacional mostrando el mallado en el plano $x=0$. Se identifica
	el mallado global como todo lo que rodea a la esfera, el mallado local como lo que compone a la esfera y la distancia $h$ entre el borde de la esfera y el borde de la caja de simulación.}
	\label{fig:CST}
\end{figure}

Dado que el parámetro $h$ define el tamaño del dominio de simulación, el número total de celdas empleado en la discretización depende directamente de su valor, tanto a través del mallado global, como del refinamiento local. Por esta razón, se analiza en primer lugar su influencia sobre $Q_{\text{ext}}$ y sobre el error%
\footnote{En este trabajo se consideran dos tipos de errores: error absoluto a una longitud de onda dada y el error absoluto promedio. El error absoluto a una longitud de onda dada se define como
	$$\text{Error}(\lambda)=|Q_{\text{CST}}(\lambda)-Q_{\text{Mie}}(\lambda)|,$$
	mientras que el error absoluto promedio se considera como 
	$$\text{Error promedio}=\frac{1}{N}\sum_i|Q_{\text{CST}}(\lambda_i)-Q_{\text{Mie}}(\lambda_i)|,$$
	donde $N$ es el número total de datos promediados.} % 
para identificar el valor de $h$ que no incurra en un tiempo de cómputo excesivo. Para este análisis se mantuvieron fijos los parámetros predeterminados de CST: un mallado global de 15 celdas por longitud de onda (cel./$\lambda$), sin refinamiento local y un nivel de reflexión estimado de $10^{-5}$, variando únicamente el valor de $h$. En la Fig.~\ref{fig:ajusteMatriz} se presenta $Q_{\text{ext}}$ en función de la longitud de onda incidente $\lambda$ para distintos valores de $h$: 500, 1,000, 1,230, 1,470, 1,720 y 2,300 nm. Los resultados se identifican mediante puntos de color negro, morado, vino, naranja oscuro, naranja claro y amarillo, respectivamente, mientras que la línea continua que los une se incluye únicamente como guía visual. Asimismo, se muestra la solución analítica obtenida mediante teoría de Mie como una línea continua verde. En la solución analítica se observan tres mínimos espectrales asociados a los centros de las principales bandas de absorción descritas por la función dieléctrica: la banda de Soret, centrada alrededor de 428 nm, y las bandas $\beta$ y $\alpha$, centradas aproximadamente en 541 nm y 575 nm, respectivamente. El espectro completo de 400 a 600 nm se presenta en la Fig.~\ref{subfig:ajusteMatrizGlobal}, mientras que en las Figs.~\ref{subfig:ajusteMatrizSoret}, \ref{subfig:ajusteMatrizBeta} y \ref{subfig:ajusteMatrizAlpha} se muestran ampliaciones en las regiones correspondientes a la banda de Soret y a las bandas $\beta$ y $\alpha$. %
% --- JAUA ---
% Propuesta para alas figuras: Sombrea las bandas de Soret, alha y beta en la principal, y en las ampliaciones pones el fondo de esos colores. Para evitar pensarlo, ponle distintos tonos de gris.
%

%
\begin{figure}[h!]
	\centering
	\sidesubfloat[First image]{\hspace{-0.1cm}{
			\includegraphics[width=0.45\textwidth]{3-FIT/ajusteMatriz.pdf}\label{subfig:ajusteMatrizGlobal}}}\hspace{0.1cm}
	\sidesubfloat[Second image]{\hspace{-0.1cm}{	\includegraphics[width=0.45\textwidth]{3-FIT/ajusteMatrizSoret.pdf}\label{subfig:ajusteMatrizSoret}}}
	
	\vspace{0.3cm}
	
	\sidesubfloat[Third image]{\hspace{-0.1cm}{	\includegraphics[width=0.45\textwidth]{3-FIT/ajusteMatrizBeta.pdf}\label{subfig:ajusteMatrizBeta}}}	\hspace{0.1cm}
	\sidesubfloat[Fourth image]{\hspace{-0.1cm}{	\includegraphics[width=0.45\textwidth]{3-FIT/ajusteMatrizAlpha.pdf}\label{subfig:ajusteMatrizAlpha}}}
	
	
	\caption{Eficiencias de extinción $Q_{\text{ext}}$ para una esfera de radio de 2,680 nm con función dieléctrica correspondiente a un eritrocito con CHCM de 30.6 g/dL iluminada con una onda plana en función de la longitud de onda $\lambda$. Se comparan los datos obtenidos mediante CST con  valores de $h$ de 500, 1,000, 1,230, 1,570, 1,720 y 2,300 nm que se identifican mediante puntos de color negro, morado, vino, naranja oscuro, marajna claro y amarillo, respectivamente. La línea continua que une los puntos es una guía visual. Se incluye también la solución analitíca de teoría de Mie (línea verde continutua)  \textbf{a)} Espectro de 400 a 600 nm. \textbf{b)} Región asociada a la banda de Soret \textbf{c)} Región asociada a la banda $\beta$. \textbf{d)} Región asociada a la banda $\alpha$.  }
	\label{fig:ajusteMatriz}
\end{figure}
%

A partir de estos resultados, se observa un corrimiento al rojo en el mínimo global entre la solución numérica y la analítica, así como oscilaciones en todo el espectro. Dado que estas oscilaciones no se observan en la solución analítica, se espera que no correspondan a un comportamiento físico del sistema. Además, se observa que las oscilaciones se encuentran más próximas entre sí a menores longitudes de onda, lo que indica una mayor frecuencia de oscilación. En cuanto a la influencia de la distancia $h$, se observa que en la región de la banda de Soret el error entre la solución de Mie con la solución analítica disminuye al incrementar $h$ alcanzando valores de hasta 0.022. En contraste, en las bandas $\beta$ y $\alpha$ el error aumenta para valores menores de $h$ alcanzando valores de 0.64. Considerando el comportamiento global en todo el rango espectral, así como el costo computacional asociado (82 minutos), se seleccionó $h=1,000$~nm como un valor conveniente para las simulaciones posteriores, dado que el ajuste de otros parámetros podría incrementar significativamente el tiempo de cómputo.

El análisis de la influencia del mallado global se realizó mediante la variación del número de celdas por longitud de onda. Para este estudio no se consideró refinamiento local y se fijó el nivel de reflexión en $10^{-5}$. En la Fig.~\ref{fig:ajusteMesh} se muestra la $Q_{\text{ext}}$ en función de la longitud de onda de la luz incidente, comparando los resultados obtenidos mediante CST con la solución analítica de Mie. El espectro completo se presenta en la Fig. \ref{subfig:ajusteMeshGlobal}, mientras que en las Figs.~\ref{subfig:ajusteMeshBeta} y \ref{subfig:ajusteMeshAlpha} se muestran ampliaciones en las regiones correspondientes a la banda de Soret y a las bandas $\beta$ y $\alpha$. El número de celdas por unidad de longitud determina directamente el tamaño mínimo de celda en la discretización. En este caso, se consideraron valores de 7, 10, 12, 15 y 17~cel./$\lambda$, correspondientes a tamaños de celda de 40.0, 28.2, 23.5, 18.7 y 16.5 nm, respectivamente. Estos puntos se representan mediante un gradiente de color rosa, que va de tonos más oscuros para los valores menores (7 cel./$\lambda$) a tonos más claros para los valores mayores (17 cel./$\lambda$).  También se muestra la solución analítica de Mie como referencia como una línea continua verde. Se observa que en las regiones de la banda de Soret, $\beta$ y $\alpha$ conforme aumenta el tamaño mínimo de celda, las desviaciones respecto a la solución analítica aumentan. Estas discrepancias se manifiestan como oscilaciones espurias y un corrimiento al rojo en la posición del mínimo global, lo que indica que la discretización es insuficiente para describir adecuadamente la respuesta electromagnética de la partícula. A medida que se refina el mallado la discrepancia entre los resultados de CST y la solución de Mie disminuye. En particular, para el mallado más fino considerado se obtiene el menor error absoluto promedio, igual a 0.49. No obstante, incluso en este caso persiste un corrimiento al rojo del mínimo global de 37 nm. Considerando el compromiso entre precisión y costo computacional, se seleccionó un mallado de 15 cel./$\lambda$, ya que la mejora obtenida al incrementar a 17 celdas es poco significativa (del orden de $10^{-3}$ en error absoluto promedio), mientras que el tiempo de simulación aumenta 50 minutos adicionales.
%
\begin{figure}[tb]
	\centering	\sidesubfloat[First image]{\hspace{-0.1cm}{
			\includegraphics[width=0.45\textwidth]{3-FIT/ajusteMesh.pdf}\label{subfig:ajusteMeshGlobal}}}\hspace{0.1cm}
	\sidesubfloat[Second image]{\hspace{-0.1cm}{	\includegraphics[width=0.45\textwidth]{3-FIT/ajusteMeshSoret.pdf}\label{subfig:ajusteMeshSoret}}}
	
	\vspace{0.3cm}
	
	\sidesubfloat[Third image]{\hspace{-0.1cm}{	\includegraphics[width=0.45\textwidth]{3-FIT/ajusteMeshBeta.pdf}\label{subfig:ajusteMeshBeta}}}	\hspace{0.1cm}
	\sidesubfloat[Fourth image]{\hspace{-0.1cm}{	\includegraphics[width=0.45\textwidth]{3-FIT/ajusteMeshAlpha.pdf}\label{subfig:ajusteMeshAlpha}}}
	
	
	\caption{Eficiencias de extinción $Q_{\text{ext}}$ para una esfera de radio de 2,680 nm con función dieléctrica correspondiente a un eritrocito con CHCM de 30.6 g/dL iluminada con una onda plana en función de la longitud de onda $\lambda$. Se comparan los datos obtenidos mediante CST con tamaños mínimos de celda iguales a 40, 28.2, 23.5, 18.7 y 16.5 nm que se identifican mediante un gradiente de color rosa que va desde tonos más oscuros a tonos más claros. La línea continua que une los puntos es una guía visual. Se incluye también la solución analitíca de teoría de Mie (línea verde continua)  \textbf{a)} Espectro de 400 a 600 nm. \textbf{b)} Región asociada a la banda de Soret. \textbf{c)} Región asociada a la banda $\beta$. \textbf{d)} Región asociada a la banda $\alpha$. }
	\label{fig:ajusteMesh}
\end{figure}
%

Para evaluar el efecto del mallado local se analizó el refinamiento en la superficie%
\footnote{Dado que en CST se implementa la técnica de %
	\textit{Perfect Boundary Approximation} (PBA), la cual permite representar superficies curvas sin necesidad de refinamientos adicionales en la discretización geométrica \cite{dassaultsystemesUnderstandingTimeDomain2010}, no se observaron variaciones ni en $Q_{\text{ext}}$ ni en el tiempo de simulación al realizar el refinamiento del mallado en la superficie, por lo que no se incluyen los resultados correspondientes.
% --- JAUA ---
% Por lo que me contaste,  más bien incluso haciendo estos cambios, CST no mofigico lo parámetros del cálculo. Si es esto cierto, hay que redactarlo de nuevo porque así suena a que hiciste cáculos y dio lo mismo.
%Sí hice los cálculos y dio lo mismo y luego me di cuenta de que cambiar eso no cambiaba en nada el numero de celdas
}
y en el volumen. Para el volumen de la esfera se fijaron los parámetros previamente optimizados: $h=1,000$ nm, un mallado global de 15 cel./$\lambda$ y un nivel de reflexión de $10^{-5}$, variando únicamente el tamaño mínimo absoluto de celda dentro de la partícula. En la Fig.~\ref{fig:ajusteLocalMesh} se muestran los resultados obtenidos de $Q_{\text{ext}}$ para tamaños de celda de 20, 16, 15, 12 y 10 nm, los cuales corresponden a tamaños efectivos mínimos dentro de la esfera de 16.17, 15.95, 14.83, 11.96 y 9.93 nm, respectivamente. Los resultados se representan mediante un gradiente de color morado, que varía desde tonos más oscuros para 16.17 nm hasta tonos más claros para 9.93 nm.  Adicionalmente, se incluye la solución analítica de la teoría de Mie como una línea continua de color verde. El espectro completo se presenta en la Fig. \ref{subfig:ajusteMeshGlobal}, mientras que en las Figs. \ref{subfig:ajusteLocalMeshSoret}, \ref{subfig:ajusteLocalMeshBeta} y \ref{subfig:ajusteLocalMeshAlpha} se muestran ampliaciones en las regiones correspondientes a la banda de Soret y a las bandas $\beta$ y $\alpha$. Se observa que, al disminuir el tamaño de celda, los resultados numéricos convergen progresivamente hacia la solución analítica. Asimismo, la inclusión del mallado local en el volumen de la partícula elimina las oscilaciones espurias observadas previamente, confirmando su origen numérico asociado a una discretización insuficiente de la geometría. Con base en estos resultados, se seleccionó un tamaño mínimo de celda de 10 nm, al representar un equilibrio entre precisión y costo computacional. Cabe señalar que refinamientos adicionales no fueron considerados debido a limitaciones de memoria RAM del equipo de cómputo empleado.
%
\begin{figure}[tb]
	\centering
	\sidesubfloat[First image]{\hspace{-0.1cm}{
			\includegraphics[width=0.45\textwidth]{3-FIT/ajusteMeshLocal.pdf}\label{subfig:ajusteLocalMeshGlobal}}}\hspace{0.1cm}
	\sidesubfloat[Second image]{\hspace{-0.1cm}{	\includegraphics[width=0.45\textwidth]{3-FIT/ajusteMeshLocalSoret.pdf}\label{subfig:ajusteLocalMeshSoret}}}
	
	\vspace{0.3cm}
	
	\sidesubfloat[Third image]{\hspace{-0.1cm}{	\includegraphics[width=0.45\textwidth]{3-FIT/ajusteMeshLocalBeta.pdf}\label{subfig:ajusteLocalMeshBeta}}}	\hspace{0.1cm}
	\sidesubfloat[Fourth image]{\hspace{-0.1cm}{	\includegraphics[width=0.45\textwidth]{3-FIT/ajusteMeshLocalAlpha.pdf}\label{subfig:ajusteLocalMeshAlpha}}}
	
	
	\caption{Eficiencias de extinción $Q_{\text{ext}}$ para una esfera de radio de 2,680 nm con función dieléctrica correspondiente a un eritrocito con CHCM de 30.6 g/dL iluminada con una onda plana en función de la longitud de onda $\lambda$. Se comparan los datos obtenidos mediante CST con tamaños mínimos de celda iguales a 16.17, 15.95, 14.83, 11.96 y 9.93 nm que se identifican mediante un gradiente de color morado que va desde tonos más oscuros a tonos más claros. La línea continua que une los puntos es una guía visual. Se incluye también la solución analitíca de teoría de Mie (línea verde continua)  \textbf{a)} Espectro de 400 a 600 nm. \textbf{b)} Región asociada a la banda de Soret. \textbf{c)} Región asociada a la banda $\beta$. \textbf{d)} Región asociada a la banda $\alpha$. }
	\label{fig:ajusteLocalMesh}
\end{figure}
%

El último parámetro analizado fue el nivel de reflexión estimado del PML considerando los valores $10^{-4}$, $10^{-6}$, $10^{-8}$ y $10^{-10}$ y manteniendo fijos los valores de los otros parámetros previamente seleccionados. En CST estos niveles de reflexión corresponden, respectivamente, a PML constituidos por 4, 6, 8 y 10 capas, cuyos espesores resultantes fueron de 68.26, 85.32, 102.39 y 136.51 nm. Como resultado, no se observaron diferencias significativas en los valores de $Q_{\text{ext}}$ obtenidos; en particular, para la banda de Soret se obtuvo un error absoluto de 0.43150 al emplear un nivel de reflexión de $10^{-4}$, mientras que para $10^{-6}$, $10^{-8}$ y $10^{-10}$ se obtuvieron errores de 0.43134, 0.43139 y 0.43134, respectivamente es decir, cambios del orden de $10^{-4}$ y $10^{-5}$. Dado que la reducción del nivel de reflexión no produce mejoras significativas en la convergencia de los resultados, pero sí incrementa el tiempo de simulación, se seleccionó finalmente un valor de $10^{-4}$ para las simulaciones posteriores.

A continuación, se sintetizan los resultados del análisis de convergencia, en términos del tiempo de simulación, con el objetivo de evaluar el equilibrio alcanzado entre precisión numérica y costo computacional. En la Fig.~\ref{fig:tiempos} se muestran los tiempos de simulación (en minutos) para cada una de las variaciones de parámetros consideradas, correspondientes a un barrido espectral de 400 a 600 nm. Este barrido espectral comprende un total de 136 longitudes de onda: de 400 a 450 nm con incrementos de 1 nm, de 455 a 530 nm con incrementos de 5 nm y de 530 a 600 nm nuevamente con incrementos de 1 nm. Se presentan los resultados conforme se realizó la variación de parámetros, es decir, de izquierda a derecha se muestran los resultados correspondientes a la variación de la distancia $h$, el mallado global, el mallado local en volumen y el nivel de reflexión asociado a la capa absorbente (PML). Cada punto está identificado mediante el mismo código de color empleado en las figuras anteriores, mientras que en color verde se indica el valor seleccionado como óptimo en cada caso. Con base en el análisis conjunto de la convergencia y el costo computacional, los parámetros finales seleccionados son: $h=1,000$ nm, mallado global con 15 cel./$\lambda$ de onda, mallado local con un tamaño mínimo de celda de 10 nm y nivel de reflexión estimado de $10^{-4}$. Esta configuración produce un error absoluto promedio de aproximadamente 0.05 respecto a la solución analítica de Mie, con un tiempo total de simulación de alrededor de 770 minutos.
%
\begin{figure}[tb]
	\centering
		\includegraphics[width=1\textwidth]{3-FIT/tiempos_convergencia.pdf}
	
	\caption{Tiempo de ejecución de las simulaciones realizadas en CST. En orden de izquierda a derecha se observa el tiempo de simulación en función de $h$, del tamaño mínimo de celda y del nivel de reflexión para cada uno de los parámetros optimizados: distancia borde-esfera, el mallado global, el mallado local y el PML.
	  }
	\label{fig:tiempos}
\end{figure}
%


Con los parámetros previamente seleccionados, se realizó una verificación adicional de los criterios de convergencia empleados. Para ello, se fijó un tamaño mínimo de celda de 10 nm y se repitió el análisis de la distancia borde-esfera con valores de $h=950$, 1,000 y 1,050 nm y del mallado global con valores de 12, 15 y 17 cel./$\lambda$. Los resultados obtenidos se compararon con la solución analítica de Mie en tres longitudes de onda representativas: 428, 529 y 563 nm. Estas longitudes corresponden, respectivamente, al mínimo asociado a la banda de Soret y a los máximos asociados a las bandas $\beta$ y $\alpha$. En la Tabla \ref{table:convergencia_h} se muestran los valores de $Q_{\text{ext}}$ obtenidos así como el error relativo con respecto a la solución analítica de Mie con un mallado global de 15~cel./$\lambda$ y valores de $h=950, 1,000$ y 1500 nm; mientras que en la Tabla \ref{table:convergencia_malla} se fija el valor de $h=1,000$ nm y se modifica el mallado global con 12, 15 y 17 cel./$\lambda$.
%
\begin{center} 
		
		% -------- BLOQUE 1 (Izquierda) --------
		\begin{minipage}[t]{0.48\textwidth}
			\centering
				
			\begin{tabular}{ccc}
				\hline
				\multicolumn{3}{c}{15 cel./$\lambda$}\\ 
				\hline
				$\lambda$ [nm] & $Q_{\text{ext}}$ & Error relativo\\
				\hline
				\multicolumn{3}{c}{$h=950$ nm} \\
				\hline
				428 & 2.07084 & 4.56847 \\
				529 & 3.07943 & 0.994313 \\
				563 & 3.17225 & 0.670544 \\
				\hline
				\multicolumn{3}{c}{$h=1,000$ nm} \\
				\hline
				428 & 2.06953 & 4.57459 \\
				529 & 3.07902 & 1.00780 \\
				563 & 3.17233 & 0.668147 \\
				\hline
				\multicolumn{3}{c}{$h=1,050$ nm} \\
				\hline
				428 & 2.06824 & 4.64005 \\
				529 & 3.07854 & 1.02350 \\
				563 & 3.17232 & 0.668147 \\
				\hline
			\end{tabular}
			\captionof{table}{Valores de $Q_{\text{ext}}$ y errores relativos respecto a la solución analítica de Mie para un mallado de 15 cel./$\lambda$ con valores de $h=950, 1,000$ y 1,500 nm.}
			\label{table:convergencia_h}
			\vspace{0.2cm} 
		\end{minipage}
		\hfill 
		% -------- BLOQUE 2 (Derecha) --------
		\begin{minipage}[t]{0.48\textwidth}
			\centering
					
			\begin{tabular}{ccc}
				\hline
				\multicolumn{3}{c}{$h=1,000$ nm}\\ 
				\hline
				$\lambda$ [nm] & $Q_{\text{ext}}$ & Error relativo\\
				\hline
				\multicolumn{3}{c}{12 cel./$\lambda$} \\
				\hline
				428 & 2.04120 & 6.02605 \\
				529 & 3.06460 & 1.48286 \\
				563 & 3.16445 & 0.918616 \\
				\hline
				\multicolumn{3}{c}{15 cel./$\lambda$} \\
				\hline
				428 & 2.06953 & 4.57459 \\
				529 & 3.07902 & 1.00780 \\
				563 & 3.17233 & 0.668140 \\
				\hline
				\multicolumn{3}{c}{17 cel./$\lambda$} \\
				\hline
				428 & 2.07888 & 4.10444 \\
				529 & 3.08384 & 0.849720 \\
				563 & 3.17484 & 0.588530 \\
				\hline
			\end{tabular}
			\captionof{table}{Valores de $Q_{\text{ext}}$ y errores relativos respecto a la solución analítica de Mie para $h=~1,000$~nm, variando el mallado global con 12, 15 y 17 cel./$\lambda$.}
			\label{table:convergencia_malla}
			\vspace{0.2cm}
		\end{minipage}
		
	\end{center}
	\vspace{0.4cm} % Espacio después de las tablas



Se observa que en la longitud de onda de 428 nm las diferencias en $Q_{\text{ext}}$ son del orden de $10^{-2}$ o menores, mientras que para 529 y 563 nm las diferencias son del orden de $10^{-3}$ o menores. Los errores relativos muestran que las tendencias observadas en los análisis iniciales se conservan tras la incorporación del mallado local. Por otra parte, el mallado de 17 cel./$\lambda$ presenta la mejor concordancia con la solución analítica de Mie en las tres longitudes de onda consideradas, alcanzando errores relativos inferiores al 4.2\%. Con base en estos resultados, se adoptó finalmente la combinación $h=1,000$ nm y un mallado global de 17 cel./$\lambda$ para las simulaciones posteriores.




% !TeX root = ../tesis.tex
\chapter{Propiedades ópticas}
\label{chapter:Resultados}


Una vez establecidos los modelos geométricos de los eritrocitos sanos y enfermos, sus funciones dieléctricas y ajustado los parámetros numéricos para la convergencia del método FIT, se calculan sus propiedades ópticas. Este análisis permite evaluar cómo las variaciones simultáneas en la morfología celular y en la concentración de hemoglobina modifican la respuesta óptica de los eritrocitos e identificar regiones espectrales y angulares en las que dichas diferencias se intensifican. En este capítulo se presentan los resultados obtenidos para el eritrocito sano y para dos alteraciones morfológicas representativas: el esferocito y el microcito. Las simulaciones se realizaron en el intervalo espectral de 400 a 600 nm. Para describir las propiedades ópticas de los eritrocitos se emplean las funciones dieléctricas complejas obtenidas en la Sección~\ref{section:optical_properties}, las cuales dependen tanto de la longitud de onda como de la concentración de hemoglobina corpuscular media (CHCM), mientras que el plasma se considera con una permitividad constante en todo el intervalo espectral. En primer lugar, se evalúa la capacidad del modelo numérico para reproducir resultados reportados en la literatura mediante la simulación de un eritrocito descrito con el modelo de Fung \cite{fungBiomechanicsMechanicalProperties1990} a una longitud de onda de $632.8~\text{nm}$. Una vez validado el modelo, se presentan los resultados correspondientes al eritrocito sano modelado mediante un óvalo de Cassini. Posteriormente, se comparan las propiedades ópticas del eritrocito sano, el esferocito y el microcito en el rango de interés por medio sus eficiencias de extinción, así como de los patrones de esparcimiento en tres longitudes de onda representativas de las bandas de absorción de la hemoglobina. Posteriormente, se estudia la distribución angular del esparcimiento mediante el factor de anisotropía, las distribuciones angulares y las contribuciones integradas en las regiones frontal, lateral y de retroesparcimiento. A partir de este análisis se identifican regiones angulares en las que el contraste entre eritrocitos sanos y enfermos aumenta.

\section{Contraste con la literatura}
\label{section:Res_literatura}

Con el fin de evaluar la capacidad del modelo para reproducir resultados reportados en la literatura, se compararon las simulaciones realizadas en este trabajo con los resultados de Tsinopoulos et al. \cite{tsinopoulosLightScatteringAggregated2002}, quienes emplearon el modelo de eritrocito propuesto por Fung \cite{fungBiomechanicsMechanicalProperties1990}, descrito por la ecuación en coordenadas cilíndricas
%
\begin{equation}
	z(\rho)=\left[1-\left(\dfrac{\rho}{a}\right)^2\right]^{1/2}\left[0.72+4.152\left(\dfrac{\rho}{a}\right)^2-3.426\left(\dfrac{\rho}{a}\right)^4\right],
	\label{eq:Tsang_eri}
\end{equation}
%
donde $\rho^2=x^2+y^2$ y $a$ es el radio del eritrocito, para el cual se empleó el valor de $3.825\,\mu$m. Esta geometría corresponde a un eritrocito con un grosor mínimo de $1.44\,\mu$m y un grosor máximo de $2.84\,\mu$m, un volumen de $97.91~\mu\text{m}^3$ y un área de $129.95~\mu\text{m}^2$. Cabe señalar que esta parametrización fue desarrollada para describir eritrocitos sanos, por lo que su aplicación se restringe a este tipo de células. En sus simulaciones, Tsinopoulos et al. emplearon el \textit{Boundary Element Method} (BEM)
y consideraron un eritrocito orientado con su eje de simetría a lo largo de la dirección $z$, iluminado por una onda plana propagándose en la misma dirección y polarizada en $x$, con una longitud de onda de $632.8~\text{nm}$. En esta longitud de onda consideraron un índice de refracción de 1.402 para el eritrocito y de 1.335 para el plasma, lo que corresponde a un índice de refracción relativo de aproximadamente 1.05. Estas mismas condiciones fueron reproducidas en CST empleando los parámetros de convergencia determinados en la Sección \ref{section:convergence}. En la Fig. \ref{fig:Tsinopoulos} se muestran los resultados de la sección transversal diferencial de esparcimiento $\dd{C_{\text{sca}}/\dd{\Omega}}$ [Ec. \eqref{C_sca_general}] en función del ángulo de esparcimiento $\theta$  obtenidos mediante CST con puntos azules y los obtenidos mediante BEM extraídos de la Ref.~\cite{tsinopoulosLightScatteringAggregated2002}. En la Fig. \ref{subfig:Tsi_par} se muestra el corte de $\varphi=0^{\circ}$ (plano paralelo a la dirección de polarización de la onda incidente), mientras que en la Fig. \ref{subfig:Tsi_perp} se muestra el corte de $\varphi=90^{\circ}$ (plano perpendicular a la dirección de polarización de la onda incidente). %
%
\begin{figure}[h!]
	\centering
	\sidesubfloat[First image]{\hspace{-0.1cm}{
			\includegraphics[width=0.45\textwidth]{4-Resultados/comparacionpaper0.pdf}\label{subfig:Tsi_par}}}\hspace{0.1cm}
	\sidesubfloat[Second image]{\hspace{-0.1cm}{	\includegraphics[width=0.45\textwidth]{4-Resultados/comparacionpaper90.pdf}\label{subfig:Tsi_perp}}}
	
	
	\caption{Sección transversal diferencial de esparcimiento $\dd{C_{\text{sca}}/\dd{\Omega}}$ en función del ángulo de esparcimiento $\theta$ para un eritrocito con geometría descrita por Fung \cite{fungBiomechanicsMechanicalProperties1990} y con radio igual a 3.825 $\mu$m. El eritrocito fue iluminado por una onda plana propagándose en la dirección $z$ y polarizada en $x$, con una longitud de onda de $632.8~\text{nm}$. Se consideró un índice de refracción relativo de 1.05. Los puntos azules muestran los resultados obtenidos mediante el método FIT y los naranjas mediante BEM en la Ref.  \cite{tsinopoulosLightScatteringAggregated2002}}
	\label{fig:Tsinopoulos}
\end{figure}
%

Se observa que ambos métodos reproducen la misma estructura angular general del patrón de esparcimiento. En particular, en el plano paralelo  ($\varphi=0^{\circ}$), la diferencia entre en la posición angular de los máximos locales obtenida en la Ref.~\cite{tsinopoulosLightScatteringAggregated2002} y la obtenida en CST es de $0.63^\circ$. De manera similar, en el plano perpendicular ($\varphi=90^{\circ}$), la diferencia promedio en la localización de los máximos locales es de  $0.46^\circ$. Por otro lado, se encuentran diferencias en la amplitud de $\dd{C_{\text{sca}}/\dd{\Omega}}$ con un error relativo promedio de 0.85 y 0.54 para $\varphi=0^{\circ}$ y $\varphi=90^{\circ}$, respectivamente.  Este tipo de diferencias tanto en la amplitud como en la posición de máximos obtenidos mediante diferentes métodos numéricos han sido reportadas previamente en las Refs.~\cite{kupresakBenchmarkingSoftwareTools2017, amirjaniFiniteBoundaryElement}, donde se señala que pueden persistir aún cuando cada método haya sido convergido de manera independiente debido a las diferencias en su formulación numérica y en la discretización espacial. En conjunto, la concordancia observada en la estructura angular y en la localización de los máximos indica que el modelo implementado en CST reproduce el patrón angular de esparcimiento reportado por Tsinopoulos et al., aunque persisten diferencias cuantitativas en su amplitud \cite{kupresakBenchmarkingSoftwareTools2017}. Estos resultados proporcionan una validación adicional de la implementación numérica empleada en este trabajo.
      
% !TeX root = ../tesis.tex

\section{Eritrocito sano}
\label{section:Res_Eritrocito_Sano}

Para identificar las propiedades ópticas de los eritrocitos en las bandas de Soret, $\beta$ y $\alpha$, se consideró el rango espectral de 400  a 600 nm.  A lo largo de esta sección, por comodidad se define a una \textit{dirección} como un ángulo de esparcimiento específico ($\theta=0^\circ$ para el esparcimiento frontal, $\theta=90^\circ$ para el lateral y $\theta=180^\circ$ para el retroesparcimiento), mientras que una \textit{región} corresponde a un intervalo angular ($0^\circ\leq\theta<60^\circ$ para el esparcimiento frontal, $60^\circ\leq\theta<120^\circ$ para el lateral y $120^\circ\leq\theta\leq180^\circ$ para el retroesparcimiento).%
% --- JAUA ---
% te propongo mejor ennumerar los nombres y entrea paréntesis los valores. Por ejemplo:
% se definen la dirección frontal como ek pangulo de esparcimeinto de \theta =0, la lateral de 90° y la dirección de tresotespacimieto para 180°. Por otro lado, se defienen tambiénn las regiones de esparcimiento frontal (pones los valores), el lateral (cvalores) y.... 
% Creo que eso lo haría más claro de leer y ponemos peso en todos los nombres que utilizas.
%
  Para modelar un eritrocito sano, se empleó un óvalo de Cassini %
  % --- JAUA ---
  % con los parámetreos mostrados en la Fig XX con....
  %
   como el de la Fig. \ref{subfig:Shape_health} con parámetros de $a=2.6\,\mu\text{m}$, $b=2.7\,\mu\text{m}$ y $c=0.812\,\mu\text{m}$, los cuales resproducen el radio equivalente de 3.75 $\mu$m y el volumen de 80.63~$\mu\text{m}^3$ reportados por Ergül et al. en la Ref.~\cite{ergulComputationalStudyScattering2010}. De acuerdo con los resultados obtenidos en la Sección~\ref{section:optical_properties}, para el plasma se consideró una función dieléctrica constante de $\varepsilon_{\text{Plasma}}=1.839$, %
   % --- JAUA ---
   % en secciones anteriores habiamos modificado a 1.84; recomiendo dejarlo de esa forma y omitir que cambia para cuatro ciffras decimales.
   %
    ya que su parte imaginaria resulta despreciable y la parte real presenta variaciones únicamente en la cuarta cifra decimal dentro del rango espectral de interés. Por su parte, para el eritrocito se empleó la función dieléctrica de una CHCM de 30.6~g/dL %
	% --- JAUA ---
	% Creo que hay que enfatizar que un eritrosito sano corresponde a esa concentración y que no es una elección arbitraria; te recomiendo hacer mención de las figuras correspodientes en secciones anterior3es, o reecoradar que esto se discutió ya antes.
	%
	. El eritrocito se orientó con su eje de simetría a lo largo de la dirección $z$ y se iluminó con una onda plana que se propaga en la misma dirección, polarizada en $x$.

En la Fig. \ref{subfig:Qext_health} se muestra la eficiencia de extinción $Q_{\text{ext}}$ en función de la longitud de onda $\lambda$ de la onda incidente; se resaltan en gris las bandas que rodean los mínimos asociados a la banda de Soret y a las bandas $\beta$ y $\alpha$ localizados en 431 nm, 541 nm y 573 nm, respectivamente, los cuales se consideran longitudes de onda representativas para el análisis de los patrones de esparcimiento. Se observa que la respuesta espectral del eritrocito sigue la dependencia espectral de su función dieléctrica, reflejando la influencia de las transiciones ópticas de la hemoglobina como mínimos en $Q_{\text{ext}}$.
%
\begin{figure}[h!]
	\centering
	\sidesubfloat[First image]{\hspace{-0.3cm}{
			\includegraphics[width=0.45\textwidth]{4-Resultados/Eritrocito_sano_corte.pdf}\label{subfig:Shape_health}}}\hspace{0.1cm}
	\sidesubfloat[Second image]{\hspace{-0.5cm}{	\includegraphics[width=0.45\textwidth]{4-Resultados/Qext_EriSano.pdf}\label{subfig:Qext_health}}}
	\vspace{-0.1cm}
	\caption{\textbf{a)} Modelo de eritrocito sano basado en óvalos de Cassini con parámetros de $a=2.6\,\mu\text{m}$, $b=2.7\,\mu\text{m}$ y 
		$c=0.812\,\mu\text{m}$, que corresponden a un radio de $3.75\,\mu$m. \textbf{b)} Eficiencia de extinción $Q_{\text{ext}}$ en función de la longitud de onda $\lambda$ de la onda incidente para un eritrocito orientado con su eje de simetría en el eje $z$ iluminado con una onda en dirección del eje $z$ y polarizada en $x$. La línea continua que une los puntos es una guía visual. Se resaltan en gris las bandas que 
		rodean los mínimos asociados a la banda de Soret y a las bandas Q.}
	\label{fig:Shape_Qext_health}
\end{figure}
%

En esta sección se discute únicamente el patrón de esparcimiento de $\lambda=431$~nm asociado a la banda de Soret, mientras que los de $\lambda=541$ y $573$ nm se incluyen en el Apéndice \ref{section:apendix3}. En la Fig.~\ref{subfig:FF_431_health} se muestra la sección transversal diferencial de esparcimiento $\dd{C_{\text{sca}}}/\dd{\Omega}$ para $\lambda=431$~nm. La magnitud de $\dd{C_{\text{sca}}}/\dd{\Omega}$ se encuentra normalizada respecto a su valor máximo, que corresponde al valor más saturado. Se observa que $\dd{C_{\text{sca}}}/\dd{\Omega}$ es mayor en la dirección frontal de propagación, lo cual se refleja en un factor de anisotropía\footnote{%
	El factor de anisotropía $g$ describe la direccionalidad de la luz esparcida por una partícula y corresponde al promedio ponderado de la intensidad de esparcimiento normalizada con el coseno del ángulo de esparcimiento \cite{videenAsymmetryParameterAggregate1998}. Se define como \cite{bohrenAbsorptionScatteringLight2008}
	$$g=\dfrac{1}{C_{\text{sca}}}\int_{4 \pi}\dfrac{\dd{C_{\text{sca}}}}{\dd{\Omega}}\cos\theta\dd{\Omega},$$
	por lo que $g$ varía entre $-1$ y 1. En particular, para una partícula que esparce la luz de forma isótropa o si su esparcimiento es simétrico respecto a un ángulo de esparcimiento $\theta=90^{\circ}$, $g$ se anula. Si la partícula esparce más luz hacia la dirección frontal ($\theta=0^{\circ}$), $g>0$; si el esparcimiento se dirige más hacia la dirección posterior ($\theta=180^{\circ}$), $g<0$ \cite{bohrenAbsorptionScatteringLight2008}.} %
$g$ igual a 0.97. Se resaltan además, dos planos de observación: $\varphi=0^{\circ}$ (azul claro), que es paralelo a la polarización de la onda indicente y $\varphi=90^{\circ}$ (azul oscuro), perpendicular a la polarización de la onda incidente. Con el fin de estudiar el efecto de la polarización de la onda incidente en el esparcimiento, en la Fig. \ref{subfig:FF_polar_431_health} se muestra $\dd{C_{\text{sca}}}/\dd{\Omega}$ en decibeles en función del ángulo de esparcimiento $\theta$ para dos cortes correspondientes a los planos $\varphi=0^{\circ}$ (línea continua azul claro) y $\varphi=90^{\circ}$ (línea continua azul oscuro). En la Fig. \ref{subfig:FF_cartesian_431_health} se muestra $\dd{C_{\text{sca}}}/\dd{\Omega}$ en función del ángulo de esparcimiento para ambos cortes en esala logarítimica. Debido a que las dimensiones del eritrocito son mayores que la longitud de onda incidente, su respuesta se encuentra fuera del régimen dipolar. En consecuencia, el patrón de esparcimiento presenta una estructura angular compleja, asociada con los retardos de fase y la interferencia entre las contribuciones al campo esparcido provenientes de distintas regiones de la célula.
	%
	\begin{figure}[t]
		\centering
		\sidesubfloat[First image]{\hspace{-0.5cm}{	\includegraphics[width=7.3cm]{4-Resultados/FFdistribution431_EriSano.pdf}\label{subfig:FF_431_health}}}\hspace{0.1cm}
		\sidesubfloat[Second image]{\hspace{-0.5cm}{
				\includegraphics[width=8cm]{4-Resultados/polarplot431FF_EriSano.pdf}\label{subfig:FF_polar_431_health}}}\vspace{-0.1cm}
		\
		\vspace{0.1cm}
		
		\sidesubfloat[Third image]{\hspace{0.3cm}{
				\includegraphics[width=14cm]{4-Resultados/linearplot431FF_EriSano.pdf}\label{subfig:FF_cartesian_431_health}}}\vspace{-0.1cm}
		\caption{\textbf{a)} Sección transversal diferencial de esparcimiento $\dd{C_{\text{sca}}}/\dd{\Omega}$ normalizada respecto a su valor máximo para una longitud de onda $\lambda=431$ nm en función del ángulo de esparcimiento $\theta$ y del ángulo $\varphi$ para un eritrocito sano. Se identifican dos planos $\varphi=0^{\circ}$ en azul claro y $\varphi=90^{\circ}$ en azul oscuro. \textbf{b)} Sección transversal diferencial de esparcimiento $\dd{C_{\text{sca}}}/\dd{\Omega}$ en decibeles en función de $\theta$ para dos cortes: $\varphi=0^{\circ}$ (puntos azules claro) y $\varphi=90^{\circ}$ (puntos azules oscuro). \textbf{c)} Sección transversal diferencial de esparcimiento $\dd{C_{\text{sca}}}/\dd{\Omega}$ en función de $\theta$ para dos cortes: $\varphi=0^{\circ}$ (puntos azules claro) y $\varphi=90^{\circ}$ (puntos azules oscuro). La línea continua que une los puntos es una guía visual.}
		\label{fig:FF_431_health}
	\end{figure}
	%
	
Para cuantificar tanto la distribución angular como la magnitud de esparcimiento en cada región, se calculó el coeficiente de correlación de Pearson\footnote{
	El coeficiente de correlación de Pearson entre dos conjuntos de datos $x$ y $y$ se define como \cite{rossIntroductionProbabilityStatistics2014}
	$$r=\dfrac{\sum_{i}{(x_i-\overline{x})(y_i-\overline{y})}}{\sqrt{\sum_{i}{(x_i-\overline{x})^2} \sum_{i}{(y_i-\overline{y})^2}}},$$
	donde $\overline{x}$ indica el promedio del conjunto de datos. Por definición, $r$ es un valor entre $-1$ y 1, donde su valor absoluto es una medida de la fuerza de la relación lineal entre los dos pares de datos, mientras que su signo indica la dirección de la relación lineal \cite{rossIntroductionProbabilityStatistics2014}.} % 
	%
	% --- JAUA ---
	% Te adelanto que alo mejor vas a tener que explicar que no significa que sean iguales, pero que siguen mismas tendencias (la ecuación te puede aydar en donde además puedes aprovechar la dependencia en x. Creo que en general pensamos en mínimos cuaddrrados donde eso no ocurre.)
	%
entre ambos planos de observación y se integró $\dd{C_{\text{sca}}}/\dd{\Omega}$ en los intervalos correspondientes al esparcimiento frontal, lateral y de retroesparcimiento. Ambas cantidades proporcionan información complementaria: mientras que el coeficiente de Pearson permite comparar la estructura angular de los patrones, independientemente de diferencias globales de amplitud, la integral de la sección transversal diferencial permite cuantificar el campo esparcido. En la región frontal se obtiene $r=1$, lo que indica que las posiciones y amplitudes relativas de los máximos y mínimos producidos por interferencia presentan la misma dependencia angular en ambos planos de observación. Asimismo, las contribuciones integradas normalizadas por el área $S=4.41\times 10^{7}\,\text{nm}^2$ son $1,817.17\times 10^{-4}$ y $1,788.52\times 10^{-4}$ para $\varphi=0^{\circ}$ y $\varphi=90^{\circ}$, respectivamente, por lo que no solo la estructura angular, sino también la magnitud del esparcimiento frontal, presenta una baja dependencia con el plano de observación.  Este mismo caso ocurre en la región de retroesparcimiento, donde $r=0.99$ y las contribuciones integradas son iguales a $2.01\times 10^{-4}$ y $2.97\times 10^{-4}$ para $\varphi=0^{\circ}$ y $\varphi=90^{\circ}$, respectivamente. En contraste, en la región lateral $r= 0.67$ y las contribuciones integradas son de $11.71\times 10^{-4}$ y $21.70\times 10^{-4}$ para $\varphi=0^{\circ}$ y $\varphi=90^{\circ}$, respectivamente. En esta región, las diferencias tanto en la estructura de máximos y mínimos como en la magnitud integrada indican una mayor dependencia del patrón de esparcimiento con el plano de observación asociado a la polarización lineal. A partir de estos resultados, se observa que la contribución frontal supera hasta dos órdenes de magnitud a las contribuciones lateral y de retroesparcimiento, lo que muestra la direccionalidad del esparcimiento del eritrocito. Este comportamiento puede asociarse con que el campo esparcido por el eritrocito resulta de la superposición de las contribuciones generadas por la distribución de la polarización inducida en toda la célula y, debido a que sus dimensiones son mayores que la longitud de onda incidente, existen retardos de fase apreciables entre las contribuciones provenientes de distintas regiones del eritrocito. En la dirección frontal, las diferencias de fase son menores, favoreciendo una superposición predominantemente constructiva, mientras que al aumentar el ángulo de esparcimiento, las diferencias de fase entre las distintas contribuciones se incrementan, produciendo mayor interfencia destructiva y en consecuencia, la reducción del esparcimiento lateral y posterior.
	
	
	

En la Tabla~\ref{table:val_eri} se resumen, para las tres longitudes de onda representativas, los valores del factor de anisotropía $g$, la correlación en la región de esparcimiento lateral entre los cortes $\varphi=0^\circ$ y $\varphi=90^\circ$, así como las contribuciones de la sección transversal de esparcimiento integradas en las regiones frontal, lateral y de retroesparcimiento para ambos cortes. La correlación se presenta únicamente para la región lateral debido a que, en las regiones frontal y de retroesparcimiento, los patrones permanecen invariantes entre ambos cortes, con coeficientes de correlación de 1 y 0.99, respectivamente. Se observa que el factor de anisotropía aumenta de 0.97 a 0.99 al pasar de la banda de Soret a la banda $\alpha$, lo que es consistente con el predominio del esparcimiento frontal en todo el intervalo espectral analizado. Asimismo, la contribución integrada en la región frontal aumenta con la longitud de onda, mientras que las regiones lateral y de retroesparcimiento continúan aportando una fracción menor al esparcimiento total. Por otra parte, la correlación en la región lateral no presenta una tendencia creciente con la longitud de onda, alcanzando su valor mínimo en la banda $\beta$ ($r=0.49$), lo que indica que en esta región espectral la distribución angular del esparcimiento es la más sensible al plano de observación. Sin embargo, la diferencia entre las contribuciones integradas entre los dos planos es inferior (38.93 \%) para esta misma longitud de onda que para $\lambda=431$~nm (85.31 \%) y $\lambda=573$~nm (98.35 \%). Esto evidencia que una menor correlación entre dos planos de observación no implica necesariamente una mayor diferencia en su contribución en $C_{\text{sca}}$.

%
\begin{table}[h!]
	\begin{tabular}{ccccccccc}
		\hline
		\multirow{3}{*}{$\lambda$ [nm]} & \multirow{3}{*}{$g$} & \multirow{3}{5em}{Correlación lateral} & \multicolumn{6}{c}{$C_{\text{sca}}/S\times 10^{-4}$ } \\
		\cline{4-9}
		& &  &   \multicolumn{2}{c}{Frontal} & \multicolumn{2}{c}{Lateral} & \multicolumn{2}{c}{Posterior}\\
		\cline{4-9}
		&  &  &  $\varphi=0^{\circ}$ & $\varphi=90^{\circ}$ &$\varphi=0^{\circ}$ & $\varphi=90^{\circ}$ &$\varphi=0^{\circ}$ & $\varphi=90^{\circ}$ \\ \hline
		431 (Soret) & 0.97  & 0.67 & 1,817.17 & 1,788.52 & 11.71 & 21.70 & 2.01 & 2.97  \\
		541 ($\beta$) & 0.98 & 0.49 & 4,324.20 & 4,296.07 & 15.08 & 20.95 & 10.78 & 13.77\\
		573 ($\alpha$) & 0.99 & 0.71 & 4,650.05 & 4,626.42 & 12.75 & 25.29 & 2.46 & 3.41\\
		\hline
	\end{tabular}
	\caption{Valores del factor de anisotropía $g$, la correlación en la región de esparcimiento lateral entre los cortes $\varphi=0^\circ$ y $\varphi=90^\circ$ y las contribuciones de la sección transversal de esparcimiento integradas en las regiones frontal, lateral y de retroesparcimiento para $\varphi=0^{\circ}$ y $\varphi=90^{\circ}$ para $\lambda=431,541$ y 573 nm. }
	\label{table:val_eri}
	\vspace{0.2cm} 
\end{table}
%




       
% !TeX root = ../tesis.tex

\section{Esferocito}
\label{section:Res_Esferocito}

%
% --- JAUA ---
%
%Mismos cambios que en laseccion anterior sobre el epsilon y lo de la  concentrciobn
Para identificar las propiedades ópticas de los esferocitos en las bandas de Soret, $\beta$ y $\alpha$ se consideró el rango espectral de 400  a 600 nm. El esferocito se modeló como una esfera de radio $2.68\,\mu\text{m}$ y volumen igual a 80.63~$\mu\text{m}^3$, equivalente al caso límite de un óvalo de Cassini con valores de $a=0\,\mu\text{m}$, $b=2.68\,\mu\text{m}$ y $c=1\,\mu\text{m}$ como se muestra en rojo en la Fig. \ref{subfig:Esferocito_Shape}, donde además se muestra en azul la comparación de la sección transversal con la de un eritrocito sano con el mismo volumen en color azul analizado en la Sección \ref{section:Res_Eritrocito_Sano}. Al igual que en el caso del eritrocito sano, se consideró al plasma con una función dieléctrica constante de $\varepsilon_{\text{Plasma}}=1.839$ como el medio circundante, mientras que para el esferocito se empleó la función dieléctrica correspondiente a una CHCM de 34~g/dL. El esferocito se iluminó con una onda plana polarizada en $x$ y que se propaga en el eje $z$.

En la Fig. \ref{subfig:Qext_esf} se muestra la eficiencia de extinción $Q_{\text{ext}}$  en función de la longitud de onda $\lambda$ de la onda incidente para el esferocito (puntos rojos) y del eritrocito sano (puntos azules), donde la línea continua que los une es una guía visual. Se resaltan en gris las bandas que rodean los mínimos asociados a la banda de Soret y a las bandas $\beta$ y $\alpha$ del eritrocito sano, cuyos mínimos se localizan en 431 nm, 541 nm y 573 nm, respectivamente. Se observa que la respuesta espectral del esferocito conserva la misma estructura general, pero los mínimos de $Q_{\text{ext}}$ presentan un corrimiento hacia el rojo con respecto a los del eritrocito sano, localizándose en 438 nm, 542 nm y 577 nm. Asimismo, se observa que en el rango espectral de 400 nm a 561 nm el esferocito presenta valores de $Q_{\text{ext}}$ superiores al eritrocito sano. Estas diferencias no pueden atribuirse de manera independiente a un único parámetro, ya que entre ambos modelos cambian simultáneamente la geometría celular y la CHCM. Por un lado, la modificación de la CHCM altera el índice de refracción complejo y, con ello, tanto la absorción como el contraste óptico entre la célula y el medio. Por otro lado, el cambio de geometría modifica tanto el área proyectada empleada en la definición de $Q_{\text{ext}}$ como la distribución espacial y de fase del campo electromagnético en el interior de la célula. Por tanto, la respuesta espectral observada resulta del efecto conjunto del cambio de las propiedades ópticas y geométricas del esferocito.
%
\begin{figure}[t]
	\centering
	\sidesubfloat[Second image]{\hspace{-0.3cm}{
			\includegraphics[width=0.45\textwidth]{4-Resultados/Eritrocito_sano_esferocito_corte.pdf}\label{subfig:Esferocito_Shape}}}\hspace{0.1cm}
	\sidesubfloat[First image]{\hspace{-0.5cm}{	\includegraphics[width=0.47\textwidth]{4-Resultados/Qext_Esf.pdf}\label{subfig:Qext_esf}}}
	\vspace{-0.1cm}
	\caption{\textbf{a)} Comparación entre el modelo del eritrocito sano (azul) con valores de $a=2.6\,\mu\text{m}$, $b=2.7\,\mu\text{m}$ y $c=0.812\,\mu\text{m}$ y el modelo de esferocito con radio igual a $2.68\,\mu\text{m}$, que es equivalente a un óvalo de Cassini con valores de  $a=0\,\mu\text{m}$, $b=2.68\,\mu\text{m}$ y 
		$c=1\,\mu\text{m}$. \textbf{b)} Eficiencia de extinción $Q_{\text{ext}}$ en función de la longitud de onda $\lambda$ de la onda incidente para un eritrocito (puntos azules) y un esferocito (puntos rojos) orientados con su eje de simetría en el eje $z$ iluminado con una onda en dirección del eje $z$ y polarizada en $x$. La línea continua que une los puntos es una guía visual. Se resaltan en gris las bandas que 
		rodean los mínimos asociados a la banda de Soret y a las bandas Q.}
	\label{fig:Shape_Qext_esf}
\end{figure}
%

En esta sección se discute únicamente el patrón de esparcimiento de $\lambda=431$~nm asociado a la banda de Soret, mientras que los de $\lambda=541$ y $573$ nm se incluyen en el Apéndice \ref{section:apendix4}. En la Fig.~\ref{subfig:FF_431_esf} se muestra la sección transversal diferencial de esparcimiento $\dd{C_{\text{sca}}}/\dd{\Omega}$ para $\lambda=431$ nm. La magnitud de $\dd{C_{\text{sca}}}/\dd{\Omega}$ se encuentra normalizada respecto a su valor máximo, que corresponde al valor más saturado. A pesar del cambio simultáneo de la geometría bicóncava del eritrocito sano a una geometría esférica y del incremento de la CHCM, el esparcimiento predomina hacia adelante. En particular, el factor de anisotropía aumenta de $g=0.97$ a $g=0.98$.

En la Fig. \ref{subfig:FF_polar_431_esf} se muestra $\dd{C_{\text{sca}}}/\dd{\Omega}$ en decibeles en función del ángulo de esparcimiento $\theta$ para los cortes $\varphi=0^{\circ}$ (puntos rojo claro) y $\varphi=90^{\circ}$ del esferocito (puntos rojo oscuro), comparados con los del eritrocito sano (puntos rojo claro y puntos rojo escuro para $\varphi=0^{\circ}$ y $\varphi=90^{\circ}$, respectivamente). Al igual que lo observado con el eritrocito, el patrón de esparcimiento depende de la polarización de la onda incidente. En particular, la correlación entre los cortes $\varphi=0^{\circ}$  y $\varphi=90^{\circ}$ del esferocito en la región lateral aumenta a $r=0.78$, lo que indica una mayor similitud entre las estructuras angulares observadas en ambos planos en comparación con el eritrocito sano. En contraste, en la región de retroesparcimiento la correlación disminuye de $r=0.99$ a $r=0.95$.

En la Fig. \ref{subfig:FF_cartesian_431_esf} se muestra $\dd{C_{\text{sca}}}/\dd{\Omega}$ en función del ángulo de esparcimiento $\theta$ para ambos cortes. En la región lateral, el patrón de esparcimiento del eritrocito sano presenta una envolvente entre $\theta=93.5^{\circ}$ y $\theta=120^{\circ}$ sobre la que se superponen pequeñas oscilaciones, mientras que en el esferocito dicha envolvente disminuye y las oscilaciones se atenúan.
%
\begin{figure}[h!]
	\centering
	\sidesubfloat[First image]{\hspace{-0.5cm}{	\includegraphics[width=7.3cm]{4-Resultados/FFdistribution431_Esf.pdf}\label{subfig:FF_431_esf}}}\hspace{0.1cm}
	\sidesubfloat[Second image]{\hspace{-0.5cm}{
			\includegraphics[width=7.5cm]{4-Resultados/polarplot431FFComEsfEri.pdf}\label{subfig:FF_polar_431_esf}}}\vspace{-0.1cm}
	\\
	\vspace{0.7cm}
	
	\sidesubfloat[Third image]{\hspace{0.3cm}{
			\includegraphics[width=14cm]{4-Resultados/linearplot431FFEsf.pdf}\label{subfig:FF_cartesian_431_esf}}}\vspace{-0.1cm}
	\caption{\textbf{a)} Sección transversal diferencial de esparcimiento $\dd{C_{\text{sca}}}/\dd{\Omega}$ normalizada respecto a su valor máximo para una longitud de onda $\lambda=431$ nm en función del ángulo de esparcimiento $\theta$ y del ángulo $\varphi$ para un esferocito. Se identifican dos planos $\varphi=0^{\circ}$ en rojo claro y $\varphi=90^{\circ}$ en rojo oscuro. \textbf{b)} Sección transversal diferencial de esparcimiento $\dd{C_{\text{sca}}}/\dd{\Omega}$ en decibeles en función de $\theta$ para un esferocito (en rojo) y un eritrocito sano (azul) para dos cortes: $\varphi=0^{\circ}$ (puntos color claro) y $\varphi=90^{\circ}$ (puntos color oscuro). \textbf{c)} Sección transversal diferencial de esparcimiento $\dd{C_{\text{sca}}}/\dd{\Omega}$ en función de $\theta$ para un esferocito (en rojo) y un eritrocito sano (azul) para dos cortes: $\varphi=0^{\circ}$ (puntos color claro) y $\varphi=90^{\circ}$ (puntos color oscuro). }
	\label{}
\end{figure}
%

Para cuantificar las diferencias en los patrones de esparcimiento se integró $\dd{C_{\text{sca}}}/\dd{\Omega}$ en las regiones frontal, lateral y de retroesparcimiento, obteniéndose contribuciones normalizadas por el área $S=2.25\times 10^{7}\,\text{nm}^2$ de $2,030.46\times 10^{-4}$, $7.61\times 10^{-4}$ y $0.47\times 10^{-4}$ para el corte $\varphi=0^{\circ}$ y de $2,019.00\times~10^{-4}$, $14.45~\times~10^{-4}$ y $0.90\times 10^{-4}$ para $\varphi=90^{\circ}$. En comparación con el eritrocito sano, las contribuciones lateral y de retroesparcimiento disminuyen, en promedio, por un factor de 0.66 y 0.27, respectivamente, mientras que el cambio en el esparcimiento frontal es menor (factor de cambio de 1.12). Esto indica que el incremento del factor de anisotropía no se debe a un aumento significativo del esparcimiento frontal, sino a una reducción  del esparcimiento hacia ángulos laterales y posteriores.

Con el fin de determinar si las modificaciones observadas para $\lambda=431$ nm se mantienen en las bandas $\beta$ y $\alpha$, en la Tabla~\ref{table:val_esf} se resumen, para las tres longitudes de onda representativas, los valores del factor de anisotropía $g$, la correlación en la región de esparcimiento lateral y posterior entre los cortes $\varphi=0^\circ$ y $\varphi=90^\circ$, así como las contribuciones de la sección transversal de esparcimiento integradas en las regiones frontal, lateral y de retroesparcimiento para ambos cortes. En las tres longitudes de onda analizadas, el esferocito mantiene el esparcimiento predominantemente hacia adelante con valores de $g$ mayores a 0.98. Estos resultados muestran que las modificaciones relativas más importantes entre el eritrocito sano y el esferocito se localizan fuera de la región frontal.

Por otra parte, los coeficientes de correlación lateral del esferocito son superiores a los obtenidos para el eritrocito sano en las tres longitudes de onda consideradas, lo que indica una mayor similitud entre las estructuras angulares de los planos $\varphi=0^\circ$ y $\varphi=90^\circ$ en esta región. A diferencia del eritrocito sano, para el esferocito se incluye también la correlación en la región de retroesparcimiento, ya que los valores obtenidos se presentan variaciones entre las longitudes de onda consideradas. Aunque las mayores diferencias relativas en la magnitud del esparcimiento entre ambas morfologías se observan en la región posterior, esta región presenta también las menores contribuciones absolutas. Por tanto, para identificar intervalos angulares potencialmente útiles para la discriminación celular es necesario considerar conjuntamente la diferencia entre su magnitud y su amplitud misma.

%
\begin{table}[h!]
	\setlength{\tabcolsep}{4pt}
	\begin{tabular}{cccccccccc}
		\hline
		\multirow{3}{*}{$\lambda$ [nm]} & \multirow{3}{*}{$g$} & \multicolumn{2}{c}{Correlación} & \multicolumn{6}{c}{$C_{\text{sca}}/S\times 10^{-4}$ } \\
		\cline{5-10}
		& & \multirow{2}{*}{Lateral} &  \multirow{2}{*}{Posterior}&  \multicolumn{2}{c}{Frontal} & \multicolumn{2}{c}{Lateral} & \multicolumn{2}{c}{Posterior}\\
		\cline{5-10}
		&  &  &  &  $\varphi=0^{\circ}$ & $\varphi=90^{\circ}$ &$\varphi=0^{\circ}$ & $\varphi=90^{\circ}$ &$\varphi=0^{\circ}$ & $\varphi=90^{\circ}$ \\ \hline
		431 (Soret) & 0.98  & 0.78 & 0.95 & 2,030.46 & 2,019.00 & 7.61 & 14.45 & 0.47 & 0.90  \\
		541 ($\beta$) & 0.99 & 0.90 & 0.93 & 4,271.44 & 4,253.33 & 7.89 & 17.66 & 0.63 & 1.12\\
		573 ($\alpha$) & 0.99 & 0.86 & 0.79 & 4,616.60 & 4,599.09 & 7.38 & 17.12 & 0.58 & 1.09\\
		\hline
	\end{tabular}
	\caption{Valores del factor de anisotropía $g$, la correlación en la región de esparcimiento lateral entre los cortes $\varphi=0^\circ$ y $\varphi=90^\circ$ y las contribuciones de la sección transversal de esparcimiento integradas en las regiones frontal, lateral y de retroesparcimiento para $\varphi=0^{\circ}$ y $\varphi=90^{\circ}$ para $\lambda=431,541$ y 573 nm para un esferocito con CHCM de 34 g/dL y 2,680 nm de radio. }
	\label{table:val_esf}
	\vspace{0.2cm} 
\end{table}
%

% !TeX root = ../tesis.tex

\section{Microcito}
\label{section:Res_Microcito}

Para modelar el microcito se empleó un óvalo de Cassini con parámetros de $a=2.418\,\mu\text{m}$, $b=2.511\,\mu\text{m}$ y $c=0.756\,\mu\text{m}$, los cuales están escalados 0.98 veces en comparación con los del eritrocito sano. La geometría resultante presenta un radio equivalente de 3.67 $\mu$m, un volumen de 64.31~$\mu\text{m}^3$ y una sección transversal como la de la Fig. \ref{subfig:Esferocito_Shape} (verde), donde además se muestra su comparación con la sección transversal del eritrocito sano (azul) analizado en la Sección \ref{section:Res_Eritrocito_Sano} y con la del esferocito (rojo) analizado en la Sección \ref{section:Res_Esferocito}. Al igual que en los casos anteriores, se consideró al plasma como medio circundante con una función dieléctrica constante $\varepsilon_{\text{Plasma}}=1.839$, mientras que al microcito se le asignó la función dieléctrica de una CHCM de 28.7~g/dL. El microcito se orientó con su eje de simetría a lo largo del eje $z$ y se iluminó con una onda plana polarizada en $x$ y que se propaga en el eje $z$.

En la Fig. \ref{subfig:Qext_micro} se muestra la eficiencia de extinción $Q_{\text{ext}}$  en función de la longitud de onda $\lambda$ de la onda incidente para el microcito (puntos verdes), el esferocito (puntos rojos) y el eritrocito sano (puntos azules); donde la línea continua es una guía visual. También se indican en gris las bandas que rodean los mínimos asociados a la banda de Soret y a las bandas $\beta$ y $\alpha$ del eritrocito sano, cuyos mínimos se localizan en 431 nm, 541 nm y 573 nm, respectivamente. Se observa que el mínimo asociado a la banda de Soret tiene un corrimiento hacia el azul, localizándose en 418 nm, mientras que los asociados a las bandas $\beta$ y $\alpha$ presentan un corrimiento hacia el rojo, localizándose en 547 nm y 580 nm, respectivamente. Asimismo, la dependencia espectral de $Q_{\text{ext}}$ conserva la forma general observada para el eritrocito sano, aunque desplazada hacia valores mayores de $Q_{\text{ext}}$. Además, se observa que los mínimos asociados a las bandas $\beta$ y $\alpha$ son menos profundos.
%
\begin{figure}[h!]
	\centering
	\sidesubfloat[Second image]{\hspace{-0.3cm}{
			\includegraphics[width=0.45\textwidth]{4-Resultados/Eritrocito_sano_esferocito_microcito_corte.pdf}\label{subfig:Microcito_Shape}}}\hspace{0.1cm}
	\sidesubfloat[First image]{\hspace{-0.5cm}{	\includegraphics[width=0.47\textwidth]{4-Resultados/Qext_Micro.pdf}\label{subfig:Qext_micro}}}
	\vspace{-0.1cm}
	\caption{\textbf{a)} Comparación entre el modelo del microcito (verde) con valores de $a=2.418\,\mu\text{m}$, $b=2.511\,\mu\text{m}$ y $c=0.756\,\mu\text{m}$, que corresponden a valores de un eritrocito sano (azul) escalados 0.98 veces y el modelo de esferocito con radio igual a $2.68\,\mu\text{m}$, que es equivalente a un óvalo de Cassini con valores de  $a=0\,\mu\text{m}$, $b=2.68\,\mu\text{m}$ y 
		$c=1\,\mu\text{m}$. \textbf{b)} Eficiencia de extinción $Q_{\text{ext}}$ en función de la longitud de onda $\lambda$ de la onda incidente para un eritrocito (puntos azules), un esferocito (puntos rojos) y un microcito (puntos verdes) orientados con su eje de simetría en el eje $z$ iluminado con una onda en dirección del eje $z$ y polarizada en $x$. La línea continua que une los puntos es una guía visual. Se resaltan en gris las bandas que 
		rodean los mínimos asociados a la banda de Soret y a las bandas Q.}
	\label{fig:Shape_Qext_micro}
\end{figure}
%

Como se observa en la Fig.~\ref{subfig:Cext_micro}, el microcito también presenta la mayor sección transversal de extinción $C_{\text{ext}}$, lo que indica que el incremento de $Q_{\text{ext}}$ no puede atribuirse únicamente a la normalización respecto al área proyectada. Para comprender este comportamiento, en la Figs.~\ref{subfig:Csca_micro} y \ref{subfig:Cabs_micro} se muestran las secciones transversales de esparcimiento $C_{\text{sca}}$ y de absorción $C_{\text{abs}}$ en función de la longitud de onda de la luz incidente para el eritrocito sano (azul), el esferocito (rojo) y el microcito (azul). Se observa que en el intervalo de 400 a 566 nm el microcito presenta los mayores valores de $C_{\text{sca}}$, mientras que $C_{\text{abs}}$ es menor en todo el intervalo espectral. Este comportamiento se explica por la disminución de la CHCM, que reduce la componente imaginaria del índice de refracción y, por tanto, la atenuación que experimenta la onda electromagnética al propagarse dentro de la célula. Como consecuencia, una mayor fracción del campo atraviesa el eritrocito y contribuye al esparcimiento, incrementando $C_{\text{sca}}$. En contraste, el incremento de la CHCM favorece la absorción de energía por la hemoglobina, dando lugar a un aumento de $C_{\text{abs}}$. En el caso del microcito, el aumento del esparcimiento supera 62.15 veces en promedio la disminución de la absorción, lo que explica que presente la mayor sección transversal de extinción $C_{\text{ext}}$ y que al normalizar por el área proyectada $Q_{\text{ext}}$ sea mayor.

%
\begin{figure}[t]
	\centering
	\sidesubfloat[Second image]{\hspace{-0.3cm}{
			\includegraphics[width=0.47\textwidth]{4-Resultados/Cext_micro.pdf}\label{subfig:Cext_micro}}}\hspace{0.1cm}
	\sidesubfloat[First image]{\hspace{-0.5cm}{	\includegraphics[width=0.47\textwidth]{4-Resultados/Cscamicro.pdf}\label{subfig:Csca_micro}}}
	\vspace{-0.1cm}
	\\
	\vspace{0.7cm}
	
	\sidesubfloat[Third image]{\hspace{0.3cm}{
			\includegraphics[width=0.47\textwidth]{4-Resultados/Cabs_micro.pdf}\label{subfig:Cabs_micro}}}\vspace{-0.1cm}
	\caption{Secciones transversales en función de la longitud de onda $\lambda$ de la onda incidente para un eritrocito (puntos azules), un esferocito (puntos rojos) y un microcito (puntos verdes) orientados con su eje de simetría en el eje $z$ iluminado con una onda en dirección del eje $z$ y polarizada en $x$. \textbf{a)} Sección transversal de extinción $C_{\text{ext}}$. \textbf{b)} Sección transversal de esparcimiento $C_{\text{abs}}$. \textbf{c)} Sección transversal de absorción $C_{\text{abs}}$.}
	\label{fig:ST_micro}
\end{figure}
%

En la Fig.~\ref{subfig:FF_431_micro} se muestra la sección transversal diferencial de esparcimiento $\dd{C_{\text{sca}}}/\dd{\Omega}$ para $\lambda=431$ nm normalizada respecto a su valor máximo, que corresponde al valor más saturado. Al igual que en el eritrocito sano y el esferocito, el esparcimiento se concentra predominantemente en la dirección frontal de propagación. No obstante, el microcito exhibe el mayor factor de anisotropía ($g=0.99$), lo que indica que el campo esparcido se concentra aún más alrededor de la dirección de propagación. Lo anterior se debe a que al disminuir la CHCM, el contraste de índice de refracción entre el medio y la célula se reduce, incrementando la transparencia de la célula.

En la Fig. \ref{subfig:FF_polar_431_micro} se muestra $\dd{C_{\text{sca}}}/\dd{\Omega}$ en decibeles en función del ángulo de esparcimiento $\theta$ para los cortes $\varphi=0^{\circ}$ (puntos verde claro) y $\varphi=90^{\circ}$ de la esfera (puntos verde oscuro), comparados con los del eritrocito sano (puntos azul claro y puntos azul escuro para $\varphi=0^{\circ}$ y $\varphi=90^{\circ}$, respectivamente). A diferencia del esferocito, el microcito presenta una mayor dependencia del patrón de esparcimiento con la polarización de la onda incidente. En la región lateral, la correlación entre $\varphi=0^{\circ}$  y $\varphi=90^{\circ}$ disminuye hasta $r=0.56$, lo que indica que las distribuciones angulares correspondientes a ambas polarizaciones son menos similares que en el eritrocito sano. Esta diferencia es aún más evidente en la región de retroesparcimiento, donde la correlación se reduce de $r=0.99$ a $r=0.76$, evidenciando que la respuesta angular del microcito es significativamente más sensible a la polarización. 

En la Fig. \ref{subfig:FF_cartesian_431_esf} se muestra $\dd{C_{\text{sca}}}/\dd{\Omega}$ en función del ángulo de esparcimiento $\theta$ para ambos cortes. En comparación con el eritrocito sano, el microcito conserva la estructura general del patrón de esparcimiento; por ejemplo, en la región lateral ambos patrones presentan una envolvente amplia entre  $\theta=90^{\circ}$ y $\theta=120^{\circ}$ aproximadamente. Físicamente, la disminución de la CHCM reduce tanto la absorción como el contraste del índice de refracción entre la célula y el plasma que, aunados a la reducción del volumen y del radio interior, generan una menor acumulación de fase en la onda transmitida. Consecuentemente, las condiciones de interferencia se modifican, produciendo variaciones locales en la distribución angular (tales como ligeros desplazamientos angulares) sin alterar la estructura global del patrón de interferencia, que sigue dominado por la morfología bicóncava.
%
\begin{figure}[h!]
	\centering
	\sidesubfloat[First image]{\hspace{-0.5cm}{	\includegraphics[width=7.3cm]{4-Resultados/FFdistribution431_Micro.pdf}\label{subfig:FF_431_micro}}}\hspace{0.1cm}
	\sidesubfloat[Second image]{\hspace{-0.5cm}{
			\includegraphics[width=7.5cm]{4-Resultados/polarplot431FFComEsfMicro.pdf}\label{subfig:FF_polar_431_micro}}}\vspace{-0.1cm}
	\\
	\vspace{0.7cm}
	
	\sidesubfloat[Third image]{\hspace{0.3cm}{
			\includegraphics[width=14cm]{4-Resultados/linearplot431FFMicro.pdf}\label{subfig:FF_cartesian_431_micro}}}\vspace{-0.1cm}
	\caption{\textbf{a)} Sección transversal diferencial de esparcimiento $\dd{C_{\text{sca}}}/\dd{\Omega}$ normalizada respecto a su valor máximo para una longitud de onda $\lambda=431$ nm en función del ángulo de esparcimiento $\theta$ y del ángulo $\varphi$ para un microcito. Se identifican dos planos $\varphi=0^{\circ}$ en verde claro y $\varphi=90^{\circ}$ en verde oscuro. \textbf{b)} Sección transversal diferencial de esparcimiento $\dd{C_{\text{sca}}}/\dd{\Omega}$ en decibeles en función de $\theta$ para un microcito (en verde) y un eritrocito sano (azul) para dos cortes: $\varphi=0^{\circ}$ (puntos color claro) y $\varphi=90^{\circ}$ (puntos color oscuro). \textbf{c)} Sección transversal diferencial de esparcimiento $\dd{C_{\text{sca}}}/\dd{\Omega}$ en función de $\theta$ para un microcito (en verde) y un eritrocito sano (azul) para dos cortes: $\varphi=0^{\circ}$ (puntos color claro) y $\varphi=90^{\circ}$ (puntos color oscuro). }
	\label{fig:FF_431_micro}
\end{figure}
%

Para cuantificar el efecto de estas variaciones locales en la distribución angular, se integró $\dd{C_{\text{sca}}}/\dd{\Omega}$ en las regiones frontal, lateral y de retroesparcimiento, obteniéndose contribuciones normalizadas por el área $S$ de $3557.07\times 10^{-4}$, $11.56\times 10^{-4}$ y $3.57\times 10^{-4}$ para el corte $\varphi=0^{\circ}$, y de $3542.14\times 10^{-4}$, $19.21\times 10^{-4}$ y $5.99\times 10^{-4}$ para $\varphi=90^{\circ}$. En comparación con el eritrocito sano, la contribución lateral disminuye, en promedio, por un factor de 0.38, mientras que las contribuciones frontal y de retroesparcimiento aumentan por factores de 1.67 y 1.61, respectivamente. En consecuencia, el incremento del factor de anisotropía del microcito se asocia principalmente con un aumento del campo esparcido en la dirección frontal, acompañado de una reducción del esparcimiento lateral. 


La Fig.~\ref{subfig:STLF0Pi2} resume las contribuciones integradas de la sección transversal diferencial de esparcimiento normalizadas por el área en las regiones frontal (STF, Sección Transversal Frontal) y lateral (STL, Sección Transversal Lateral) para $\lambda=431$ nm. Los marcadores cuadrados y los circulares corresponden a los cortes $\varphi=0^{\circ}$ y $\varphi=90^{\circ}$, respectivamente. El color de cada marcador identifica el tipo de eritrocito (azul: sano, rojo: esferocito y verde: microcito), mientras que el tamaño del marcador es proporcional a la contribución integrada del retroesparcimiento.
%
\begin{figure}[h!]
	\centering
	\sidesubfloat[Second image]{\hspace{-0.3cm}{
			\includegraphics[width=0.47\textwidth]{4-Resultados/STLF0Pi2.pdf}\label{subfig:STLF0Pi2}}}\hspace{0.1cm}
	\sidesubfloat[First image]{\hspace{-0.3cm}{	\includegraphics[width=0.43\textwidth]{4-Resultados/FSCpaper.pdf}\label{subfig:STLFPi2}}}
	\vspace{-0.1cm}
	\caption{\textbf{a)} \textbf{b)} Eficiencia de extinción $Q_{\text{ext}}$ en función de la longitud de onda $\lambda$ de la onda incidente para un eritrocito (puntos azules), un esferocito (puntos rojos) y un microcito (puntos verdes) orientados con su eje de simetría en el eje $z$ iluminado con una onda en dirección del eje $z$ y polarizada en $x$. La línea continua que une los puntos es una guía visual. Se resaltan en gris las bandas que 
		rodean los mínimos asociados a la banda de Soret y a las bandas Q.}
	\label{fig:STLF}
\end{figure}
%

En ambos cortes se observa que el microcito presenta la mayor contribución en la región frontal, seguido del esferocito y del eritrocito sano. Este resultado es consistente con el incremento del factor de anisotropía obtenido para el microcito, indicando que una mayor fracción de la energía es esparcida en la dirección de propagación. En contraste, la contribución lateral muestra el efecto de la geometría celular: el eritrocito sano y el microcito, ambos con geometría bicóncava, presentan contribuciones laterales comparables, mientras que el esferocito exhibe una reducción significativa de esta contribución como consecuencia de la pérdida de la concavidad central. Esto sugiere que el esparcimiento lateral es particularmente sensible a la morfología de la célula. Finalmente, en cuanto al tamaño de los marcadores, se observa que el retroesparcimiento aumenta desde el eritrocito sano hasta el esferocito y alcanza su valor máximo en el microcito. Esta representación muestra que las modificaciones en la CHCM afectan principalmente el esparcimiento frontal y el retroesparcimiento, mientras que la geometría ejerce mayor influencia sobre la contribución lateral. Basado en lo obtenido, estas magnitudes constituyen parámetros útiles para la diferenciación de eritrocitos sanos y enfermos.






% !TeX root = ../tesis.tex

\section{Intervalos angulares}
\label{section:angle_interval}

Una vez caracterizados los patrones de esparcimiento y las propiedades ópticas de los eritrocitos sanos y enfermos, se realizó un primer análisis con el fin de identificar las regiones angulares en las que las diferencias entre las distintas morfologías son mayores. %
% --- JAUA ---
% Hay que ligar con lo anterior, no es un párrafo aparte sin o una secicón disitna.
% Usas palabras de relleno y hay que qiuitarlas, por ejemplo:
%		se realizó un primer análisis con el fin de identificar las regiones angulares en las que las diferencias entre las distintas morfologías son mayores
%  -> 	se identificaron regiones angulares con base en diferencias apreciables según la morfología de los eritrocitos.
%
% Aunque aún en esa frase no es claro de qié jablas: lo explicas más adelante entonces hay que reestrucurar el párrafo.
 Este análisis permite obtener una primera aproximación de los intervalos de $\theta$ potencialmente relevantes antes de considerar ventanas angulares tanto en $\theta$ como en $\varphi$. Para ello, se integró $\dd{C_{\text{sca}}}/\dd{\Omega}$ sobre ventanas de $6^{\circ}$ en $\theta$, desplazadas en pasos de $1^{\circ}$, %
 % --- JAUA ---
 % se calcullo  \dd{C_sca}\dd\varphi (\theta) considerando una integración alrededor de \theta de \pm 3 grados. Eso es más preciso, demás de que se nos pasó cmabiar la cantidad que graficas, no?? y el 'phi está evaluado según yo. EWse texto es más claro que lo de la ventana que descibes y además te ahorras la explicación. Según yo en estas no integraste con varphi entonces hay que dejar esa dependencia aún en el cociente que mencionas. no lo había pensado hasta ahora.
 %
 manteniendo fijos los planos de observación $\varphi=0^\circ$ y $\varphi=90^\circ$. Para cada ventana se calculó el cociente entre la contribución de esparcimiento correspondiente a la célula enferma y la obtenida para el eritrocito sano. Este cociente se define mediante el parámetro $R$ %
 % --- JAUA ---
 %  Selecciona una letra griega si queres, pero R puedde confundirse con reflectiancia.
 %
%
\begin{equation}
	R_j(\theta,\varphi)=\frac{C_{\text{sca}}^{\text{j}}}{C_{\text{sca}}^{\text{sano}}},\quad \text{con}\quad j=\{\text{esferocito, microcito}\}.
\end{equation}
%
De esta forma, $R_j=1$ corresponde a una contribución de esparcimiento igual a la del eritrocito sano, mientras que valores alejados de la unidad indican un mayor contraste entre ambas morfologías.

En la Fig.~\ref{fig:theta_esf_micro} se muestran los valores de $R_j$ obtenidos para el esferocito (rojo) y el microcito (verde) en los planos $\varphi=0^\circ$ y $\varphi=90^\circ$. \sout{Con el fin de identificar regiones de alto contraste,} se resaltan en rojo los intervalos angulares para los cuales $R_{\text{esferocito}}<0.5$, es decir, aquellos en los que la contribución de esparcimiento del esferocito es inferior a la mitad de la correspondiente al eritrocito sano.%
% --- JAUA ---
% yo pondría toda la explicación dela mitas y entre paréntesis al final que es R<.5
%
 De manera análoga, se resaltan en verde los intervalos para los que $R_{\text{microcito}}>2$, donde la contribución del microcito supera por al menos un factor de dos a la del eritrocito sano. Ambos criterios corresponden, por tanto, a diferencias de al menos un factor de dos respecto a la respuesta del eritrocito sano. %
 % --- JAUA ---
 % puedes resumir todo lo anterior si inicias con esya ultima frase el párrafo y pones las ecuaciones en línea. 
 %
%
\begin{figure}[h!]
	\centering
	\sidesubfloat[First image]{\hspace{-0.5cm}{	\includegraphics[width=0.45\textwidth]{4-Resultados/contriInt431.pdf}\label{subfig:theta_esf_micro_431}}}\hspace{0.1cm}
	\sidesubfloat[Second image]{\hspace{-0.5cm}{
			\includegraphics[width=0.45\textwidth]{4-Resultados/contriInt541.pdf}\label{subfig:theta_esf_micro_541}}}\vspace{-0.1cm}
	\\
	\vspace{0.7cm}
	
	\sidesubfloat[Third image]{\hspace{0.3cm}{
			\includegraphics[width=0.45\textwidth]{4-Resultados/contriInt573.pdf}\label{subfig:theta_esf_micro_573}}}\vspace{-0.1cm}
	\caption{Cociente $R_j$ entre la contribución angular de esparcimiento de los eritrocitos patológicos y la del eritrocito sano, obtenido mediante ventanas de $6^{\circ}$ en $\theta$ desplazadas en pasos de $1^{\circ}$, para los planos $\varphi=0^\circ$ y $\varphi=90^\circ$. Se muestran los resultados para el esferocito (rojo) y el microcito (verde) para las bandas de \textbf{a)} Soret, \textbf{b)} $\beta$ y \textbf{c)} $\alpha$. Las Las regiones sombreadas en rojo indican los intervalos para los que $R_{\text{esferocito}}<0.5$, mientras que las regiones sombreadas en verde corresponden a
		$R_{\text{microcito}}>2$}%
	% --- JAUA ---
	% Tienes espacio para poner un diagrama de la integración que estás haciendo; te recomiendo que la hagas t cambies la nomevkatura ahí para centrat en \theta e integrar en \pm 3 deg
	%	
	\label{fig:theta_esf_micro}
\end{figure}
%
%
% --- JAUA ---
% A lo mejor hacer mención que las grafias anteriores corresponden a lo que se espera de sistmas estándares de citometría, discutido en la Sección Xx (de  hecho, sería una excelente forma de iniciar esta sección). Luego, mencionar que para identificar ventanas en otros angulos expandes la busqueda en mas voles de varphi. Uy ya continuas lo que llevas aca.
%
Para la banda de Soret, las condiciones $R_{\text{esferocito}}<0.5$ y $R_{\text{microcito}}>2$ %
% --- JAUA ---
% mejor 1/2 y 2, para que sea obivio el criterio que propones
% 
se satisfacen simultáneamente entre $95^{\circ}$
y $98^{\circ}, 105^{\circ}$ y $106^{\circ}$, $120^{\circ}$ y $122^{\circ}$, y $130^{\circ}$ y $153^{\circ}$. %
% --- JAUA ---
% solo menciona que identificas una región enttre 130 7 153 donde se cumplen las dos conficiones; si bien para valores menores hay ventaas donde eso ocrrue (por ejemplo 95°) su angosta extensión la excluye para un esquema de medición (o algo así.) 
%
En la banda $\beta$, en cambio, ambas condiciones coinciden únicamente en el intervalo comprendido entre $129^{\circ}$ y $132^{\circ}$. Para la banda $\alpha$, la región de coincidencia se extiende aproximadamente entre $128^{\circ}$ y $150^{\circ}$. Al comparar las tres bandas, se observa que existe una región común entre $\theta=129^{\circ}$ y $\theta=132^{\circ}$, en la que las tres bandas presentan un contraste elevado entre el eritrocito sano y las dos morfologías de eritrocitos enfermos. Este resultado muestra que a esta región angular es una candidata para un análisis más detallado de discriminación celular. %
% --- JAUA ---
%
%Si quieres resaltar este punto de 129 común a los tres, márcalo con un color distinto.

%%A partir de aqui solo revisar las propuestas de las figuras

%
% --- JAUA ---
% en egenal bien. En Veo etiqyetas blancas al rededor del texto lo que se ve mal. También inclute en grande o arriba a qué banda (y longitu de onda) se refiere cada gráfica. Resatula con morado y magenta el pnuto comú nde 129 que mencipnaste. En el texto debería venir r_j, no r nada más. Según yo las etiquetas, un valro debe ser 90°, ero tines tres con 0°. Inssito que pueddes poner un diagrama de tus claculos trasladados a la medicion.

%
% --- JAUA ---
% de las otras gráficas. Siguiendo el esquema de colore, estas deberían ser verdes y rojas. Puedes hacer un blend con un color verde y la escala de blanco a negro, y lo mismo con uno rojo. Pero el morado desentona mucho y recueerda ak úbdice del plasma.
%
%

Con el fin de estudiar una configuración angular más general, el análisis se extendió a regiones con extensión finita en ambos ángulos. Para ello, se integró $\dd{C_{\text{sca}}}/\dd{\Omega}$ sobre ventanas angulares de $6^{\circ}$ en $\theta$ y $6^{\circ}$ en $\varphi$ y se realizó un barrido de dichas ventanas sobre los dos ángulos. Para cada posición se calculó nuevamente el cociente $R_j$.
%
\begin{figure}
	\includegraphics[width=0.75\textwidth]{4-Resultados/plot3DContourEsfSoret.pdf}
		\caption{Cociente $R_j$ entre la contribución angular de esparcimiento de los eritrocitos patológicos y la del eritrocito sano, obtenido mediante ventanas de $6^{\circ}$ en $\theta$ desplazadas en pasos de $1^{\circ}$, para los planos $\varphi=0^\circ$ y $\varphi=90^\circ$.}
\end{figure}
%


%
\begin{figure}
	\includegraphics[width=0.75\textwidth]{4-Resultados/plot3DContourMicroSoret.pdf}
	\caption{Cociente $R_j$ entre la contribución angular de esparcimiento de los eritrocitos patológicos y la del eritrocito sano, obtenido mediante ventanas de $6^{\circ}$ en $\theta$ desplazadas en pasos de $1^{\circ}$, para los planos $\varphi=0^\circ$ y $\varphi=90^\circ$.}
\end{figure}
%


% !TeX root = ../tesis.tex
\chapter*{Conclusiones}\addcontentsline{toc}{chapter}{\protect\numberline{}Conclusiones}
\label{chapter:concl}

En este trabajo se estudiaron numéricamente las propiedades ópticas de eritrocitos sanos y enfermos mediante el método de \textit{Finite Integration Technique} implementado en CST Studio Suite. Para ello, se modelaron eritrocitos sanos, esferocitos y microcitos considerando tanto sus diferencias morfológicas como las variaciones en la concentración de hemoglobina, incorporadas mediante funciones dieléctricas dependientes de la longitud de onda. Estas funciones se obtuvieron a partir de datos experimentales reportados en la literatura para distintas concentraciones de hemoglobina, empleando relaciones de Kramers-Kronig sustractivas con el fin de obtener la respuesta óptica en el rango espectral de 400 a 600 nm. La geometría de los eritrocitos se describió mediante óvalos de Cassini y los sólidos de revolución correspondientes se construyeron en SolidWorks.

Para implementar el modelo en CST Studio, primero se verificó la capacidad del modelo numérico para reproducir resultados con solución analítica. Para ello se realizó un análisis de los principales parámetros de simulación con el fin de identificar la configuración 	ue proporcionara la mejor concordancia con la solución analítica de teoría Mie para una esfera. Como resultado, se obtuvieron errores relativos inferiores al 5\%. Posteriormente, se comparó el patrón de esparcimiento de un eritrocito con resultados numéricos reportados en la literatura, encontrando que el modelo implementado en CST reproduce el patrón angular de esparcimiento. Los resultados del análisis de variación paramétrica permitieron establecer una configuración numérica adecuada de los parámetros principales del método computacional para el estudio de las propiedades ópticas de las distintas morfologías.

Al implementar los modelos y las funciones dieléctricas propuestas en CST, el análisis espectral mostró que la respuesta óptica de los eritrocitos está fuertemente modulada por las respuesta óptica de la hemoglobina, particularmente por tres bandas representativas: la banda de Soret, situada alrededor de 431 nm, la banda $\beta$ alrededor de 541 nm y la banda $\alpha$ alrededor de 573 nm. En particular, la región asociada a la banda de Soret presenta las mayores diferencias en la eficiencia de extinción entre las morfologías consideradas, mientras que las bandas $\beta$ y $\alpha$ conservan la estructura espectral general pero con variaciones menores. Por otro lado, los patrones de esparcimiento mostraron que, para las tres morfologías, la mayor parte de la luz esparcida se concentra en la dirección frontal. Sin embargo, las diferencias relativas entre células sanas y enfermas se vuelven más pronunciadas en las regiones lateral y posterior. En el caso del esferocito, se observó una reducción del esparcimiento lateral y de retroesparcimiento respecto al eritrocito sano. Por su parte, el microcito conserva en mayor medida la estructura angular característica de la geometría bicóncava, aunque presenta variaciones locales en la amplitud y posición de máximos y mínimos.

El análisis mediante ventanas angulares permitió identificar regiones en las que el la diferencia en el esparcimiento entre las morfologías de eritrocitos enfermos y el eritrocito sano aumenta de manera significativa. En particular, se encontró una región alrededor de $\theta\approx130^\circ$ que aparece de forma consistente para las tres longitudes de onda analizadas. Al extender el análisis a ventanas finitas tanto en $\theta$ como en $\varphi$, se observó que las bandas de Soret y $\alpha$ presentan regiones de alto contraste más amplias que la banda $\beta$, lo que proporciona una mayor tolerancia en la selección de los ángulos de detección.

Como resultados de este trabajo de tesis, se encontró que la respuesta óptica de los eritrocitos depende del efecto combinado de su geometría y de sus propiedades dieléctricas, y que las diferencias entre morfologías pueden intensificarse mediante una selección adecuada de la longitud de onda y del intervalo angular de detección. En particular, la banda de Soret y la región angular cercana a $\theta\approx130^\circ$ constituyen condiciones potencialmente relevantes para futuros esquemas de discriminación óptica entre eritrocitos sanos y patológicos en citometría de flujo.

 

   
% !TeX root = ../tesis.tex


\chapter{Trabajo a futuro}

\label{section:apendix1}



En cuanto a la parte experimental, se realizó la comparación con los resultados obtenidos de et al. quienes realizaron mediciones experimentales de la FSC para una muestra de eritrocitos sanos de 34.4 g/dL de CHCM de un citómetro de flujo, asumiendo incidencia en la cara circular del eritrocito ($\mathbf{k_1}$) en función de una incidencia lateral ($\mathbf{k_2}$) iluminada a $\lambda=632.8$ nm. En la Fig. \ref{subfig:STLFPi2} se muestran la FSC correspondiente a la polarización $\mathbf{k_2}$ en funcion de la $\mathbf{k_1}$, donde cada punto representa la medición correspondiente a un eritrocito y las zonas de color amarillo representan una mayor densidad de eritrocitos. En azul claro se resalta el punto correspondiente al cálculo de CST para el modelo de eritrocito empleado en este trabajo con una concentración de 30.6 g/dL. Se observa que los valores calculados que se integraron en $\theta$ de 2.2 a $8.2^{\circ}$ y en $\varphi$ de modo que $|\sin(\varphi)|\geq 2.2^{\circ}/\theta$, para simular las condiciones experimentales bajo las que se hicieron las mediciones. Se observa que el punto obtenido de manera numérica está dentro de las mediciones experimentales.

%
\begin{figure}[h!]
	\centering
	\includegraphics[width=0.43\textwidth]{4-Resultados/FSCpaper.pdf}\label{subfig:STLFPi2}
	\vspace{-0.1cm}
	\caption{ Eficiencia de extinción $Q_{\text{ext}}$ en función de la longitud de onda $\lambda$ de la onda incidente para un eritrocito (puntos azules), un esferocito (puntos rojos) y un microcito (puntos verdes) orientados con su eje de simetría en el eje $z$ iluminado con una onda en dirección del eje $z$ y polarizada en $x$. La línea continua que une los puntos es una guía visual. Se resaltan en gris las bandas que 
		rodean los mínimos asociados a la banda de Soret y a las bandas Q.}
	\label{fig:STLF}
\end{figure}
%

            
%-------------------------------------------------------------------------------
%                               References                                   |
%-------------------------------------------------------------------------------

\appendix
 % \input{7-Apéndices/1-Lorentz.tex} 
 % \input{7-Apéndices/2-SSKK.tex} 
 % \input{7-Apéndices/3-Res_Eritrocito_Sano.tex} 
 % \input{7-Apéndices/4-Res_Esferocito.tex} 
 % \input{7-Apéndices/5-Res_Microcito.tex} 


\setlength\bibitemsep{.1\itemsep}

\printbibliography

%\printindex
%-------------------------------------------------------------------------------
%                              Appendix                                   |
%-------------------------------------------------------------------------------


           
\end{document}
